\documentclass[color]{tumphthesis}

\usepackage{duckuments}
\usepackage{epigraph}
\usepackage{lipsum}
\usepackage{dsfont}
\usepackage{listings}

% %%% --- PACKAGES TO ADD --- %%%
% \usepackage{stmaryrd}
% \let\stmaryrdLightning\lightning
% \usepackage[english]{babel}
% \usepackage{listings}
% \usepackage{tikz}
% \usetikzlibrary{decorations.pathreplacing,calligraphy}
% \usepackage{gensymb}
% \usepackage{caption}
% \usepackage[normalem]{ulem}
% \usepackage{siunitx}  
% \usepackage{enumitem}
% \usepackage[list=true, font=small, labelfont=bf, labelformat=brace, position=bottom]{subcaption}
% \usepackage{mathtools}
% \DeclarePairedDelimiter\bra{\langle}{\rvert}
% \DeclarePairedDelimiter\ket{\lvert}{\rangle}
% \DeclarePairedDelimiterX\braket[2]{\langle}{\rangle}{#1 \delimsize\vert #2}
% \usepackage{mdframed}

% \usepackage[autostyle]{csquotes}

% %%% --- SETUP & FORMATTING --- %%%
% \numberwithin{equation}{section}
% \setlist[enumerate,1]{label=\roman*)}
% \sisetup{locale = DE}
% \mdfdefinestyle{box}{}

% \DeclarePairedDelimiter{\norm}{\lVert}{\rVert}

% \newcommand{\dd}{\ \mathrm{d}}
% %%% --- CODE BLOCK STYLING (listings) --- %%%
% % Note: xcolor is already loaded by the class, so we just define the colors
% \definecolor{codegreen}{rgb}{0,0.6,0}
% \definecolor{codegray}{rgb}{0.5,0.5,0.5}
% \definecolor{codepurple}{rgb}{0.58,0,0.82}
% \definecolor{backcolour}{rgb}{0.97,0.97,0.97}

% \lstdefinestyle{mystyle}{
%     backgroundcolor=\color{backcolour},   
%     commentstyle=\color{codegray},
%     keywordstyle=\color{codepurple},
%     numberstyle=\tiny\color{codegray},
%     stringstyle=\color{codegreen},
%     basicstyle=\ttfamily\footnotesize,
%     breakatwhitespace=false,         
%     breaklines=true,                
%     captionpos=b,                    
%     keepspaces=true,                 
%     numbers=left,                    
%     numbersep=5pt,                  
%     showspaces=false,                
%     showstringspaces=false,
%     showtabs=false,                  
%     tabsize=2
% }
% \lstset{style=mystyle}

% %%% --- CUSTOM MATH COMMANDS --- %%%
% \newcommand{\olsi}[1]{\,\overline{\!{#1}}} % overline short italic

% \newcommand{\tr}{\mathrm{tr}}
% \newcommand{\trr}[1]{\mathrm{tr[#1]}}
% \newcommand{\Tr}{\mathrm{Tr}}

% \newcommand{\m}{\mathcal}
% \newcommand{\id}{\mathrm{id}}
% \newcommand{\B}{\mathcal B}
% \newcommand{\R}{\mathbb{R}}
% \renewcommand{\L}{\mathcal{L}}
% \newcommand{\N}{\mathbb{N}}

% \newcommand{\Z}{\mathbb{Z}}
% \newcommand{\C}{\mathbb{C}}
% \newcommand{\End}{\text{End}}
% \newcommand{\Q}{\mathbb{Q}}
% \newcommand{\mi}{^{-1}}
% \newcommand{\de}{Delta}
% \renewcommand{\P}{\mathcal{P}}
% \newcommand{\PM}{\mathbb{P}}
% \newcommand{\hi}{\mathcal{H}}
% \newcommand{\e}{\mathrm{e}}
% \renewcommand{\Im}{\mathrm{Im}}
% \newcommand{\Ker}{\mathrm{Ker}}
% \renewcommand{\dim}{\mathrm{dim}}
% \newcommand{\pr}[1]{\ket{#1}\bra{#1}}
% \newcommand{\artanh}{\text{artanh}}
% \newcommand{\1}{\mathds{1}}
% \newcommand{\rk}{\mathrm{rk}}
% \newcommand{\Hom}{\mathrm{Hom}}
% \newcommand{\supp}{\mathrm{supp}}
% \newcommand{\inter}{\mathrm{int}}
% \newcommand{\sep}{\mathrm{SEP}}

%%% --- PACKAGES TO ADD --- %%%
\usepackage{stmaryrd}
\let\stmaryrdLightning\lightning
\usepackage[english]{babel}
\usepackage{listings}
\usepackage{tikz}
\usetikzlibrary{decorations.pathreplacing,calligraphy}
\usepackage{gensymb}
\usepackage{caption}
\usepackage[normalem]{ulem}
\usepackage{siunitx}  
\usepackage{enumitem}
\usepackage[list=true, font=small, labelfont=bf, labelformat=brace, position=bottom]{subcaption}
\usepackage{mathtools}
\usepackage{mdframed}
\usepackage[autostyle]{csquotes}
%%% --- SETUP & FORMATTING --- %%%
\numberwithin{equation}{section}
\setlist[enumerate,1]{label=\roman*)}
\sisetup{locale = DE}
\mdfdefinestyle{box}{}
\renewcommand{\abstractname}{Danksagung}

%%% --- CODE BLOCK STYLING (listings) --- %%%
% Note: xcolor is already loaded by the class, so we just define the colors
\definecolor{codegreen}{rgb}{0,0.6,0}
\definecolor{codegray}{rgb}{0.5,0.5,0.5}
\definecolor{codepurple}{rgb}{0.58,0,0.82}
\definecolor{backcolour}{rgb}{0.97,0.97,0.97}

\lstdefinestyle{mystyle}{
    backgroundcolor=\color{backcolour},   
    commentstyle=\color{codegray},
    keywordstyle=\color{codepurple},
    numberstyle=\tiny\color{codegray},
    stringstyle=\color{codegreen},
    basicstyle=\ttfamily\footnotesize,
    breakatwhitespace=false,         
    breaklines=true,                
    captionpos=b,                    
    keepspaces=true,                 
    numbers=left,                    
    numbersep=5pt,                  
    showspaces=false,                
    showstringspaces=false,
    showtabs=false,                  
    tabsize=2
}
\lstset{style=mystyle}

%%% --- CUSTOM MATH COMMANDS --- %%%
\newcommand{\olsi}[1]{\,\overline{\!{#1}}} % overline short italic
\DeclarePairedDelimiter{\norm}{\lVert}{\rVert}
\DeclarePairedDelimiter\bra{\langle}{\rvert}
\DeclarePairedDelimiter\ket{\lvert}{\rangle}
\DeclarePairedDelimiterX\braket[2]{\langle}{\rangle}{#1 \delimsize\vert #2}

\newcommand{\dd}{\ \mathrm{d}}
\newcommand{\tr}{\mathrm{tr}}
\newcommand{\trr}[1]{\mathrm{tr[#1]}}
\newcommand{\Tr}{\mathrm{Tr}}
\newcommand{\SP}{\mathrm{S}\mathbb{P}}
\newcommand{\m}{\mathcal}
\newcommand{\id}{\mathrm{id}}
\newcommand{\B}{\mathcal B}
\newcommand{\R}{\mathbb{R}}
\renewcommand{\L}{\mathcal{L}}
\newcommand{\N}{\mathbb{N}}
\newcommand{\mQ}{\mathcal{Q}}
\newcommand{\Z}{\mathbb{Z}}
\newcommand{\C}{\mathbb{C}}
\newcommand{\End}{\text{End}}
\newcommand{\Q}{\mathbb{Q}}
\newcommand{\mi}{^{-1}}
\newcommand{\de}{\Delta}
\renewcommand{\P}{\mathcal{P}}
\newcommand{\PM}{\mathbb{P}}
\newcommand{\hi}{\mathcal{H}}
\newcommand{\e}{\mathrm{e}}
\renewcommand{\Im}{\mathrm{Im}}
\newcommand{\Ker}{\mathrm{Ker}}
\renewcommand{\dim}{\mathrm{dim}}
\newcommand{\pr}[1]{\ket{#1}\bra{#1}}
\newcommand{\artanh}{\text{artanh}}
\newcommand{\1}{\mathds{1}}
\newcommand{\rk}{\mathrm{rk}}
\newcommand{\Hom}{\mathrm{Hom}}
\newcommand{\supp}{\mathrm{supp}}
\newcommand{\inter}{\mathrm{int}}
\newcommand{\sep}{\mathrm{SEP}}
\newcommand{\finedef}{\hfill$\lrcorner$}


\usepackage{xcolor}
\usepackage{soul}

\thesistitle
{Entanglement-Breaking Through Iterated Quantum Operations}

\thesisauthor{Malte Zehnpfennig}{Malte Zehnpfennig}
\thesisdate{2021}{8}{30}
\matriculationnumber{03789321}
\email{malte.zehnpfennig@gmx.de}
\supervisor{Prof. Dr. Michael Wolf}
\examiner{Prof. Dr. }
\date{\today}


\addbibresource{library.bib}

\usepackage[withpage, nohyperlinks]{acronym}

\begin{document}

\maketitle

% \abstractpage
% %\clearpage
%\chapter*{Abstract}
%\addcontentsline{toc}{chapter}{\protect Abstract}
Completely positive trace-preserving maps between matrices describe the most general evolution of quantum systems. These are called quantum channels. This thesis investigates the concept of entanglement-breaking quantum channels, with a particular focus on their behavior under composition. After reviewing the necessary preliminaries for quantum information theory and the study of quantum channels, some recent developments on entanglement-breaking channels, completely copositive channels (i.e., PPT channels), and their composition are reviewed. The so-called PPT$^2$ conjecture is investigated by searching for a possible violation via an SDP program, although none is found. Other necessary and separately sufficient conditions for entanglement-breaking depending on the matrix norm of the quantum channel are introduced, which provide some bounds on their entanglement-breaking indices. Furthermore, a sufficient condition for entanglement-breaking by using SIC-POVMs is provided. Lastly, the entanglement-breaking index of a set of quantum channels is considered, and some conditions are provided. 
 

\clearpage
\chapter*{Abstract}
\addcontentsline{toc}{chapter}{\protect Abstract}
%\clearpage
%\chapter*{Abstract}
%\addcontentsline{toc}{chapter}{\protect Abstract}
Completely positive trace-preserving maps between matrices describe the most general evolution of quantum systems. These are called quantum channels. This thesis investigates the concept of entanglement-breaking quantum channels, with a particular focus on their behavior under composition. After reviewing the necessary preliminaries for quantum information theory and the study of quantum channels, some recent developments on entanglement-breaking channels, completely copositive channels (i.e., PPT channels), and their composition are reviewed. The so-called PPT$^2$ conjecture is investigated by searching for a possible violation via an SDP program, although none is found. Other necessary and separately sufficient conditions for entanglement-breaking depending on the matrix norm of the quantum channel are introduced, which provide some bounds on their entanglement-breaking indices. Furthermore, a sufficient condition for entanglement-breaking by using SIC-POVMs is provided. Lastly, the entanglement-breaking index of a set of quantum channels is considered, and some conditions are provided. 
 


\acknowledgments
I would like to express my gratitude to Prof. Wolf for supervising this thesis. I am also thankful to my fellow students for their continuous encouragement throughout my Master's studies. Finally, I owe a special thank you to my family for their unwavering support.


\tableofcontents

% \addcontentsline{toc}{chapter}{\protect List of Figures}
% \listoffigures

% \addcontentsline{toc}{chapter}{\protect List of Tables}
% \listoftables

\begin{comment}
\clearpage
\chapter*{List of Abbreviations}
\addcontentsline{toc}{chapter}{\protect List of Abbreviations}

\begin{acronym}[SIC-POVM] % Ordered by first-use page for ascending page tracking.
    \acro{AEB}{Asymptotically Entanglement-Breaking}
    \acro{AES}{Asymptotically Entanglement-Saving}
    \acro{CP}{Completely Positive}
    \acro{CPTP}{Completely Positive Trace-Preserving}
    \acro{EB}{Entanglement-Breaking}
    \acro{EEB}{Eventually Entanglement-Breaking}
    \acro{HS}{Hilbert--Schmidt}
    \acro{LMI}{Linear Matrix Inequality}
    \acro{NP}{Nondeterministic Polynomial Time}
    \acro{POVM}{Positive Operator-Valued Measure}
    \acro{PPT}{Positive Partial Transpose}
    \acro{PVM}{Projection-Valued Measure}
    \acro{QDS}{Quantum Dynamical Semigroup}
    \acro{SDP}{Semidefinite Programming}
    \acro{SIC}{Symmetric Informationally Complete}
    \acro{SIC-POVM}{Symmetric Informationally Complete Positive Operator-Valued Measure}
    \acro{SN}{Schmidt Number}
    % \acro{TP}{Trace-Preserving}
    % \acro{BFGS}{Broyden--Fletcher--Goldfarb--Shanno}
    % \acro{COO}{Coordinate}
    % \acro{CVXPY}{Convex Optimization in Python}
    % \acro{ES}{Entanglement-Saving}
    % \acro{MOSEK}{MOSEK Optimization Software}
    % \acro{PT}{Partial Transpose}
    % \acro{QR}{QR Decomposition}
    % \acro{SCS}{Splitting Conic Solver}
    % \acro{SEP}{Separable}
    % \acro{SR}{Schmidt Rank}
\end{acronym}
\end{comment}

\mainmatter\setcounter{page}{1}
\chapter{Introduction}
\label{chap:intro}

\subsubsection{Background and Motivation}

In practice, quantum key distribution is severely limited by distance because signals inevitably degrade due to noise in the communication channel (like optical fibers or free space). To overcome this, standard quantum repeaters are used to establish long-distance secure links by performing entanglement distillation (extracting pure, perfect entangled states from many noisy ones) and reliably swapping them across multiple intermediate nodes. However, there are bound entangled states from which no pure entanglement can be distilled, but from which a secure cryptographic key can still be extracted \cite{horodecki2005secure}. This raised an important question: Could a specialized quantum key repeater be built to bypass the need for pure entanglement distillation, allowing one to swap and extend the secret keys from these highly noisy states over long distances? \cite{limitations2015} 

Quantum channels provide the fundamental mathematical framework for describing the evolution of quantum systems and the transmission of quantum information. One particular class of noise is modeled by completely copositive (PPT) channels, which, when applied to a subsystem, ensure that the global state has a positive semidefinite partial transpose. Consider a quantum key repeater scenario where the point-to-point communication links are modeled by such PPT channels. A minimal requirement for successful key distribution in this scenario is the preservation of entanglement across the shared bipartite states. Interestingly, in \cite{christandl2012ppt} the PPT$^2$ conjecture was formulated, stating that the composition of any such PPT channels is an entanglement-breaking channel, which completely annihilates the desired quantum correlations when applied to a subsystem. If true, this implies that a signal traversing two such sequential noise channels would have its quantum correlations completely annihilated, rendering the resulting end-to-end state fully separable.

Beyond its immediate practical implications for quantum network architectures, determining the exact number of iterations required for a quantum channel to become entanglement-breaking is also of mathematical interest. Significant research has recently been dedicated to characterizing the behavior of the entanglement-breaking property under compositions. Notably, the PPT$^2$ conjecture was proven to be true for qutrit systems in \cite{christandlWhenComposedMaps2019}, its general validity for arbitrary finite dimensions remains to be an open problem in quantum information theory. 

\subsubsection{Outline}
To establish the necessary framework for analyzing these phenomena, Chapter 2 provides a brief introduction to the core concepts of quantum information theory, alongside the basic mathematical foundations required. Chapter 3 focuses on the concepts of entanglement and positive maps. Chapter 4 subsequently offers a more detailed overview of quantum channels and some subclasses with a focus on their matrix representations. Chapter 5 introduces the core notions of eventually entanglement-breaking channels and entanglement-breaking indices. The chapter reviews some of the existing literature and highlights some recent theoretical developments in the field. Finally, Chapter 6 establishes some necessary and separately sufficient analytical conditions for determining whether a given channel is entanglement-breaking.

\chapter{Preliminaries}
\label{chap:2_general_preliminaries}

\section{Mathematics}\label{sec:math}

To start studying quantum channels, some fundamental concepts from analysis and linear algebra are needed and will be briefly presented in this section. As a reference, see for example \cite{Axler2024} for linear algebra or \cite{Abbott2015} for undergraduate analysis. The focus is on finite-dimensional spaces.

% \textcolor{green}{Before delving into the physical framework of quantum channels, we first establish the underlying mathematical formalism. This section outlines the essential concepts from linear algebra, functional analysis, and convexity required for the remainder of this thesis. For comprehensive treatments of these topics, we refer the reader to \cite{Axler2024} for linear algebra and \cite{Abbott2015} for analysis.}

% \textcolor{green}{Crucially, throughout this work, we restrict our focus exclusively to finite-dimensional complex vector spaces. This restriction significantly simplifies the topological considerations—for instance, ensuring that all linear operators are continuous and all norms are equivalent—while remaining perfectly sufficient for describing discrete-variable quantum information. We will adhere to the standard Dirac bra-ket notation to represent vectors and their duals. Furthermore, we adopt standard roman typography for operators and maps to distinguish them from variables, denoting the identity map as $\id$ (or the identity matrix as $\1$) and the trace as $\tr$.}

% \textcolor{green}{The following subsections will build our mathematical toolkit layer by layer. We begin by reviewing fundamental linear algebraic structures, including inner product spaces, operator bases, orthogonal projections, and fundamental subspace relations such as the orthogonal complement of kernels and ranges. Finally, we introduce the necessary analytical and geometric concepts—such as compactness, Lipschitz continuity, fixed-point theorems, and the structure of convex sets—which are indispensable for characterizing the state spaces of quantum systems and their evolution.}

\subsection{Linear Algebra}
A vector space over a field $\mathbb F$ is a set $V$ together with an addition $+:V \times V \rightarrow V$ and a scalar multiplication $\cdot : \mathbb F \times V \rightarrow V$ fulfilling the following properties: The addition operation $+$ needs to be commutative and associative, as well as having an inverse and identity element in $V$. The multiplication $\cdot$ needs to have an inverse and fulfill the distributive property with respect to addition on $V$. One usually writes $\cdot(\alpha, v)=:\alpha\cdot v=: \alpha v$ for $\alpha \in \mathbb F$ and $v \in V$. The additive identity is termed $0 \in V$, the multiplicative one $1\in V$, and for each $v \in V$ its additive inverse is $-v$. Elements $v$ of the vector space $V$ are called vectors. 

A subset of a vector space that is again a vector space is called a (linear) subspace. For example, it may be readily checked that the linear span of a set $\{v_i\}_{i\in I} \subset V$ of vectors, $\mathrm{span}\{v_i\}_{i\in I}:=\{\sum_{i\in I}\alpha_i v_i :  \alpha_i \in \mathbb F\  \forall \  i \in I\} $, is again a subspace of $V$. A set $\{v_i\}_{i\in I}$ of vectors is said to be a spanning system for a vector space $V$ if\footnote{An `if' in a definition is to be read as an `if and only if'.} its span equals $V$. On the other hand, it is defined to be linearly independent if the only set of scalars $\{\alpha_i \in \mathbb F\}_{i\in I}$ fulfilling $\sum_{i \in I}\alpha_i v_i = 0$ is $\alpha_i = 0$ for all $i \in I$. A set of vectors is a basis for a vector space if it is a spanning system and linearly independent. It can be shown that every vector space has a basis. This thesis focuses particularly on finite-dimensional vector spaces, which have a basis of finite length, i.e., a finite number of elements. In this case, every basis has the same length, which is defined to be the dimension of the vector space. Let $\{b_i\}_{i=1,\dots,d}$ be a basis for the vector space $V$ of dimension $d$. Then any element $v \in V$ may be written as 
\begin{align}\label{eq:coefficient_expansion}
v = \sum_{i=1}^d \alpha_i b_i
\end{align}
for suitably chosen coefficients $\alpha_i \in \mathbb F$. In general, to find the coefficients of a vector with respect to a basis, one has to solve a system of linear equations.


In quantum information theory, the underlying field is usually the set of complex numbers $\C$, in which case the vector space is called a complex vector space. $\C$ itself is a vector space with the usual operations. Thus, from now on, only complex vector spaces will be considered. It can be helpful to have a notion of length and angles for vectors. A norm $\lvert \lvert \cdot \rvert \rvert : V \rightarrow \R_{\geq 0}$ maps vectors to non-negative real numbers such that 
\begin{itemize}
    \item$\norm {v+w} \leq \norm v + \norm w $ for all $v,w \in V$ (Triangle inequality)
    \item$\norm{\alpha v}=\lvert \alpha \rvert \norm v$ for all $\alpha \in \C$ and $v \in V$ (Absolute Homogeneity)
    \item For all $v \in V$: ($\norm v = 0 \Rightarrow v = 0)$ (Definiteness) 
\end{itemize}

An inner product over a complex vector space $V$ is a map $\langle \cdot, \cdot \rangle: V \times V \rightarrow C $, that is linear in the second argument\footnote{Note that in the mathematical literature the inner product is often linear in the first and anti-linear in the second argument.}, conjugate symmetric in the sense that $\langle v, w \rangle = \overline{\langle w, v \rangle}$ holds for all $v,w \in V$, and positive semidefinite, i.e. $\langle v, v \rangle \geq 0$ $\forall v \in V$. Here, $\bar \cdot$ denotes complex conjugation. For example, the standard inner product on $\C^n$ is $\langle v,w\rangle = \sum_{i=1}^n v_i w_i$ for vectors $(v_i)_{i=1}^n$ and $(w_i)_{i=1}^n$ in $\C^n$. Every inner product induces a norm by setting $\norm \cdot \coloneqq \sqrt{\langle \cdot, \cdot \rangle}$. A vector space $V$ together with an inner product $\langle \cdot, \cdot \rangle$ is called an inner product space $(V, \langle \cdot, \cdot \rangle)$ and $(V, \norm \cdot)$ is a normed space whenever $\norm \cdot $ is a norm. Elements of an inner product space are called orthogonal if their inner product is zero. Similarly, elements of a normed space are said to be normalized or unit vectors if their norm is one. Thus, in inner product spaces there is the notion of an orthonormal basis $\{e_i\}_{i=1,\dots, d}$ whose elements are both pairwise orthogonal and normalized, that is, $\langle e_i, e_j\rangle = \delta_{i,j}$ for all $i,j \in \{1, \dots, d\}$. Let $v \in V$ be a vector, then $v = \sum_{i=1}^d \langle e_i , v \rangle e_i$, which may be checked by applying the map $\langle e_i, \cdot \rangle$ onto equation \eqref{eq:coefficient_expansion}.  

There are several interesting ways to construct new vector spaces. Of particular importance in quantum mechanics is the tensor product: Let $V_1$ and $V_2$ be two complex inner product vector spaces of finite dimension and $v_1\in V_1$ and $v_2 \in V_2$. The tensor product of $v_1$ and $v_2$ is defined as $v_1 \otimes v_2: V_1 \otimes V_2 \rightarrow \C, (w_1,w_2) \mapsto \langle w_1,v_1\rangle \langle w_2, v_2\rangle$. 

One defines the tensor product $V_1 \otimes V_2$ of the vector spaces as all linear combinations of such tensor products of vectors, i.e., $V_1 \otimes V_2 \coloneqq \{\text{bilinear forms }V_1'\times V_2' \rightarrow \C\}$. Consequently, if $\{e_i\}_{i=1}^n$ and $\{f_j\}_{j=1}^m$ are bases of $V_1$ and $V_2$, respectively, then $\{e_i \otimes f_j\}_{i=1}^n {}_{j=1}^m$ is a basis for $V_1 \otimes V_2$. Hence, $\dim(V_1 \otimes V_2) = \dim(V_1)\dim(V_2)$. Similarly, one defines the tensor operators of matrices $A \in \L(V_1)$ and $B \in \L(V_2)$ on the basis as $(A \otimes B)(e_i\otimes f_j) \coloneqq Ae_i \otimes Bf_j$ and extends the definition linearly onto the whole space. If $A$ and $B$ have eigenvalues $\lambda_i$ and $\mu_j$ with eigenvectors $a_i$ and $b_j$, respectively, then $A \otimes B$ has eigenvalues $\lambda_i \mu_j$ with corresponding eigenvectors $a_i \otimes b_j$.

\subsection{(Functional) Analysis}
\textcolor{red}{motivation for analysis} Consider a normed space $(V, \norm \cdot)$. For each point $v \in V$ the open $\epsilon$-ball around $v$ is defined to be the set $B_\epsilon(v):= \{w \in V: \norm{v-w}<\epsilon\}$ of all points whose distance to $v$ is less than $\epsilon$. A subset $O \subset V$ is called open if any of its points has an $\epsilon$-ball that is contained in $O$. On the contrary, a set is closed if its complement is open. A sequence $\{v_n\}_{n \in \N}$ is said to converge to a vector $v$ in $V$ if for every $\epsilon$-ball there is a sufficiently big index $N \in \N$ such that for all $n>N$ it holds that $v_n \in B_\epsilon(v)$. In that case, $v$ is called the limit of the sequence and one writes $v_n \to v$ as well as $v=:\lim_{n \to \infty} v_n$. A point $w \in V$ is an accumulation point of the sequence if for all $\epsilon > 0$ there is $n \in \N$ such that $v_n \neq w$ and $ \in B_\epsilon(w) $. Note that an accumulation point is a weaker notion than a limit point, since the sequence might still move away from it. A closed set $C$ is indeed 'closed' in the sense that the limit of every sequence that converges inside the vector space $V$ lies again in $C$. Consider a map $T: V \rightarrow V$. $T$ is called continuous at a point $v \in V$ if for every sequence $(v_n)_{n\in \N}$ that converges to $V$ it holds that $T(v_n) \to T(v)$. If a map $T$ is continuous in all points in its domain, then $T$ itself is said to be continuous. It can be shown that this is equivalent to every pre-image of open sets being open again. A set is compact if for any open covering of the set there is a finite open subcovering.

A map $T$ between complex vector spaces is linear if it preserves linearity, that is, $T(\alpha v + \beta w ) = \alpha T(v) + \beta T(w)$ holds for all scalars $\alpha,\beta$ and vectors $v,w$. For linear maps $T:V \rightarrow V$ on an inner product space $V$ with orthonormal basis $\{e_i\}_{i\in I}$ one defines the trace of $T$ as $\tr[T]\coloneqq \sum_{i\in I}\langle e_i, T e_i\rangle$. This definition is independent of the choice of the orthonormal basis. The image $T(V)$ is a vector space, and its dimension is the rank of $T$, $\rk(A)$. The kernel $\ker(T)$ is made up by all vectors in $V$ that are mapped to zero. The space of linear maps becomes itself a complex vector space again by defining scalar multiplication and addition point-wise. More specifically, let $(V, \norm\cdot)$ and $(W, \lvert\lvert \lvert \cdot \rvert \rvert \rvert)$ be normed spaces and $\mathcal L(V,W):=\{T:V \rightarrow W: \text{$T$ is linear}\}$ the vector space of linear operators. $\mathcal L(V,W)$ may be equipped with a norm by taking
\begin{align}
    \norm T := \sup\{\lvert \lvert \lvert T(v)\rvert \rvert \rvert: v \in V, \ \norm v = 1\}.
\end{align}
Whenever this mapping is finite, it defines a norm, and thus
\begin{align}
    \mathcal B(V,W) \coloneqq \{T \in \mathcal L(V,W): \norm T < \infty\}\text{ and }\mathcal B(V):=\mathcal B(V,V).
\end{align}
are normed spaces. Linear maps between vector spaces are also called operators, so that the norm defined above is called the operator norm and also depicted as $\norm\cdot _\text{op}$ or $\norm \cdot _{\infty}$. The operator norm is the infimum of all positive numbers $C$ such that $\lvert \lvert \lvert T(v)\rvert \rvert \rvert \leq C\norm v$ for all $v\in V$. Operators with finite operator norm are called bounded.

Linear maps are continuous if and only if they are bounded, so that $\mathcal B(V, W)$ equivalently constitutes the space of linear bounded continuous maps. In finite-dimensional spaces, all norms are equivalent, which in particular implies that convergence does not depend on the choice of norm and, hence, the boundedness of an operator does not depend on the norms of the underlying spaces. The space $V^*:=\mathcal B(V,\C)$ is termed the dual space of $V$ and its elements are called functionals. 


All linear maps between finite-dimensional vector spaces are equivalent to a matrix multiplication of the coefficients of the input vector with respect to some basis. The space of complex $n\times m$ matrices will be denoted by $\m M_{n\times m}$ and $\m M_n\coloneqq \m M_{n \times n}$. A basis for the space $\m M_{n \times m}$ is for example given by $\{E_{i,j}: i\in \{1,\dots,n\},  j \in\{1,\dots, m\}\}$ where $E_{i,j}$ is the matrix that has only zeros except for a one at the $(i,j)$ entry. Hence for finite-dimensional vector spaces $V$ and $W$ it holds that $\dim(\mathcal L(V,W)) = \dim(V) \cdot \dim(W)$.

Sequences $(v_n)_{n\in \N} \subset V$ whose elements come asymptotically closer, that is, for every $\epsilon > 0$ there is an $N\in \N$ so that for all $n,m>N$ it is $\norm{v_n - v_m}<\epsilon$, are called Cauchy sequences. If all Cauchy sequences in a normed space converge, then the vector space is complete. Finally, a Hilbert space is a complex inner product space that is complete with respect to the norm induced by the inner product. Hilbert spaces are essential in quantum information theory and usually denoted by the symbol $\hi$. The Riesz representation theorem states that, if $\hi$ is a Hilbert space, any functional $x^* \in \hi^*$ is of the form $x^*= \langle x, \cdot \rangle$ for some unique $x\in \hi$. In the physics literature, the bra-ket-notation is often used: Let $x^* = \langle x, \cdot \rangle$. Elements $x \in \hi$ are written as $\ket x \coloneqq x$ and functionals are written as $\bra x\coloneqq x^*$. The inner product then becomes $\langle x, y \rangle = x^*(y)\eqqcolon \braket{x}{y}$. Note that the norm usually still reads $\norm x$ instead of $\norm {\ket{x}}$.

Due to the Riesz representation theorem any bounded linear operator $T:X\rightarrow Y$ between Hilbert spaces $(X,\langle \cdot, \cdot \rangle_X)$ and $(Y, \langle \cdot , \cdot \rangle_Y)$ has a Hilbert adjoint $T^* \in \mathcal B(Y,X)$ that is uniquely defined via 
\begin{align}
    \langle Tx,y\rangle_Y = \langle x, T^* y \rangle_X \text{ for all } x\in X \text{ and }y\in Y.
\end{align} The mapping $T\mapsto T^*$ is an involution and is operator-norm-preserving. The Hilbert adjoint of a matrix with respect to the standard inner product on $\C^n$ is its complex conjugate transpose, that is, $(A^*)_{i,j} = (\overline A)_{j,i}$ for all complex matrices $A$. Note that the Hilbert adjoint depends on the inner product. An example of a finite-dimensional Hilbert space is the space of bounded linear operators $\mathcal B(\hi)$ on a complex Hilbert space of dimension $\dim(\hi)=:d$, that is, simply complex $d\times d$ matrices $\m M_d$, together with the Hilbert-Schmidt inner product $\langle A, B \rangle _\text{HS}\coloneqq \tr[A^*B]$. Operators $T \in \B(\hi)$ are
\begin{itemize}
    \item Hermitian or self-adjoint if $T^* = T$,
    \item positive semidefinite if there exists $B \in \B(\hi)$ such that $T = B^*B$,
    \item normal if $T$ and $T^*$ commute,
    \item a projection if $T^2 = T$, an orthogonal projection if $T=T^*=T^2$ and
    \item unitary if $T^*T = TT^*$, that is, $T^* = T^{-1}$.
\end{itemize}
The spaces of Hermitian, positive semidefinite and unitary complex $d\times d$ matrices will be denoted by $\text{Herm}(\C^d)$, $\m M_d^+$ and $\m U(d)$, respectively. Semidefinite positivity of matrices induces a partial ordering via $A\geq B:\Leftrightarrow A-B\geq  0 $ for matrices $A$ and $B$ of the same dimension. Clearly, any Hermitian operator is normal.  An operator $T\in \B(V,W)$ mapping between normed spaces $V$ and $W$ is called an isometry if it preserves the norm. If $V$ and $W$ are additionally Hilbert spaces, then an equivalent definition of an isometry is that $T^*T = \1_{V}$. A partial isometry is an operator that is an isometry on the orthogonal complement of its kernel (often called the initial space). Similarly, an operator is Hermitian if and only if it preserves the inner product, which is true if and only if $\langle v, T v\rangle \in \R$ for all $v\in \hi$. $T$ is positive semidefinite if and only if $\langle v, T v\rangle \geq 0$ for all $v\in \hi$, implying that they are in particular also Hermitian.

Additionally, one may consider the space $\mathcal B(\mathcal B(\hi_1),\mathcal B(\hi_2))$ where both $\mathcal B(\hi_1)$ and $\mathcal B(\hi_2)$ are equipped with the Hilbert-Schmidt inner product. Let $T \in \mathcal B(\mathcal B(\hi_1),\mathcal B(\hi_2)) $, its Hilbert adjoint satisfies $\tr[A^* T(B)] = \tr[(T^*(A))^*B]$ for all $B \in \mathcal B(\hi_1)$ and $A \in \mathcal B(\hi_2)$. In this context, $T^*$ is called the dual of the map $T$. Since the adjoint $^*:\B(\hi_2) \rightarrow \B(\hi_2)$ is bijective, this is equivalent to $\tr[A T(B)] = \tr[(T^*(A^*))^*B]$ and if $T$ (and thus $T^*$) preserves the adjoint, then the dual is uniquely defined via
\begin{align}
    \tr[A T(B)] = \tr[T^*(A)B] \text{ for all }B \in \mathcal B(\hi_1)\text{ and }A \in \mathcal B(\hi_2).
\end{align}

An important concept from linear algebra is eigenvalues: Let $A \in \m M_n$ and $\lambda \in \C$. If there exists a non-zero vector $v \in \C^n$ such that $Av = \lambda v$, then $\lambda$ is an eigenvalue of $A$ and $v$ a corresponding eigenvector. Lists of eigenvectors corresponding to different eigenvalues are linearly independent and if the matrix is normal, then that list is even pairwise orthogonal. More generally, if $T \in \B(X)$ is a linear bounded operator on a complete\footnote{The definition is provided for spaces of arbitrary dimension, where the reason for completeness of $X$ is this implying the resolvent being bounded again.} complex normed space $X$, then its spectrum is the set $\sigma(T):=\{\lambda \in \C: \lambda \1 - T \text{ is not bijective}\}$. Here, $\1$ is the identity on $X$. For maps on finite-dimensional spaces, injectivity and surjectivity are equivalent, so that the spectrum is exactly the set of eigenvalues in this case. It can be shown that the spectrum is non-empty and compact, hence, in particular, closed and bounded. The spectral radius is the supremum of the absolute value of all spectral values and it is upper bounded by the operator norm $\norm T$ of $T$. The peripheral spectrum is sometimes used to refer to all spectral values whose absolute value is equal to the spectral radius. The algebraic multiplicity of an eigenvalue is the order of the corresponding zero in the minimal polynomial, and its geometric multiplicity is the number of linearly independent eigenvectors corresponding to said eigenvalue. If $\lambda$ is an eigenvalue of a complex matrix $T\in \m M_n$, then $v\in \C^n, v\neq 0$ is a generalized eigenvalue if there is a $k\in \C^n $ such that $(T-\lambda \1)^k v = 0$. Again, every generalized eigenvector corresponds to only one eigenvalue.

Consequently, for any matrix $T \in \m M_n$ with eigenvalues $\lambda_1, \dots, \lambda_n$ (counted with algebraic multiplicity), its trace and determinant are directly given by the sum and product of its spectrum, that is, $\tr[T] = \sum_i \lambda_i$ and $\det(T) = \prod_i \lambda_i$.

One may define the square root of a quadratic matrix $T \in \m M_n$ as a matrix $\sqrt{T}$ such that $\sqrt{T}\cdot \sqrt{T} = T$. Any positive semidefinite matrix has such a square root, for example.
There are a lot of useful structure theorems for (subclasses of) linear operators depending on their spectral properties.
\begin{theorem}[Structure theorems for operators]
\
    \begin{itemize}
    \item (Polar Decomposition) Let $\hi$ be a Hilbert space of arbitrary dimension, $T \in \B(\hi)$, then there are partial isometries $U,V \in \B(\hi)$ such that $T= U|T| = |T|V$, where $|T|:= \sqrt{T^*T}$
    \item (Singular Value Decomposition) Let $T\in \m M_{n\times m}$ be a complex $n\times m$ matrix, then there exits an $n\times n$ unitary $U$, an $m\times m$ unitary $V$ and a diagonal $n \times m$ matrix $\Sigma$ such that $T = U\Sigma V^* = \sum_{i=1}^{\rk(A)} \sigma_i u_i \langle v_i, \cdot \rangle$, where the $\sigma_i$ are the eigenvalues of $\sqrt{T^*T}$ and $u_i$ and $v_i$ some suitably chosen unit vectors.
    \item (Diagonalization/Spectral Decomposition) Let $T\in \m M_n$ be a matrix whose eigenvalues $\{\lambda_i\}_{i=1}^n$ all have the same algebraic and geometric multiplicity. Then $T$ is of the form 
    \begin{align*}
        T = \sum_{i=1}^n \lambda_i \ket{R_i}\bra{L_i},
    \end{align*}
    where $\ket{R_k}$ and $\ket{L_k}$ are the right and left eigenvectors, respectively. They are biorthogonal, i.e., $\braket{L_i}{R_j} = \delta_{ij}$. If $T$ is normal, then the left and right eigenvectors are the same and need to be normalized in the diagonalization.
    \item (Jordan decomposition) For any complex square matrix $T\in \m M_n$ there is an invertible $X \in \m M_n$ and a block diagonal matrix $J \in \m M_n$ such that $T = XJX^{-1}$. The matrix $J$ is called the Jordan normal form of $T$ and of the form
    \begin{align*}
        J = \bigoplus_{k=1}^K J_k(\lambda_k), \ J_k(\lambda) \coloneqq \begin{pmatrix}
\lambda & 1       & 0      & \cdots  & 0      \\
        & \lambda & \ddots & \ddots  & \vdots \\
        &         & \ddots & \ddots  & 0      \\
        &         &        & \lambda & 1      \\
\smash{\makebox[0pt][l]{\hspace{0.75em}\raisebox{1em}{\text{\huge $0$}}}} & & & & \lambda
\end{pmatrix} \in \m M_{d_k}.
    \end{align*}
    Here, the $\lambda_k$ are the eigenvalues and $K$ is the number of linearly independent eigenvectors of $T$. $J_k$ is called a Jordan block corresponding to $\lambda_k$. The number of Jordan blocks corresponding to a given eigenvalue equals its geometric multiplicity. Their joint dimension $\sum_{k:\lambda_k = \lambda}d_k$ is the eigenvalue's algebraic multiplicity.  Adding the similarity transformation $X$, the matrix may be decomposed as 
    \begin{align*}
        T = \sum_{k=1}^K \lambda_k P_k + N_k,
    \end{align*}
    where $P_k$ is a (not necessarily orthogonal) projection and $N_k$ is nilpotent:
    \begin{align*}
        N_k^{d_k} = 0, N_kP_k =P_k N_k = N_k, P_kP_l = \delta_{kl}P_k , \tr[P_k]=d_k, \sum_{k=1}^K P_k = \1.
    \end{align*}
    {} 
\end{itemize}
\end{theorem}\noindent
The $\sigma_i$ in the singular value decomposition are called singular values of $T$. All spectral values of unitary operators have absolute value one and all eigenvalues of Hermitian and positive semidefinite operators are real and non-negative, respectively. The converse of all three statements holds if the operator is normal, which may e.g. be seen by considering the spectral decomposition.


The spectral decomposition allows a useful definition of functions $f: \C \rightarrow \C$ onto a matrix $T = \sum_{i=1}^d \lambda_i \ket{R_i}\bra{L_i}\in \m M_d$ via $f(T)\coloneqq \sum_{i=1}^df(\lambda_i)\ket{R_i}\bra{L_i}$. This definition works well since for example
\begin{itemize}
    \item it respects the algebra of functions, that is, $(f+g)(T)=f(T)+g(T)$, $(fg)(T) = f(g(T))$ and $f(g(T))=(f\circ g)(T)$,
    \item is consistent with power series
    \item and is invariant under similarity transformations, that is, $f(XTX^{-1}) = Xf(T)X^{-1}$ for all invertible matrices $X$.
\end{itemize}
There are other types of functional calculus, like holomorphic or Borel calculus. However, all of these coincide for normal matrices, see, e.g., \cite{Rudin1990-xf}, chapters 10 and 12.

When operators describe an evolution over time, they often form a semigroup. Let $V$ be a set and $T_t: V \rightarrow V$ be a mapping for every real $t\geq 0$. $\{T_t\}_{t\geq 0}$ is called a one-parameter semigroup if $T_0 = \1$ and $T_t T_s = T_{t+s}$ for all $s,t \geq0 $. If $V$ is upgraded to a normed space, then a one-parameter semigroup $\{T_t:V\rightarrow V\}_{t\geq 0}$ is strongly continuous if for all $x \in X$ it holds that $\vert \vert T_tx-x\vert\vert \overset{t \rightarrow 0}{\rightarrow} 0$. In this case, it is called a $C_0$-semigroup. A one-parameter semigroup is called uniformly continuous if $\vert \vert T_t - \1\vert \vert \overset{t \rightarrow 0}{\rightarrow} 0$. Uniform implies strong continuity and the two are equivalent in finite dimensions. The infinitesimal generator $A$ of a $C_0$-semigroup $\{T_t\}_{t\geq 0}$ is defined via
\begin{align}
    Ax \coloneqq \lim_{t \rightarrow 0}\frac{T_t x - x}{t}
\end{align}
for all $x \in V$ for which the limit above exists. A semigroup is uniformly continuous if and only if the generator is bounded, in which case the semigroup elements have the form $T_t = \exp(At)$, and the domain of $A$ equals the total space, see e.g., theorem 3.7 in \cite{semigroups} for the non-trivial direction.

Finally, some results on convexity are presented. For reference, see for example \cite{convex}. Let $V$ be a vector space. If a $v,w \in V$ then $\alpha v+(1-\alpha)w$ is called a convex combination of $v$ and $w$ whenever $\alpha \in [0,1]$. A set $S \subset V$ is called convex if it is closed under convex combinations. The convex hull $\text{co}(S)$ is defined to be the smallest convex set containing $S$. It can be shown that the convex hull is exactly the set of all convex combinations of elements in the set. By iteration, if $K$ is a convex set and $v_1, \dots, v_n \in K$ then $\sum_{i=1}^n \alpha_i v_i  \in K$ whenever the $\alpha_i$ make up a probability distribution. Elements that can not be written as a non-trivial convex combination of other elements are said to be extremal in that set. The Krein-Milman theorem states that a compact convex subset of a (Hausdorff locally convex topological) vector space is equal to the closure of the convex hull of its extremal points. Note that in finite dimension the span is already closed. In quantum information theory, one usually deals with complex vector spaces which are isomorphic to real spaces, so that Caratheodory's theorem may be useful: Any element in the convex hull of a set in an $n$-dimensional real vector space may be written as the convex combination of $n+1$ elements in said set. Notably, Caratheodory's theorem is constructive. 

A useful result in fixed-point analysis is Brouwer's fixed-point theorem which states that any continuous map mapping a non-empty,compact and convex set into itself has a fixed point.


\section{Basics of Quantum Information Theory}


This chapter aims to provide an overview of the basic concepts of quantum information theory, see for example \cite{watrousTheoryQuantumInformation2018} for more comprehensive literature. Consider an experiment in a physical laboratory with possible measurements. One notable example, that one can address in this regard is the Stern-Gerlach experiment \cite{Gerlach}. Here, a beam of silver atoms is directed through an inhomogeneous magnetic field $B$. The deflecting force acting on the magnetic dipole moment $\mu$ of the atoms is given by $F = \nabla({\mu}\cdot {B})$. While classical mechanics would predict a continuous spread of the beam along the gradient axis, the beam instead splits cleanly into two discrete paths, revealing that the intrinsic magnetic moment (and hence the spin) is space-quantized and takes only discrete values. By shooting a single silver atom through the magnetic field, this framework allows us to define the measurement as whether that silver atom is deflected to the top or towards the bottom. A fundamental assumption is that an experiment may be repeated identically, arbitrarily often. Each time the said measurement is performed, it produces an outcome that is recorded. The outcome of the measurement may change, so that after any number of measurements, there is a probability distribution regarding the outcomes. For example, after $n\in\N$ many implementations a distribution of the measurement outcomes would be $k$ times 'up' and $n-k$ times 'down' for some $k\in\{1,\dots, n\}$. The \textit{preparation} of a system is the set of all actions that determine this probability distribution. The collection of identically prepared systems is what physicists call an ensemble. A \textit{state} specifies the effect that this preparation has on the system. Mathematically, a state is modeled via a density operator, see for example \cite{Von_Neumann1996-bt} for a motivation.
\begin{definition}[Density operator] A state is described by a density operator, which is a positive semidefinite operator $\rho \in \B(\hi)$ on a Hilbert space $\hi$ such that $\tr[\rho] = 1$. The set of all density operators will be denoted by $\m D(\hi)$. A state is called pure if there exists a $\psi \in \hi$ such that $\rho = \pr\psi$.
\end{definition}\noindent
Note that the set of density operators $\m D(\hi)$ is convex, closed and bounded. By the spectral decomposition any density operator $\rho \in \m D(\hi)$ may be written as a convex combination of pure states, $\rho=\sum_{i=1}^{\dim(\hi)}p_i \pr{\psi_i}$, where $p_i$ are the eigenvalues of $\rho$ with corresponding eigenvectors $\psi_i$. The definition of density operators implies that the eigenvalues of $\rho$ make up a probability distribution. Take the Stern-Gerlach experiment as an example again. The degrees of freedom are 'up' and 'down' so that the corresponding Hilbert space is two-dimensional: $\hi = \text{span}\{e_\uparrow,e_\downarrow\}$. Hence, any state is of the form 
\begin{align}\label{eq:stern-gerlach}
\rho = p_\uparrow \pr{e_\uparrow}+p_\downarrow\pr{e_\downarrow}.
\end{align}

Mathematically, a measurement result needs to be an element of a sigma algebra $\m F$ with corresponding measurable space $(\Omega, \m F)$. It is postulated that each measurement $X \in \m F$ is assigned one operator $M(X)$ so that the probability $p(X|\rho)$ of getting the result $X$ from measuring a state $\rho \in \m D(\hi)$ is given by Born's rule $p(X|\rho) = \tr[\rho M(X)]$. 
Clearly, $p(\cdot|\rho)$ needs to be a probability measure on $(\Omega, \m F)$ for every density operator $\rho$. Hence the mapping $X \mapsto M(X)$ needs to have the following properties:
\begin{definition}[POVM]
    Let $(\Omega, \m F)$ be a measurable space and $\hi$ a Hilbert space. A POVM (short for Positive Operator Valued Measure) is a map $M:\m F \rightarrow \B(\hi)$ such that
    \begin{itemize}
        \item[i)]$M(X)\geq 0$ for all $X \in \m F$ and 
        \item[ii)] for every disjoint countable union $\Omega = \uplus_{n\in \N}X_n$, $X_n \in \m F$, it holds that $\sum_{n\in \N}M(X_n) = \1$.
    \end{itemize}
    The operators $M(X)$ are termed effect operators and the POVM is called sharp or a projective valued measurement (PVM) if additionally $M(X)$ is an orthogonal projection for every $X \in \m F$.{}
\end{definition}
If the number of measurement outcomes is countable, then define $M_n\coloneqq M(X_n)$ and write the ordered set $(M_n)_n$ for the POVM. For example, in the Stern-Gerlach experiment with Hilbert space $\hi = \text{span}(e_\uparrow,e_\downarrow)$ one has $\Omega = \{\text{'up'},\text{'down'}\}$. A possible PVM for the Stern-Gerlach measurement is $(\pr {e_\uparrow}, \pr {e_\downarrow})$. Thus, the probability of measuring 'up' and 'down' are given by $p(\text{'up'}|\rho) = \tr[\rho \pr{e_\uparrow}]=p_\uparrow$ and  $p(\text{'down'}|\rho) =p_\downarrow$, respectively, where $p_\uparrow$ and $p_\downarrow$ were defined in equation \eqref{eq:stern-gerlach}. It becomes apparent that pure states $\ket{\psi}\bra{\psi}$ are exactly those states for which there is a PVM, namely ($\pr \psi, \1 - \pr \psi)$, which has a unit probability outcome.

\subsubsection{Composite Systems}
Suppose there are two finite-dimensional Hilbert spaces $\hi_A$ and $\hi_B$ that describe systems that do not exclude each other. For example, the excitation and spin of a particle. Both partite states are described by density operators $\rho_A \in \m D(\hi_A)$ and $\rho_B \in \m D(\hi_B)$ on the partite systems. The joint state is described by the density matrix $\rho_A \otimes \rho_B$ on the composite Hilbert space $\hi_A\otimes \hi_B$. On the other hand, the states of a particle before and after a collision are mutually exclusive, and their joint state is thus described by the direct sum, as is known from classical physics.

Suppose there is a bipartite system $\hi_{AB}\coloneqq \hi_A \otimes \hi_B$ and a state $\rho \in \m D(\hi_{AB})$ on it. By Born's rule, the probability of measuring an outcome associated with an effect operator $X_{AB}=X_{AB}^* \in \m B(\hi_{AB})$ is $\tr[\rho X_{AB}]$. If the measurement only affects one subsystem, say system A, the effect operator is of the form $X \otimes \1$. In this case, there is an effective density operator describing the state on subsystem $A$ in the following sense:
\begin{theorem}[Partial Trace]\label{thm:partial_trace}
    There is a unique linear map $\tr_B: \m B(\hi_A \otimes \hi_B) \rightarrow \m B(\hi_A), \ \rho \mapsto \rho^A$ such that
    \begin{align*}
        \forall X \in \m B(\hi_1): \tr[\rho(X \otimes \1)] = \tr[\rho^A X].
    \end{align*}
    $\tr_B$ is called the partial trace and if $\rho \in \m D(\hi_A \otimes \hi_B)$ is a state, then $\rho^A$ is called its reduced state on system A. The partial trace has the following properties:
    \begin{enumerate}
        \item $\tr[\rho] = \tr[\rho^A]$ for all $\rho \in \m B(\hi_A \otimes \hi_B)$ (Trace preserving)
        \item $\rho \geq 0  \Rightarrow \rho^A \geq 0$ for all $\rho \in \m B(\hi_A \otimes \hi_B)$ (Semidefinite positivity preserving)
        \item For all $v,w \in \hi_A$ and every orthonormal basis $\{b_i\}_{i=1}^{d_B}$ it is $\langle v, \rho^A w\rangle = \sum_{i=1}^{d_A} \langle v \otimes b_i, \rho(w \otimes b_i)\rangle$.
        \item For all $A \in \m B(\hi_A)$ and $B \in \m B(\hi_B)$ it is $\tr_B[A \otimes B] = \tr[B]A$. {}
    \end{enumerate}
\end{theorem}\noindent
Properties i) and ii) of theorem \ref{thm:partial_trace} imply that the partial trace of a density operator is indeed again a density operator. 

Consider a product observable $A \otimes B \in \m M_{d_Ad_B}$ and a bipartite state $\rho \in \m D(\C^{d_A}\otimes \C^{d_B})$. The expectation value of the corresponding measurement is $\tr[\rho A \otimes B]$ and the expectation values that Alice and Bob measure on their subsystems are $\tr[A\rho^A]$ and $\tr[B\rho^B]$, respectively. If the product of the partite expectation values yields the global expectation, then the global state is called uncorrelated or a product state:
\begin{prop}
    Let $\rho \in \m D(\C^{d_A}\otimes \C^{d_B})$ be a bipartite density matrix. It is called a product state if 
    \begin{align*}
        \rho = \rho^A \otimes \rho^B \Leftrightarrow \left(\tr[\rho A \otimes B] = \tr[\rho^A A ]\tr[\rho^B B] \text{ for all }A\in \B(\C^{d_A}),B\in \B(\C^{d_B})\right)
    \end{align*}{}
\end{prop}
\begin{proof}
    The direction $\Rightarrow$ is straightforward. So assume the right hand side holds, then $\langle (\rho -\rho^A \otimes \rho ^B)^*, A \otimes B\rangle_\text{HS}$ holes for all $A$ and $B$. By linearity $\rho -\rho^A \otimes \rho ^B$ is orthogonal on $\m M_{d_Ad_B}$, which implies $\rho = \rho^A \otimes \rho^B$.
\end{proof}
All other states are correlated. However, in some states, this correlation is not quantum in nature but stems from a classical probabilistic preparation of the state. Such states are called separable. The definition is provided for positive semidefinite operators of arbitrary trace:
\begin{definition}
    A positive semidefinite bipartite matrix $O \in (\m M_{d_1}\otimes \m M_{d_2})^+$ is called separable if there is $n\in \N$ and $A_1, \dots, A_n \in \m M_{d_1}^+$, $B_1, \dots, B_n \in \m M_{d_2}^+$ such that $O = \sum_{i=1}^n A_i \otimes B_i$. Otherwise, the matrix is called entangled. The set of all such separable matrices will be denoted by SEP$(\C^{d_1},\C^{d_2})$ and the separable density matrices by $\m D_\text{sep}(\C^{d_1},\C^{d_2})$
\end{definition}\noindent
Equivalently, a bipartite density matrix is separable if and only if it is a classical mixture of product states. Note that $\m D_\text{sep}(\C^{d_1},\C^{d_2})$ is a convex closed and bounded set. There are several necessary or sufficient conditions for a state to be entangled. An overview is given in chapter \ref{ch:entanglement}. 

The following is a very useful result:
\begin{prop}[Hilbert-Schmidt isomorphism \cite{schatten1950theory}]\label{hilbert-schmidt}
Let 
\begin{align*}
I: \C^{d_1} \otimes \C^{d_2} \to \m M_{d_2\times d_1}, 
\ket \psi \otimes \ket \phi \mapsto \ket\psi\bra{\bar{\phi}}
\end{align*}
extended linearly. Then $I$ is bijective and
\begin{enumerate}
\item[(i)] For all $ \ket Phi, \ket \Psi \in \C^{d_1} \otimes \C^{d_2}$ it is $
\braket{ \Phi}{\Psi} = \operatorname{tr}[I(\ket \Phi)^* I(\ket \Psi)].
$
    \item[(ii)] For $A \in \m M_{d_1}$, $B \in \m M_{d_2}$ and $\Psi \in \C^{d_1}\otimes \C^{d_2}$ it is $
    I[(A \otimes B)\ket \Psi] = B I(\ket\Psi) A^T$
    where the transpose is taken in the basis used in the complex conjugation in the definition of $I$.
    
    \item[(iii)] If $d_1=d_2$ and $I(\ket\Psi)$ is invertible, then for all $A \in \C^{d_1}$ there exists a $B \in \m M_{d_2}$ such that
    \[
    (A \otimes \1)\ket\Psi = (\1\otimes B)\ket \Psi.
    \]{}
\end{enumerate}
\end{prop}
\begin{proof}\mbox{} % \mbox{} forces the list to start on a new line
(i) and (ii) follow from plugging in the definition. For (iii), apply (ii) to $(A\otimes \1)\Psi = (\1 \otimes B)\Psi$ yielding $B=I(\Psi)A^TI(\Psi)^{-1}$.

\end{proof}\noindent
Let $\ket \psi \in \hi_1 \otimes \hi_2$ be a bipartite vector. Applying the Hilbert-Schmidt isomorphism on it yields a matrix that has a singular value decomposition $I(\ket \psi) = \sum_i s_i \ket{f_i}\bra{\overline e_i}$. The inverse $\ket \psi = \sum_i s_i \ket{e_i}\otimes \ket{f_i}$ is called the Schmidt decomposition:
\begin{prop}[Schmidt decomposition]
    Let $\hi_1$ and $\hi_2$ be Hilbert spaces, $\ket \psi \in \hi_1 \otimes \hi_2$ and $r$ the number of non-zero singular values $s_i >0$ of $\ket \psi$. There are orthonormal bases $\{\ket{e_i}\}_{i=1}^{\dim(\hi_1)}$ and $\{\ket{f_i}\}_{i=1}^{\dim (\hi_2)}$ for $\hi_1$ and $\hi_2$, respectively, such that $
        \ket \psi = \sum_{i=1}^r s_i \ket{e_i}\otimes \ket{f_i}$. Hence, $\norm \psi ^2 = \sum_{i=1}^r s_i^2$. $r\eqqcolon \text{SR}(\ket \psi)$ is called the Schmidt-rank of $\ket \psi$ and the $s_i$ are the Schmidt values of $\ket \psi$.{}
 \end{prop}\noindent
Practically, the Schmidt decomposition can also be computed in the following way: With respect to chosen orthonormal bases expand $\ket \psi = \sum_{i=1}^{d_A}\sum_{j=1}^{d_B}C_{ij}\ket{ij}$. $C$ is a $d_A\times d_B$ complex matrix with singular value decomposition $C=U\Sigma V^*$ so that $C_{ij}=\sum_k U_{ik}\lambda_kV_{kj}^*$ and 
\begin{align}
    \ket \psi = \sum_k \lambda_k \underbrace{\left(\sum_i U_{ik}\ket i\right)}_{\eqqcolon \ket{e_i}}\otimes \underbrace{\left(\sum_j V_{kj}^*\ket j\right)}_{\eqqcolon \ket{f_i}}.
\end{align}
The Schmidt rank may equivalently be written as SR$(\ket\psi) = \text{rk}((id\otimes \tr)(\pr \psi))$ and is upper bounded by $d_A$ and $d_B$. Clearly, a pure state is separable if and only if its Schmidt rank is equal to one. Note that $\m D(\hi)$ is also a Hilbert space together with the $\norm \cdot _2$ norm, so that there is also a Schmidt decomposition for mixed states $\rho = \sum_{i=1}^r s_i E_i \otimes F_i$, where the $E_i$ and $F_i$ are each orthonormal with respect to the Hilbert-Schmidt inner product. Since there are mixed separable states, a Schmidt rank greater than one is not a criterion for entanglement. For example, the separable state $\rho =  \ket{00}\bra{00}/2+\ket{11}\bra{11}/2$ is already written in Schmidt decomposition and thus has Schmidt rank $\text{SR}(\rho) = 2$. To find a more fitting analog to the Schmidt rank for mixed states in the context of separability, the Schmidt number of a density matrix is introduced:
\begin{definition}[Schmidt number\cite{Terhal2000}]
    The Schmidt number $\text{SN}(\rho)$ of a density matrix $\rho \in \m D(\C^d)$ is the smallest $n\in \N$ such that $\rho$ admits a convex decomposition into pure states whose Schmidt ranks, each, do not exceed $n$. That is,
    \begin{align*}
        \text{SN}(\rho) \coloneqq \min\left\{\max_i \text{SR}(\ket{\psi_i}) : (p_i, \ket{\psi_i})_i \text{ is a pure state decomposition of }\rho\right\}
    \end{align*}{}
\end{definition}
Clearly, the Schmidt number and Schmidt rank of a pure state are the same: $\text{SN}(\pr \psi) = \text{SR}(\ket{\psi})$. A state is separable if and only if its Schmidt number is one.

\subsubsection{Evolution}





Suppose that after the silver atom passes through the magnetic field, it passes through another magnetic field with the same orientation as the first, then another, and so on. Interestingly, the results of the subsequent measurements are equal to those of the first one, with unit probability. See, for example, \cite{Frisch1933} for an experimental realization of a multi-stage Stern-Gerlach experiment. Consequently, a measurement may change the state. Such a change of state is not described by POVMs. Any map that describes a state change needs to map a density operator to another density operator (on a possibly different Hilbert space). Hence, one considers mappings $T:\m D(\C^{d_1}) \rightarrow \m D(\C^{d_2})$ that preserve positive-semidefiniteness and the trace. By the spectral decomposition of an Hermitian matrix into a positive and negative part, $T$ also preserves Hermiticity. Since any matrix $X = (X+X^*)/2+(X-X^*)/(2i)\eqqcolon A+B$ may be decomposed into its Hermitian and skew-Hermitian parts, it also is $T(X^*)=T(A)-T(B)=T(X)^*$. A state $\rho \in \m D(\C^d)$ may be the reduced state of a state $\sigma \in \m D(\C^d \otimes \C^{d_A})$ on a bigger system. In this case, the evolution of the composite system is given by $T\otimes \id_{d_A}$, which also needs to preserve the positive-semidefiniteness (it is trace-preserving if $T$ is). This leads to the following definitions:
\begin{definition}An operator $T\in \m B(\m M_{d_1},\m M_{d_2})$ is called 
\begin{itemize}
    \item trace-preserving if $\tr[T(X)]= \tr[X]$ for all $X\in \m M_{d_1}$,
    \item unital if $T(\1) = \1$,
    \item positive if $T(X)\geq 0$ for all $X \geq 0$,
    \item $n$-positive if $T\otimes \id_n$ is positive and
    \item completely positive if $T$ is $n$-positive for all $n\in \N$.
\end{itemize}
The set of completely positive linear maps from $\m M_{d_1}$ to $\m M_{d_2}$ will be denoted by $ {CP}(\C^{d_1}, \C^{d_2})$. $T$ is called a quantum channel if it is linear, trace-preserving and completely positive. The set of all quantum channels mapping a state on system $\C^{d_1}$ to $\C^{d_2}$ is written as $C(\C^{d_1}, \C^{d_2})$ and $C_d \coloneqq C(\C^d) \coloneqq  C(\C^d,\C^d)$. 
\end{definition}
It should not matter whether a composite system is tensored on the left or right side, and indeed $T\otimes \id_n$ is positive if and only if $\id_n \otimes T$ is positive. To see this, note that $T\otimes \id_n = (S_\text{out} \cdot S_\text{out}^*)\circ T \circ (S_\text{in} \cdot S_\text{in}^*)$, where $S_\text{in}:\C^d\otimes \C^{n}\rightarrow \C^{n}\otimes \C^{d}$ and $S_\text{out}:\C^{n}\otimes \C^{d}\rightarrow \C^{d}\otimes \C^{n}$ are the swap operators. Since swapping of an operator preserves the spectrum, it preserves positive semidefiniteness, from which the statement follows.

Thus, if $T$ is a quantum channel and $\rho$ a density matrix, then $\rho \mapsto T(\rho)$ is a valid evolution of a density operator. If one wants to describe an evolution of a quantum system in continuous time a $C_0$-semigroup of quantum channels is needed:
\begin{definition}
    A quantum dynamical semigroup (QDS) $\{T_t \in \m C(\C^d)\}_{t\geq 0}$ is a $C_0$-semigroup in which each element is a quantum channel.{}
\end{definition}

As explained in section \ref{sec:math}, every quantum channel $T$ has a dual map $T^*$ uniquely defined by the equation
\begin{align}\label{eq:dual}
\tr[AT(B)] = \tr[T^*(A)B]\text{ for all }A,B \in \B(\hi).
\end{align}
Suppose a system is in state $\rho \in \m D(\hi)$ and an outcome with corresponding effect operator $M$ is measured with probability $\tr[\rho M]$. Now the state evolves into $T(\rho)$ so that the new probability is $\tr[T(\rho)M]$. By using the dual, this probability is equal to $\tr[\rho T^*(M)]$, where $M \mapsto T^*(M)$ may be viewed as the evolution of the measurement process. The former, where the state evolves, is known as the Schrödinger picture and the latter, where observables evolve, as the Heisenberg picture. There are the following relations between properties of a linear map on matrix spaces and its dual:
\begin{prop}
    Let $T \in \m B(\C^{d\times d},\C^{d' \times d'})$ and $T^*$ its dual, then
    \begin{itemize}
        \item[(i)]$T\geq0\Leftrightarrow T^* \geq 0$
        \item[(ii)]$T$ is completely positive $\Leftrightarrow$ $T^*$ is completely positive
        \item[(iii)]$T$ is trace preserving $\Leftrightarrow$ $T^*$ is unital{}
    \end{itemize}
\end{prop}
\begin{proof}
    (iii) is proven by chosing $A=\1$ in equation \eqref{eq:dual} from which follows $\langle \1-T^*(A), B\rangle_{HS}=0$ for all $B\in \m M_d$ and $A\in \m M_{d'}$, so $\{\1-T^*(A)\}\perp \C^{d \times d}\Rightarrow T^*(A)=\1$ for all $A \in \C^{d'\times d'}$. For (i) choose $A=\ket{a}\bra{a}$ or $B=\ket{b}\bra{b}$ and (ii) by applying (i) on $T \otimes \1_n$ for every $n\in\N$.
\end{proof}
From this it is straightforward, that, if $(M_n)_n$ is a POVM, then so is $(T(M_n))_n$. Consider the following examples:


Finally, the state change due to measurements is described by instruments:
\begin{definition}[Quantum Instrument]Let $(\Omega, \m F)$ be a measurable space corresponding to a measurement. A quantum instrument is a map $I:\m F \rightarrow \m{CP}{\hi_1, \hi_2}, X \mapsto I_X$ such that for any disjoint parition, $\Omega = \uplus_k X_k$, $\sum_k I_{X_k}$ is trace-preserving.{}
\end{definition}
The post-measurement state after outcome $Y$ is then given by $I_Y(\rho)/\tr[I_Y(\rho)]$.

\subsubsection{Choi matrix and Kraus and environment representation}
For the sake of convenience, orthonormal bases for $\C^d$ will be denoted by $\{\ket i\}_{i=0}^{d-1}$. Another way to represent a linear mapping by a matrix is via the Choi-Jamilkowski-bijection. For this characterization, first consider the following lemma:
\begin{prop}[Cyclicity]\label{prop:cyclicity}
    The vector
    $\ket \Omega \coloneqq \sum_{i=1}^{d}\ket{i}\otimes\ket{i}$ is called the (unnormalized) maximally entangled vector. Let $\ket{\psi}\in \C^d\otimes \C^{d'}$, then $\ket \psi = (\1_d \otimes I(\psi))\ket \Omega$, where $I$ is the Hilbert-Schmidt map with respect to the basis used in $\ket \Omega$.
\end{prop}
\begin{proof}
    Decompose $\ket \psi = \sum_{i=1}^d\sum_{j=1}^{d'}c_{ij}\ket{ij}$, then $I(\psi) = \sum_{i=1}^d\sum_{j=1}^{d'}c_{ij} \ket{j}\bra{i}$ and thus 
    \begin{align*}
    (\1_d \otimes I(\psi))\ket \Omega = \sum_{ijk}c_{ij}\ket{k}\otimes \ket{j}\braket{i}{k} = \sum_{ij}c_{ij}\ket{ij}=\ket\psi
    \end{align*}
\end{proof}
By noting that $I(\ket\Omega) = \1$, which is invertible, one has $\ket \psi = (\1 \otimes A)\ket \Omega = (A^\text T\otimes \1)\Omega$, when $A\coloneqq I(\psi)$.
\begin{theorem}[Choi matrix\cite{Choi1975}]
    Let $T \in \m B(\m M_{d_1}, \m M_{d_2})$. Its Choi matrix with respect to an orthonormal basis $\{\ket{i}\}_{i=1}^{d_1}$ of $\C^{d_1}$ is $C_T\coloneqq (\id_{d_1} \otimes T)(\pr \Omega)=\sum_{i,j=1}^{d_1} \ket{i}\bra{j}\otimes T(\ket{i}\bra{j}) \in \m M_{d_1 \times d_2}$.
    \begin{enumerate}
        \item$T\mapsto C_T$ is a bijection with inverse $T(\cdot) = \tr_1[C(\cdot^T \otimes \1)]$, where the transpose is to be taken in the basis used in $\ket \Omega$. $T$ is also uniquely defined via $\tr[T(A)B]=\tr[C(A^\text T \otimes B)]$
        \item$C = C^* \Leftrightarrow T$ preserves the adjoint
        \item$C \geq 0 \Leftrightarrow T$ is completely positive.
        \item$\tr_2[C] = \1 \Leftrightarrow T$ is trace-preserving
        \item$\tr_1[C] = \1 \Leftrightarrow T$ is unital
    \end{enumerate}
\end{theorem}
\begin{proof}
    \begin{enumerate}
        \item a straightforward calculation confirms the inverse relations.
        \item Assume the Choi matrix is Hermitian. Then for all $B \in \m M_{d_2}$ it is $\tr[T(A^*)B] = \tr[C(\overline A\otimes B)] = \overline{\tr[C(A^\text T \otimes B^*)]}=\overline{\tr[T(A)B^*]}=\overline{\tr[T(A)^*B]}$, which then via the Hilbert-Schmidt inner product implies $T(A)^*=T(A^*)$ for all $A \in \m M_{d_2}$. The other direction is straightforward.
    \item Suppose $T$ is completely positive, then $C$ is clearly positive. For the other direction, one must show $(\id_n \otimes T)(\pr \psi)\geq 0$ for all $\psi \in \C^n \otimes \C^{d_2}$ since this extends by the spectral decomposition to all positive operators. Define $A\coloneqq I(\psi)^\text T$, then $\ket \psi = (A \times \1)\ket \Omega$ and thus
    \begin{align*}
        (\id_n \otimes T)(\pr \psi) = (\id_n \otimes T)((A \otimes \1)(\pr \Omega) (A^* \otimes \1)) = (A\otimes \1)C_T(A^*\otimes \1),
    \end{align*}
    which is positive if $C_T$ is positive.
    \end{enumerate}
    4. and 5. are straightforward by using the first point.
\end{proof}
In particular, 3. implies that a map $T \in \m B(\m M_{d_1}, \m M_{d_2})$ is completely positive if and only if it is $d_1$-positive.

With the Choi matrix, one can prove an important structural theorem for completely positive linear maps:
\begin{theorem}[Kraus representation \cite{kraus1971general}, \cite{Choi1975}]\label{thm:kraus}
A linear map $T \in \mathcal{B}(\mathcal{M}_d, \mathcal{M}_{d'})$ is completely positive if and only if it is of the form
\begin{align*}
    T(\cdot) = \sum_{i=1}^r K_i \cdot K_i^*.
\end{align*}
The operators $K_i$ are called Kraus operators of the map and the minimal number of them is the Kraus rank of $T$. The following statements hold:
\begin{enumerate}
    \item $T$ is trace preserving if and only if $\sum_{i=1}^r K_i^*K_i = \1$ and unital iff $\sum_{i=1}^r K_iK_i ^*= \1$.
    
    \item The Kraus rank is the rank of the Choi matrix.
    
    \item There is a set of Kraus operators that is mutually orthogonal with respect to the Hilbert-Schmidt inner product.
    
    \item \textit{Freedom:} Two sets of Kraus operators $\{K_j\}$ and $\{\tilde{K}_i\}$ represent the same map $T$ if and only if there is a unitary $U$ so that $K_j = \sum_i U_{ki}\tilde{K}_i$ (where the smaller set is padded with zeros).
\end{enumerate}
\end{theorem}

\begin{proof}
If $T$ is completely positive, then the Choi matrix $C_T$ is positive and consequently may be written as $C_T = \sum_{i=1}^r \pr {\psi_i}$, where $r = \rk(C_T)$. According to \ref{prop:cyclicity} it is $\ket  {\psi_i} = (\1\otimes K_i)\Omega$ for a suitably chosen $K_i$. Thus, $C_T = \sum_{i=1}^r (\1\otimes K_i)\pr \Omega (\1 \otimes K_i^*) = (\id \otimes T)\pr \Omega$, from which the desired Kraus decomposition of $T$ follows. To see that any map in Kraus decomposition is completely positive, one may simply consider the Choi matrix again. Eigenvectors of normal matrices may be chosen mutually orthogonal, i.e., $\braket{\psi_i}{\psi_j} = \bra\Omega \1\otimes K_i^*K_j  \ket \Omega \sim \delta_{ij}$. Using that $\bra{\Omega}A\otimes B\ket{\Omega}=\tr[A^\text T B]$, with transpose in the basis used in $\ket \Omega$, shows the possible orthogonality of the Kraus operators. Two spectral decompositions of the Choi matrix are the same, i.e. $\sum_i \pr {\psi_i} = \sum_j \pr{\tilde{\psi}_j}$, if and only if $\ket{\psi_i} = \sum_jU_{ij} \ket{\tilde\psi_j}$ \cite{Hughston1993}. Applying this to the Choi matrix proves the freedom of the Kraus operators.
\end{proof}\noindent
If a map $T$ has Kraus operators $\{K_i\}_i$ then its dual $T^*$ has Kraus operators $\{K_i^*\}_i$.

\begin{ex}[Some quantum channels]The following maps are examples of quantum channels, which may be verified by considering the Choi matrix e.g.:
\begin{itemize}
    \item $\m D(\C^d)\ni\rho \mapsto \rho \otimes \sigma$, where $\sigma \in \m D(\C^{d'})$ for some $d'\in \N$.
    \item $\m D(\C^d)\ni\rho \mapsto U\rho U^*$, where $U$ is a $d\times d$ complex unitary.
    \item $\m D(\C^d\otimes \C^{d'})\ni\rho \mapsto \tr_2[\rho]$\footnote{Or alternatively the partial trace over the first subsystem.}
\end{itemize}
\end{ex}
Note that compositions of completely positive and trace-preserving maps are again completely positive and trace-preserving, respectively. Consequently, the set of quantum channels $\m C(\hi_1, \hi_2)$ is closed under compositions. It turns out that any quantum channel is a composition of the three maps introduced above. This holds true for any quantum channel, but here the result is provided for channels between spaces of the same dimension:
\begin{theorem}[Environment representation, e.g. \cite{wolf2012quantum} theorem 2.5]
    Let $T \in C_d$ be a quantum channel. There is an auxiliary dimension $k$ and a unitary $U\in \m M_{d}\otimes \m M_{k}$ and a unit vector $\ket\phi \in \C^{k}$ so that
    \begin{align*}
        T(\cdot) = \text{tr}_{\C^k}[U(\cdot \otimes \pr \phi)U^*]
    \end{align*}
\end{theorem}



\chapter{Entanglement and positive Maps}\label{ch:entanglement}
\subsubsection{$n$-positive maps}
If a state $\rho =\sum_i p_i \rho_i^A\otimes \rho_i^B \in \m D(\hi_A \otimes \hi_B)$ is separable and $T$ is a positive map from $\B(\hi_A)$ to $\B(\hi_B)$, then $(T\otimes \id)(\rho)$ is again positive semidefinite. Consequently, every positive map $T$ provides a necessary condition for the separability of a density matrix. Clearly, for the criterion to be non-trivial, the map $T$ should not be completely positive since in this case $(T\otimes \1)(\rho)$ is positive even for entangled states. Again, it does not matter whether one considers $T \otimes \id $ or $\id \otimes T$, in the sense that both maps detect the same entangled states. To see this, one can again consider the swapping mechanism from above. For sufficiency, all positive maps must be considered:
\begin{theorem}\label{thm:positive-maps}[\cite{Horodecki1996} theorem 2]
    Let $\rho \in \m D(\C^d \otimes \C^{d'})$, then $\rho$ is separable if and only if $(T\otimes \id_{d'})(\rho)\geq 0$ for all positive maps $T \in \m B(\m M_d, \m M_{d'})$.{}
\end{theorem}
More generally, it can be shown that
\begin{theorem}[$n$-positive maps and entanglement \cite{Terhal2000}]\label{thm:n-pos}
    A density matrix has Schmidt number at most $n$ if and only if for all $n$-positive maps $T$ it is $(T \otimes \1)\rho \geq 0$. {}
\end{theorem}

An example of a positive, but not completely positive map is the transposition $\vartheta: X\mapsto X^T$ map since it preserves the spectrum, but for example
\begin{align}
    (\vartheta \otimes \id)(\pr \Omega) = \sum_{ij}\ket{j}\bra{i}\otimes \ket{i}\bra{j}
\end{align}
has a negative eigenvalue of $-1$ on the subspace spanned by $\{\ket{01},\ket{10}\}$ \cite{Peres1996}. The transposition criterion is closely linked to the concept of decomposable maps:
\begin{definition}[Decomposability]
    A positive map $T\in \m B(\m M_d , \m M_{d'})$ is called decomposable if there are completely postive maps $T_1, T_2 \in \m{CP}(\C^d,\C^{d'})$ such that $T = T_1+T_2 \circ \vartheta_{d_1}$. Otherwise, it is indecomposable.
\end{definition}
If $T$ is decomposable, then the only way that $(T\otimes \id)(\rho)$ can have a negative eigenvalue is if $(\vartheta \otimes \id)(\rho)$ has a negative eigenvalue. Hence, a decomposable positive map can detect entanglement if and only if it is also detected by partial transposition. The transposition criterion was the first separability condition for mixed states. For sufficiently small systems it is also sufficient for separability:
\begin{theorem}[Positive Partial Transposition (PPT) criterion \cite{Horodecki1996}]Let $\rho \in \m D(\C^{d_A}\otimes \C^{d_B})$ be a density matrix. If $(\vartheta_{d_A}\otimes \id_{d_B})(\rho)$ is not positive, then $\rho$ is entangled. The converse holds true if $d_Ad_B \leq 6$.{}\label{thm:PPT}
\end{theorem}
The proof of sufficiency in \cite{Horodecki1996} relies on results from \cite{Stormer1963} and \cite{Woronowicz1976} that show that every positive map $\m M_2 \rightarrow \m M_2$ and $\m M_2 \rightarrow \m M_3$ is decomposable, respectively. The rough idea in both of these papers was to consider extremal rays of positive maps and show that these are decomposable, and then using Krein-Milman theorem to extend this to all positive maps\footnote{This is not straightforward since the set of positive map is e.g. not compact, for more details the reader is referred to \cite{Horodecki1996}.}. By then using theorem \ref{thm:positive-maps}, the PPT theorem \ref{thm:PPT} is proven. In systems of size $3\times 3$ (and thus also in even higher-dimensional systems), there are entangled states with a positive partial transpose, this is also shown in \cite{Woronowicz1976}.

The notion of positive maps is related to that of so-called witnesses:
\begin{prop}[Entanglement witnesses]
A density operator $\rho \in \mathcal{M}_d(\mathbb{C}) \otimes \mathcal{M}_{d'}(\mathbb{C})$ has Schmidt number at least $n+1$ iff there exists a Hermitian operator $W \in \mathcal{M}_d(\mathbb{C}) \otimes \mathcal{M}_{d'}(\mathbb{C})$ such that $\operatorname{tr}[W\rho] < 0$ but $\bra{\psi}W\ket{\psi} \geq 0$ for all vectors $\ket{\psi} \in \mathbb{C}^d \otimes \mathbb{C}^{d'}$ with Schmidt rank $n$.
\end{prop}. Such a matrix $W$ is said to be an witness for entanglement in the case $n=1$.
\begin{proof}
Use theorem \ref{thm:n-pos} together with the hyperplane separation version of the Hahn-Banach theorem.
\end{proof}

\begin{prop}
Let $\mathcal{E}$ be positive but not completely positive. Let $\rho \in \mathcal{D}(\mathcal{H}_{AB})$ be such that $(\text{id} \otimes \mathcal{E})(\rho) \not\geq 0$. Then the corresponding entanglement witness to $\mathcal{E}$ for $\rho$ is 
\begin{align*}
    W &:= (\text{id} \otimes \mathcal{E}^*)(\ket{\eta}\bra{\eta})
\end{align*}
where $\ket{\eta}$ is an eigenvector of $(\text{id} \otimes \mathcal{E})(\rho)$ for an eigenvalue $\lambda < 0$.
\end{prop}

\begin{proof}
It is $
    \text{tr}[W\rho] = \bra{\eta}(\text{id} \otimes \mathcal{E})(\rho)\ket{\eta} = \lambda < 0$.
\end{proof}


\begin{definition}
An entanglement witness $W'$ is finer than $W$ if and only if it detects all entangled states detected by $W$ and at least one additional entangled state. It is optimal if there is no finer witness. Correspondingly, a positive map is optimal if its associated witness is.
\end{definition}

Another example of a positive, but not completely positive map is the Breuer-Hall map:
\begin{definition}[Breuer-Hall map \cite{Breuer2006, Hall2006}]If $U \in \m M_d$ is a skew-symmetric matrix, i.e. $U^\text T = -U$, and $U^*U\leq 1$ holds, then a map $T_\text{BH}:\m M_d \rightarrow \m M_d, X \mapsto \tr[X]\1_d - X - UX^\text TU^*$ is called a Breuer-Hall map.
{}
\end{definition}\noindent
Note that skew-symmetric unitary matrices only exist in even dimensions since $\det(U^\text T)=\det(U)=\det(-U)=(-1)^d\det(U)$ and unitarity sets $\vert \det(U)\vert = 1$. Consequently, Breuer-Hall maps with proper unitaries $U$ also only exist in even dimensions.
\begin{prop}[Breuer--Hall map \cite{Breuer2006, Hall2006}]\label{prop:breuer}
    Breuer-Hall maps $T_\text{BH}(X) = \tr[X]\1_d - X - UX^\text TU^*$ are positive, but not completely positive if $d>2$. If $U$ is additionally unitary, then $T_\text{BH}$ is indecomposable.
\end{prop}
\begin{proof}
    It needs to be shown that $T_\text{BH}(\pr \psi)=\1- \pr \psi - U\pr{\overline \psi}U^* \geq 0$ for all unit vectors $\ket \psi$. Show that $\pr \psi$ and $U\pr{\overline \psi}U^*$ are Hilbert-Schmidt orthogonal. This implies, for $\ket \phi\coloneqq U\ket{\overline \psi}$, that $\pr \psi + \pr \phi \leq \1$ since due to $U^*U \leq \1$ it is $\braket{\phi}{\phi} \leq 1$. The Choi matrix of the Breuer-Hall map is $C = \1_{d^2}-\pr \Omega - (\1 \otimes U)\mathbb F(\1 \otimes U^*)$, which is not positive since 
    \begin{align*}
        \bra \Omega C \ket \Omega &= d -d^2 - \sum_{ij}(\bra i \otimes \bra i U)\mathbb F (\ket j \otimes U^* \ket j) = d-d^2 - \sum_{ij}(\bra i \otimes \bra i U)(U^*\ket j \otimes j)\\
        &= d-d^2 -\sum_{ij}\braket{U^\text Tj}{i} \bra{i} U\ket{j}= d-d^2 -\sum_{j}\bra j \overline U U \ket j = 2d-d^2<0 \text{ if }d>2.
    \end{align*}
    Now let $U$ be unitary and assume that $T_\text{BH}(X) = \Phi_1(X)+\Phi_2(X^T)$ is decomposable for completely positive maps $\Phi_1$ and $\Phi_2$. Let $\ket \psi$ be an arbitrary unit vector, then $\bra \psi T(\pr \psi) \ket \psi = 1-1-0 = 0$. It follows that $\bra \psi \Phi_1(\pr \psi)\ket \psi = \bra \psi \Phi_2(\pr {\overline \psi})\ket \psi=0$ since the $\Phi_i$ are in particular positive. Since $\Phi_1$ is completely positive, there are Kraus operators $K_i$ such that $\Phi_1(\pr \psi) = \sum_i K_i\pr \psi K_i^*$. Hence, $\bra \psi \Phi_1(\pr \psi)\ket \psi = \sum_i \vert \bra \psi K_i \ket \psi\vert ^2 = 0 \Rightarrow \bra \psi K_i \ket \psi = 0$ for all unit vectors $\psi$, implying $K_i = 0$ for all $i$ and thus $\Phi_1 \equiv 0$. Consequently, $T_\text{BH}=\Phi_2\circ \vartheta$ and $T_\text{BH}\circ \vartheta = \Phi_2$ is completely positive. By using $(\id \otimes \vartheta)(\pr \Omega) = (\vartheta\otimes \id)(\pr \Omega)$ one has $C_{T_\text{BH}\circ\vartheta}=C_{T_\text{BH}}^{T_A} = \1_{d^2} - \mathbb F - (\1 \otimes U)\pr \Omega (\1 \otimes U^*)$. Test positivity on the vector $\ket{\phi}\coloneqq (\1 \otimes U)\ket \Omega$:
    \begin{align*}
        \bra{\phi}C_\text{BH}^{T_A}\ket{\phi} = d- \sum_{ij}U_{ij}U_{ij}^*-d^2 = 2d-d^2,
    \end{align*}
    which is again negative for $d>2$, resulting in the desired contradiction.
\end{proof}
To summarize, the Breuer-Hall maps make up a particularly nice positive map to test entanglement is that i) they have a 'tunable parameter': the skew-symmetric $U$, ii) they are indecomposable, that is, they can detect PPT entangled states and iii) they are optimal \cite{hall2006constructions}.

There are many more positive, but not completely positive maps, see e.g. \cite{wolf2012quantum}.

\subsubsection{Other entanglement criteria}
Besides positive maps, there are many criteria that are typically written in different ways.
\begin{prop}[Realignment criterion, \cite{chen2003matrix}]
    Let $\rho \in \m D(\C^d\otimes \C^d)$ be a bipartite density matrix. If $\sum_{i=1}^r \lambda_i >1$, then $\rho$ is entangled. The Schmidt values $\lambda_i$ are equivalently the singular values of the realigned matrix $\rho^R\coloneqq (\rho_{ij,kl})_{ik,jl}${}
\end{prop}

The Pauli matrices \cite{Pauli1927} 
\begin{align}
    \sigma_0 = \1, \sigma_x = \begin{pmatrix}
        1 & 0 \\
        0 & 1
    \end{pmatrix}, \sigma_y = \begin{pmatrix}
        0 & -i\\
        i & 0
    \end{pmatrix}\text{ and } \sigma_z = \begin{pmatrix}
        1 & 0 \\
        0 & -1
    \end{pmatrix}
\end{align}
are extremely useful in $d=2$ since they are Hermitian, unitary, and form an orthogonal basis for $\text{Herm}(\C^2)$ with respect to the Hilbert-Schmidt inner product. Generalizations of such matrices for higher dimensions are thus highly sought after. A generalization where the matrices are at least unitary is, e.g., given by the Heisenberg-Weyl matrices:
\begin{definition}[Heisenberg-Weyl basis \cite{Weyl1927}]
    Let $k,l \in \{0, \dots, d-1\}$ and $W_{kl}\coloneqq \sum_{j=0}^{d-1}\omega ^{jk}\ket{j}\bra{j\oplus l}$, where $\omega\coloneqq \exp(2\pi i/d)$ is the $d$-th root of unity and $\oplus$ is addition modulus $d$.
\end{definition}\noindent
The set $\{W_{kl}\}_{kl}$ is a unitary and orthogonal basis for $\m M_d$: $\trr{W_{kl}^*W_{k'l'}} = d \delta_{kk‘}\delta_{ll'}$. For a proof of this and more properties see for example \cite{hiesmayrBipartiteBoundEntanglement2025} for a recent review.  Higher-dimensional generalizations where the matrices are Hermitian instead of unitary are, e.g., the generalized Gell-Mann matrices \cite{Hioe1981}. 

If the rank of a bipartite density matrix is sufficiently small, then
\begin{prop}[Necessary rank condition for separability]\label{prop:rank-condition}Any separable matrix $\text{SEP}(\C^d,\C^{d'})$ fulfills $\rk(X)\geq \max\{\rk(X^A),\rk(X^B)\}$.{}
\end{prop}
\begin{proof}
    It is mentioned in theorem 1 of \cite{Horodecki2003} that if the rank condition above holds then $\rho$ is not distillable. Since separable states are in particular PPT and thus not distillable by \cite{Horodecki1998}, the statement follows.
\end{proof}
\begin{prop}[Low-Rank Separability, theorem 1 in \cite{Horodecki2000}]\label{prop:low-rank}
Let $\rho \in \mathcal{D}(\mathbb{C}^d \otimes \mathbb{C}^{d'})$ be a bipartite $d\times d'$ density matrix. If $\rk(\rho) \leq \min(d,d')$, then PPT is also sufficient for separability
\end{prop}










\subsubsection{Isotropic States}
For higher dimensions, detecting entanglement is generally a hard task. One class of states for which there are efficient necessary and sufficient criteria is that of isotropic states:
\begin{definition}[Isotropic state \cite{Horodecki1999}, chapter 5]
    A density matrix $\rho \in \m D(\C^d \otimes \C^d)$ is isotropic if for all unitaries $U\in \m U(d)$ one has $(U\otimes U^*)\rho(U\otimes U^*)^* = \rho$. The definition extends similarly to arbitrary matrices.
\end{definition}

On a bipartite system $\hi \otimes \hi$, define
    \begin{align}
        P^+\coloneqq P^\text{sym}\coloneqq \frac{1}{2}(\1+\mathbb F)=\pr{\phi^+} \text{ and }P^- \coloneqq P^\text{asy}\coloneqq \frac{1}{2}(\1-\mathbb F)=\1-\pr{\phi^+}. 
    \end{align}
Note that the swapping matrix $\mathbb F = \sum_{ij}\ket{i}\bra j\otimes \ket j \bra i$ clearly fulfills $\mathbb F^2 = \1$ and has thus only $+1$ and $-1$ as eigenvalues with respective subspaces $S_\pm$. $P^\text{sym}$ are $P^\text{asym}$ are exactly the projectors on these subspaces since $\mathbb  F = P^\text{sym}-P^\text{asy}$.
\begin{definition}[Werner state]
    If a state $W_p \in \m D(\C^d \otimes \C^d)$ is of the form 
    \begin{align*}
        W_p= p \frac{2}{d(d+1)}P^+ + (1-p)\frac{2}{d(d-1)}P^- = \frac{d-2p+1}{d(d^2-1)}\1 + \frac{2pd-d-1}{d(d^2-1)}\mathbb F
    \end{align*}
    for some $p\in[0,1]$, then it is a Werner state.
\end{definition}\noindent
$W_p \geq 0$ is obvious. Using $\trr{\mathbb F} = d$ and $\trr{P^\pm}\coloneqq (d^2\pm d)/2$ shows that they indeed have unit trace.
\begin{lemma}[\cite{werner1989quantum}] Werner states are invariant under unitary transformations of the form $U\otimes U$, i.e., for all $U \in \m U(d)$ and Werner states $W_p \in \m D(\C^d \otimes \C^d)$ it is 
\begin{align*}
    (U\otimes U)W_p (U\otimes U)^* = W_p.
\end{align*}
Additionally, any such state is a Werner state.
\end{lemma}
\begin{proof}
    It is sufficient to show the invariance of $\1$ and $\mathbb F$. The former is trivial and for the latter consider $
        (U\otimes U)\mathbb F(U\otimes U)^* = \sum_{ij}U\ket i \bra j U^* \otimes U\ket j \bra i U^* = \sum_{ij}G_{ij}\otimes G_{ij}^*$ where $U\ket i \bra j U^*\eqqcolon G_{ij}$. Now, $\{G_{ij}\}_{ij}$ is an orthonormal basis and it is $\mathbb F = \sum_{ij}G_{ij}\otimes G_{ij}^*$ for any such basis. For the converse, the reader is refered to the literature.
\end{proof}
\begin{lemma}
    A state is isotropic if and only if it is the partial transpose of a Werner state: $W_p^{T_B} = \alpha(p)\1 + \beta(p)\pr \Omega = (\alpha + \beta)\pr \Omega + \beta \pr \Omega^\perp$.
\end{lemma}
\begin{proof}
    Using that $\mathbb F^{T_B} = \pr \Omega$ it is 
    \begin{align*}
    \frac{d-2p+1}{d(d^2-1)}\1 + \frac{2pd-d-1}{d(d^2-1)}\pr \Omega.
    \end{align*}
    Clearly, $\1$ is invariant under $U\otimes U^*$. For $\ket \Omega$ one has $(U\otimes U^*)\ket \Omega = \sum_{ijk}U_{ji}\overline U_{ik} \ket{jk} = \sum_{jk}\delta_{jk}\ket{jk}=\ket \Omega$ which shows that $W_p^{T_B}$ is isotropic. See the literature once again for the converse; the basic idea is to the Schur-Weyl duality from algebra.
\end{proof}
\begin{definition}
    The fidelity $F$ of a state $\rho \in \m D(\C^d)$ is $
        F(\rho) = \bra \Omega \rho \ket \Omega/d$.
\end{definition}
\begin{lemma}\label{prop:fidelity_isotropic}
    An isotropic state of the form $\rho = \gamma \pr \Omega + \delta \pr \Omega ^\perp$ has fidelity $F\coloneqq F(\rho) = \gamma d $ so that \begin{align*}
        \rho = \frac{F}{d}\pr \Omega + \frac{1-F}{d^2}\pr \Omega^\perp =\frac{1-F}{d^2-1}\1+ \frac{Fd^2-1}{d(d^2-1)}\pr \Omega 
    \end{align*}
\end{lemma}
\begin{proof}
    Straightforward.
\end{proof}
\begin{lemma}[Eigenvalues of Werner states]
    Isotropic states $\rho = \alpha \1+\beta \mathbb F$ have eigenvalues $\alpha \pm \beta$ so that Werner states 
    \begin{align*}
        \frac{1-F}{d^2-1}\1+\frac{Fd^2-1}{d^2-1}\mathbb F
    \end{align*}
    have smallest eigenvalue $\lambda_\text{min} = (1-Fd)/(d(d-1))$. Hence, isotropic states are PPT if and only if $F\leq 1/d$.
\end{lemma}

\begin{definition}[Isotropic curling \cite{horodecki1999general}]
    The isotropic curling operator is defined as $\m T(\rho)\coloneqq \int_{\m U(d)} (U\otimes \overline U)\rho (U \otimes \overline U)^*\dd U$.
\end{definition}
\begin{lemma}
    \begin{enumerate}
        \item For all states $\rho \in \m D(\hi)$ $\m T(\rho)$ is isotropic and $\m T$ preserves the fidelity: $F(\rho) = F(\m T(\rho))$
        \item $\m T$ preserves separability
    \end{enumerate}
\end{lemma}
\begin{proof}
i) By construction, $\mathcal{T}$ outputs an isotropic state, since the Haar measure is translation-invariant. Because of $(U^* \otimes U^T)\ket{\Omega} = \ket{\Omega}$ for all unitaries $U$, the fidelity is preserved under the integral:
\begin{align*}
    F(\mathcal{T}(\rho)) &= \bra{\Omega} \mathcal{T}(\rho) \ket{\Omega} = \int \bra{\Omega} (U \otimes \bar{U}) \rho (U^* \otimes U^T) \ket{\Omega} dU \\
    &= \int \bra{\Omega} \rho \ket{\Omega} dU = F(\rho).
\end{align*}
ii) Let $\rho = \sum_k p_k \sigma_k \otimes \tau_k$ be an arbitrary separable state. The action of local unitaries yields $(U \sigma_k U^*) \otimes (\bar{U} \tau_k U^T)$, which is strictly separable. Because $\mathcal{T}(\rho)$ is an integral (a continuous convex combination) over these separable states, the map $\mathcal{T}$ inherently preserves separability.
\end{proof}
\begin{lemma}
    Let $\rho_\text{iso}(F)$ be an isotropic state with fidelity $F$, then for all $p \in [0,1]$ one has  
    \begin{align*}
        p\rho_\textbf{iso}(F)+(1-p)\rho_\text{iso}(F') = \rho_\text{iso}(pF+(1-p)F')
    \end{align*}
\end{lemma}
\begin{proof}
    Straightforward by using lemma \ref{prop:fidelity_isotropic}.
\end{proof}
\begin{prop}
    Let $\rho_\text{iso} \in \m D(\hi)$ be an isotropic state, then 
    \begin{align*}
        \rho_\text{iso} \text{ separable }\Leftrightarrow \rho_\text{iso}^{T_B}\geq 0 \Leftrightarrow F(\rho_\text{iso})\leq \frac{1}{d}
    \end{align*}
\end{prop}
\begin{proof}
$\rho_{\text{iso}} \text{ separable} \Rightarrow \rho_{\text{iso}}^{T_B} \ge 0$ is trivial. For the equivalence $\rho_{\text{iso}}^{T_B} \ge 0 \iff F(\rho_{\text{iso}}) \le \frac{1}{d}$, evaluate the spectrum of the partial transpose. The smallest eigenvalue of $\rho_{\text{iso}}^{T_B}$ is strictly $\lambda_{\min} = \frac{1 - Fd}{d(d-1)}$. Thus, $\rho_{\text{iso}}^{T_B} \ge 0$ is mathematically equivalent to $1 - Fd \ge 0$, which evaluates to $F \le \frac{1}{d}$. Finally, to prove $F \le \frac{1}{d} \implies \rho_{\text{iso}} \text{ separable}$, consider the trivially separable state $\sigma = \frac{1}{d} \1 \otimes \ket{0}\bra{0}$, which yields a fidelity $F(\sigma) = \frac{1}{d}$. By Lemma 3.0.5, its twirled version $\mathcal{T}(\sigma) = \rho_{\text{iso}}(1/d)$ strictly remains separable. Likewise, twirling the separable state $\sigma' = \frac{1}{d-1} (\1 - \ket{0}\bra{0}) \otimes \ket{0}\bra{0}$ yields the separable boundary state $\rho_{\text{iso}}(0)$. By the convexity established in Lemma 3.0.6, any isotropic state $\rho_{\text{iso}}(F)$ with $F \in [0, 1/d]$ is a convex combination of separable states and is therefore structurally separable.
\end{proof}


\chapter{Quantum Channels, their Representations and Properties} \label{ch:3_preliminaries_channels}
\section{More Matrix representations}
\subsubsection{Natural Matrix and Liouville Representation}\label{sec:matrix-rep}
As quantum channels are linear maps, they may be represented by a matrix with respect to a basis, see for example section 4.4.1 in \cite{heinosaariMathematicalLanguageQuantum2011} or section 2.3 in \cite{wolf2012quantum}. Note that the set of Hermitian matrices makes up a real vector space. Let $\{E_i\}_{i=0}^{d^2-1}$ and $\{F_j\}_{j=0}^{d'^2-1}$ be mutually orthogonal bases for $\text{Herm}(\C^d)$ and $\text{Herm}(\C^{d'})$, respectively, with respect to the Hilbert-Schmidt inner product. Then any $A \in \text{Herm}(\C^d)$ can be written as 
\begin{align}
    A = \sum_i a_i E_i \text{ with }a_i = \frac{\braket{E_i}{A}_\text{HS}}{\braket{E_i}{E_i}_\text{HS}}.\label{eq:vector-represen}
\end{align}
Let $T \in \m B(\C^d,\C^{d'})$, then
\begin{align}\label{eq:bloch-trafo}
    T(A) = \sum_j \frac{\tr[F_j^*T(A)]}{\tr[F_j^*F_j]}F_j   = \sum_{ij}T_{ji} a_iF_j = \langle \Vec{F}, \hat T\cdot \vec a\rangle \text{, where }T_{ji}\coloneqq \frac{\tr[F_j^*T(E_i)]}{\tr[F_j^*F_j]}
\end{align}
and $\hat T\coloneqq (T_{ij})_{ij}$, $\vec F\coloneqq (F_i)_i$ and $\vec a \coloneqq (a_i)_i$. That is, $\hat T$ is a $d'\times d$ matrix representation of the map $T$ and $T(A)$ has the coefficients $T(A)_j = \sum_i T_{ji}a_i$ with respect to $(F_j)_j$. Such a matrix representation will be called a natural matrix representation. Note that the natural matrix of the composition of linear maps is exactly the matrix multiplication of their natural matrix representations.

For the sake of simplicity and since this thesis is concerned mainly with powers of channels, now channels between spaces of the same dimension will be considered. I.e., $d=d'$ and let $E_i = F_i$ for all $i$. One may always choose $E_0 = \1$ and $\tr[E_i^*E_i] = d$ by the Gram-Schmidt procedure. By orthogonality all other basis elements are then traceless. An example of such a basis are the generalized Gell-Mann matrices, see e.g. \cite{Kimura2003}, which generalize the famous Pauli and Gell-Mann matrices in dimensions $d=2$ and $d=3$, respectively. In this case, if $T$ is trace-preserving, one has $
    T_{0j}=\tr[E_j]/d$ by using unitality of the dual $T^*$. In particular, $\hat T$ takes the form
    \begin{align}\label{eq:natural-matrix-rep}
        \hat T = \begin{pmatrix}
            1 & 0 \\
            v & \Delta
        \end{pmatrix},
    \end{align}
where $v \in \R^{d^2-1}$ and $\Delta$ is a $(d^2-1)\times (d^2-1)$ real matrix. 


Take a density matrix $\rho \in \m D(\C^d)$, its vector representation corresponding to equation \eqref{eq:vector-represen} is
\begin{align}
    \rho = \frac{1}{d}(\1+\vec r \cdot \vec E) \text{ with }r_i =\tr[E_i^*\rho].
\end{align}
The vector $\vec r \in \R^{d^2-1}$ is the Bloch vector of the state $\rho$. The Bloch picture is most famously used in the qubit case with respect to the Pauli basis, in which case the Bloch vector lies inside a sphere of radius one. A qubit state is pure if and only if its Bloch vector lies on the surface of the ball. In general dimension, the purity $\tr[\rho^2]$ is computed as 
\begin{align}
    \tr[\rho^2] = \frac{1}{d}(1+\vert \vert \vec r\vert \vert ^2).
\end{align}
Enforcing the upper bound of one on the purity thus yields $\vert \vert \vec r\vert \vert \leq \sqrt{d-1}$. In contrast to the qubit case, however, generally lying inside the ball of radius $\sqrt{d-1}$ is not sufficient to make up a valid density matrix.

As a concrete counterexample, consider a qutrit system ($d=3$), where the outscribed sphere of valid Bloch vectors has a radius of $\sqrt{3-1} = \sqrt{2}$. To satisfy the purity condition in equation (3.1.5), the traceless basis elements $\vec{E}$ must be normalized such that $\text{tr}[E_i^2] = 3$. So choose one of these orthogonal basis observables to be the diagonal matrix $E = \text{diag}(\frac{1}{\sqrt{2}}, \frac{1}{\sqrt{2}}, -\sqrt{2})$. Now, construct a state whose Bloch vector $\vec{r}$ points entirely along this axis with a magnitude of $||\vec{r}|| = 1$, it satisfies the bound $1 \le \sqrt{2}$. However, computing the corresponding operator yields
\begin{align}
\rho = \frac{1}{3}(\1 + E) = \frac{1}{3} \begin{pmatrix} 1 + \frac{1}{\sqrt{2}} & 0 & 0 \\ 0 & 1 + \frac{1}{\sqrt{2}} & 0 \\ 0 & 0 & 1 - \sqrt{2} \end{pmatrix}.
\end{align}
Because $\sqrt{2} > 1$, the third eigenvalue is strictly negative. Therefore, despite lying well within the bounded sphere, $\rho$ is not positive semi-definite and fails to be a valid density matrix.


As can be seen from equation \eqref{eq:bloch-trafo}, a quantum channel $T$ transforms the Bloch vector according to 
\begin{align}
\begin{split}
    r_i &\mapsto \sum_j T_{ij}r_j = \frac{1}{d}T_{i0}+ \sum_{j=1}^{d^2-1}T_{ij}r_j
    \Rightarrow \vec r \mapsto  \Delta\vec r+\vec v,
    \end{split}
\end{align}
Consequently, quantum channels are affine mappings between the Bloch vectors. Note that $\vec r = 0$ corresponds to the maximally mixed state which gets mapped as $\vec 0 \maps \vec v$. Consequently, $\vert \vert \vec v \vert \vert \leq \sqrt{d-1}$ is a necessary condition for $T$ being a quantum channel. It also follows directly that a channel is unital if and only if $\vec v = 0$. Unfortunately, in arbitrary dimension there are no simple necessary and sufficient criteria on the natural matrix representation for the channel to be completely positive. In the qubit special case, one can derive those necessary and sufficient conditions depending on $\vec v$ and the eigenvalues of $\Delta$, see \cite{ruskai}.

One may also choose a non-Hermitian basis for $\m M_d$, e.g., the standard matrices $E_{ij}\coloneqq \ket{i}\bra{j}$. The matrix of $T$ with respect to this basis is the Liouville representation. It is given by $\m L_{ij,kl}=\bra{i}T(\ket{k}\bra{l})\ket{j}$
according to \eqref{eq:bloch-trafo} since the basis $E_{ij}$ is also mutually orthogonal.



\subsubsection{$\chi$-matrix representation}
There is still another representation of a quantum channel by a matrix, called the $\chi$-matrix. This matrix was introduced in \cite{Chuang1997}; for a more recent review, see section 4.4.2 in \cite{heinosaariMathematicalLanguageQuantum2011}.
\begin{lemma}\label{lemma:basis-superoperators}
    Let $\{E_i\}_{i=1}^{d^2}$ be an orthonormal basis for $\m M_d$. For each $A \in \m M_d$, define $S_{ij}(A)\coloneqq A_i \rho A_j^*$. The maps $\{S_{ij}\}_{ij}$ form a basis for $\m B(\m M_d)$.
\end{lemma}
\begin{proof}
The Choi-Jamio\l{}kowski isomorphism establishes a bijective linear mapping between superoperators and $d^2 \times d^2$ matrices. For a general map $\mathcal{E}$, the Choi matrix is defined as $C_{\mathcal{E}} = \sum_{k,l=1}^d \ket{k}\bra{l} \otimes \mathcal{E}(\ket{k}\bra{l})$, where $\{\ket{k}\}$ is the standard computational basis. Applying this definition to the map $S_{ij}$ yields
\begin{align*}
    C_{S_{ij}} &= \sum_{k,l=1}^d \ket{k}\bra{l} \otimes \left(E_i \ket{k}\bra{l} E_j^*\right) = \left(\sum_{k=1}^d \ket{k} \otimes E_i \ket{k}\right) \left(\sum_{l=1}^d \bra{l} \otimes \bra{l} E_j^*\right) = \ket{\Psi_i}\bra{\Psi_j},
\end{align*}
where $\ket{\Psi_i} := \sum_{k=1}^d \ket{k} \otimes E_i \ket{k}$ is a vector in $\mathbb{C}^d \otimes \mathbb{C}^d$. Because $\{E_i\}_{i=1}^{d^2}$ forms a basis for $\mathcal{M}_d$, the Hilbert-Schmidt isomorphism guarantees that the corresponding set of vectors $\{\ket{\Psi_i}\}_{i=1}^{d^2}$ forms a basis for $\mathbb{C}^d \otimes \mathbb{C}^d$. Consequently, the set of rank-one outer products $\{\ket{\Psi_i}\bra{\Psi_j}\}_{i,j=1}^{d^2}$ forms a complete basis for the space of all $d^2 \times d^2$ matrices. Since the Choi mapping is an isomorphism, it inherently preserves linear independence and dimensionality. Therefore, the preimages $\{S_{ij}\}$ span the entire superoperator space and form a complete basis for $\mathcal{B}(\mathcal{M}_d)$.
\end{proof}
\begin{definition}[$\chi$-matrix]
Let $\mathcal{E}\in \m B(\m M_d)$ be a linear map between matrix spaces, $\{E_i\}_{i=1}^{d^2}$ be an orthonormal basis for $\m M_d$ and the $S_{ij}$ as defined in lemma \ref{lemma:basis-superoperators}. By the statement of said lemma, there exist scalars $\chi_{ij}$ such that for all $A \in \m M_d$ it is
\begin{align*}
\mathcal{E}(A) = \sum_{i,j=1}^{d^2} \chi_{ij} E_i A E_j^*.
\end{align*}
The matrix $\chi \eqqcolon \chi^{\m E}$ is called the $\chi$-matrix of the map $\m E$.{}
\end{definition}
Consider the important special case, where $\m E$ is a quantum channel (or at least completely positive), i.e., it has a Kraus decomposition $\m E = \sum_i K_i \cdot K_i^*$. In this case, the $\chi$-matrix can be constructed in the following way\footnote{Of course, the $\chi$-matrix can also be explicitly constructed for any other linear map.}: Any $A_i \in \m M_d$ can be written as $
    A_i = \sum_j a_{ij}E_j \text{ where }a_{ij}=\trr{E_j^*A_i}$. Then
\begin{align}
\mathcal{E} &= \sum_{i=1}^{\text{rk}(C_{\m E})} \sum_{j,l} a_{ij} \bar{a}_{il} E_j \cdot E_l^* = \sum_{j,l=1}^{d^2} \underbrace{\left( \sum_{i=1}^{\text{rk}(C_{\m E})} a_{ij} \bar{a}_{il} \right)}_{= \chi_{jl}} E_j \cdot E_l^*\label{eq:chi}
\end{align}

\begin{prop}
Different Kraus representations of the same channel yield the same $\chi$-matrix, so that the $\chi$-matrix of a channel is well-defined.
\end{prop}
\begin{proof}
Let $\{K_j\}_j$ and $\{\tilde{K}_k\}_k$ be two distinct sets of Kraus operators for the channel $\mathcal{E}$. By the unitary freedom of Kraus operators, see theorem \ref{thm:kraus} iv), there exists a unitary matrix $U$ such that $
K_j = \sum_k U_{jk} \tilde{K}_k$. Now, expand both sets of Kraus operators in the same orthonormal basis $\{E_m\}$: $
K_j = \sum_m a_{jm} E_m$ and $ \tilde{K}_k = \sum_m \tilde{a}_{km} E_m$.
Substituting this into the unitary relation between the different Kraus operators yields:
\begin{align*}
\sum_m a_{jm} E_m = \sum_k U_{jk} \left( \sum_m \tilde{a}_{km} E_m \right) = \sum_m \left( \sum_k U_{jk} \tilde{a}_{km} \right) E_m
\end{align*}
Since $\{E_m\}_m$ forms a basis, the coefficients on both sides must be equal: $
a_{jm} = \sum_k U_{jk} \tilde{a}_{km}$. Using this together with the definition of the $\chi$-matrix yields
\begin{align*}
\chi_{mn} &= \sum_j \left( \sum_k U_{jk} \tilde{a}_{km} \right) \left( \sum_l \bar{U}_{jl} \bar{\tilde{a}}_{ln} \right) 
= \sum_{k,l} \tilde{a}_{km} \bar{\tilde{a}}_{ln} \left( \sum_j U_{jk} \bar{U}_{jl} \right)\\  &= \sum_{k,l} \tilde{a}_{km} \bar{\tilde{a}}_{ln}\underbrace{(U^T \bar{U})_{kl}}_{=\delta_{kl}} =\sum_k \tilde{a}_{km} \bar{\tilde{a}}_{kn} = \tilde \chi _{mn}.
\end{align*}
\end{proof}
In general, the concatenation of channels does not correspond to the concatenation of the $\chi$-matrices (as is the case with the natural matrix representation). Furthermore, there are no simple criteria for determining whether a given matrix is the $\chi$-matrix of some quantum channel (as is the case with the Choi matrix). At least there is the following property of $\chi$-matrices:

\begin{prop}[\cite{heinosaariMathematicalLanguageQuantum2011} 4.45]
    Let $\m E \in \m B(\m M_d)$ be a linear map between matrix spaces and $\chi^{\m E}$ be its $\chi$-matrix. Then the following relations hold:
    \begin{enumerate}
        \item $\chi^{\m E} \geq 0 \Leftrightarrow$ $\m E$ is completely positive.
        \item $\m E$ is trace-preserving $\Rightarrow$ $\trr{\chi} = d$. The converse direction is, in general, not true.
    \end{enumerate}
\end{prop}
\begin{proof}
\begin{enumerate}
\item Let $\m E$ be completely positive, then using equation \eqref{eq:chi} yields for all $v \in \mathbb{C}^{d^2}$:
\begin{align*}
    \bra v \chi \ket v   &= \sum_{j,l} \bar{v}_j \chi_{jl} v_l 
    = \sum_{j,l} \bar{v}_j \left( \sum_{i=1}^{\text{rk}(C_{\m E})} a_{ij} \bar{a}_{il} \right) v_l \\
    &= \sum_i \left( \sum_j a_{ij} \bar{v}_j \right) \overline{\left( \sum_l {a}_{il} \bar{v}_l \right) }\geq 0.
\end{align*}
Conversely, assume $\chi \ge 0$. By the spectral theorem there is a diagonalization $\chi_{ij} = \sum_k C_{ik} \lambda_k \bar{C}_{jk}$ with non-negative eigenvalues $\lambda_k \ge 0$. Define $K_k := \sum_i \sqrt{\lambda_k} C_{ik} E_i$. Substituting this decomposition back into the general expansion $\mathcal{E}(A) = \sum_{i,j} \chi_{ij} E_i A E_j^*$ yields $
    \mathcal{E}(\rho) = \sum_k K_k \rho K_k^*$. Thus, $\m E$ is completely positive.
\item 
Computing the trace yields
\begin{align*}
    \text{tr}[\chi^{\m E}] &= \sum_j \sum_{i=1}^{\rk(C_{\m E})} a_{ij} \bar{a}_{ij} = \sum_{i,j} \text{tr}[E_j^* A_i] \text{tr}[E_j A_i^*] = \sum_{i,j} \langle E_j, A_i \rangle_\text{HS} \overline{\langle E_j, A_i \rangle}_\text{HS} \\
    &= \sum_{i,j} \langle A_i, E_j \rangle_\text{HS} \langle E_j, A_i \rangle = \sum_i \langle A_i, A_i \rangle_\text{HS} = \sum_i \text{tr}[A_i^* A_i] \quad \\
    &= \text{tr}[\1] = d.\footnotemark
\end{align*}\footnotetext{There is a minor printing error in the proof of proposition 4.45 in \cite{heinosaariMathematicalLanguageQuantum2011}, where it says H-S instead of HS for the Hilbert-Schmidt norm.}
To see that the converse of ii) is, in general, false, consider the following example: Define a completely positive map on two-dimensional spaces by the single Kraus operator 
\begin{align*}
K = \begin{pmatrix} \sqrt{2} & 0 \\ 0 & 0 \end{pmatrix}
\end{align*}

Because the map is defined via a Kraus operator, it is completely positive, ensuring its $\chi$-matrix is positive semi-definite. The trace of the process matrix is equal to the sum of the Hilbert-Schmidt norms of the Kraus operators:
\begin{align*}
    \text{tr}[\chi] &= \text{tr}[K^* K] = \text{tr}\left[ \begin{pmatrix} 2 & 0 \\ 0 & 0 \end{pmatrix} \right] = 2 = d
\end{align*}
Thus, the scalar trace condition is perfectly satisfied. However, the exact operator condition for trace-preservation requires $\sum_k K_k^* K_k = \1$. In this case, $K^* K \neq \1$. Applying this channel to the pure state $|1\rangle\langle 1|$ yields the zero matrix, which has a trace of 0. Because the channel destroys probability, it is not trace-preserving, despite satisfying $\text{tr}[\chi] = d$.

\end{enumerate}
\end{proof}


\begin{prop}[Matrix representation $\leftrightarrow$ $\chi$, section 4.4.2 in \cite{heinosaariMathematicalLanguageQuantum2011}]
    Let $\m E$ be a quantum channel, $\{E_i\}_i$ be an orthonormal basis, and $\hat{\m E}$ and $\chi_{\m E}$ the natural matrix and $\chi$-matrix with respect to $\{E_i\}_i$, respectively. The following relation holds true:
\begin{align*}
\mathcal{E}_{jk} &= \sum_{r,s} M_{jk,rs} \chi_{rs} \text{ where }M_{jk,rs} \coloneqq \text{tr}[E_j^* E_r E_k E_s^*].
\end{align*}
\end{prop}
\begin{proof}
Remember that $\mathcal{E}_{jk} = \text{tr}[E_j^* \mathcal{E}(E_k)]$. 
Next, we substitute the standard action of the channel expanded via the $\chi$-matrix representation, $\mathcal{E}(\rho) = \sum_{r,s} \chi_{rs} E_r \rho E_s^*$, using the basis element $E_k$ as the input state:
\begin{align*}
\mathcal{E}_{jk} &= \text{tr}\left[ E_j^* \left( \sum_{r,s} \chi_{rs} E_r E_k E_s^* \right) \right]
\end{align*}
By leveraging the linearity of the trace, we pull the summation and the scalar coefficients $\chi_{rs}$ outside of the trace operation:
\begin{align*}
\mathcal{E}_{jk} &= \sum_{r,s} \chi_{rs} \text{tr}[E_j^* E_r E_k E_s^*] = \sum_{r,s} M_{jk,rs} \chi_{rs},
\end{align*}
which yields the desired relation.
\end{proof}
\begin{prop}\label{prop:norm-chi}[Norm of the $\chi$-matrix]Let $\chi^{\m E}$ be the $\chi$-matrix of a quantum channel $\m E \in C_d$, then $1\leq \norm\chi_F \leq d$.
\end{prop}
\begin{proof}
By definition, $\chi \ge 0$ and $\text{tr}[\chi] = d$. Thus, its $d^2$ non-negative eigenvalues $\lambda_i$ sum to $d$, and its squared Frobenius norm is given by $\|\chi\|_F^2 = \sum_i \lambda_i^2$.

For the upper bound, the sum of squares is naturally bounded by the square of the sum:
\begin{align*}
    \|\chi\|_F^2 &\le \left( \sum_i \lambda_i \right)^2 = d^2.
\end{align*}
Equality holds if and only if exactly one eigenvalue is $d$ and the rest are $0$.

For the lower bound, the sum of squares is minimized when all eigenvalues are uniformly distributed ($\lambda_i = \frac{1}{d}$):
\begin{align*}
    \|\chi\|_F^2 &\ge \sum_{i=1}^{d^2} \left(\frac{1}{d}\right)^2 = d^2 \left(\frac{1}{d^2}\right) = 1.
\end{align*}
Equality holds if and only if $\lambda_i = \frac{1}{d}$ for all $i$. Taking the square root of both extreme cases yields $1 \le \|\chi\|_F \le d$.
\end{proof}
\begin{prop}[Relation between the norms of the matrix representations]\label{prop:relation-norms}
Let $\m E = \sum_i K_i \cdot K_i^* \in C_d$ be a quantum channel and $\hat{\m E}$ and $\chi^{\m E}$ the respective natural matrix and $\chi$-matrix representation with respect to the same orthonormal basis. The Choi matrix is, as usual, denoted by $C_{\m E}$. Then the following relation holds true:
\begin{align*}\norm{C_{\m E}}_F = \norm{\hat{\m E}}_F = \norm {\chi^{\m E}}_F = \sum_{ij}\vert \trr{K_i^* K_j}\vert^2.
\end{align*}
\end{prop}
\begin{proof}
    The norms of the Choi matrix and the natural matrix representation are the same since the matrices are reshufflings of each other, see equation (2.22) in \cite{wolf2012quantum}. It is $\hat{\mathcal{E}} = \sum_i K_i \otimes \overline{K}_i$ and thus
    \begin{align*}
 \|\hat{\mathcal{E}}\|_{\text{HS}}^2 &= \text{tr}[\hat{\mathcal{E}}^* \hat{\mathcal{E}}] = \sum_{i,j} \text{tr}[(K_i^* \otimes K_i^T)(K_j \otimes \overline{K}_j)] \\
&= \sum_{i,j} \text{tr}[K_i^* K_j] \text{tr}[K_i^T \overline{K}_j] = \sum_{i,j} |\text{tr}[K_i^* K_j]|^2
\end{align*}

Also remember $\mathcal{E}(\rho) = \sum_{i,j} \chi_{ij} E_i \rho E_j^*$, so $\hat{\mathcal{E}} = \sum_{i,j} \chi_{ij} (E_i \otimes \overline{E}_j)$, which implies

\begin{align*}
\|\hat{\mathcal{E}}\|_{\text{HS}}^2 &= \sum_{i,j,k,l} \chi_{ij}^* \chi_{kl} \text{tr}[E_i^* E_k] \text{tr}[\overline{E}_j^* \overline{E}_l] = \sum_{i,j} |\chi_{ij}|^2 = \|\chi\|_{\text{HS}}^2.
\end{align*}
\end{proof}
\section{Possible properties of quantum channels}
\subsubsection{Contractivity}
\begin{comment}
\begin{definition}[\cite{Kitaev1997}]
Let $\m E \in \m B(\m M_d, \m M_{d'})$  be a linear map between matrix spaces. The induced and diamond norm of $\m E$ are defined as
\begin{align*}
\|\m E\|_1 := \sup \left\{ \|\m E(X)\|_1 : \|X\|_1 \leq 1 \right\} \text{ and }
\|\m E\|_\diamond := \sup_{n \in \mathbb{N}} \|\m E \otimes \text{id}_n\|_1,
\end{align*}
respectively. $\m E$ is said to be completely bounded if it has a finite diamond norm: $\norm{ \m E}_\diamond < \infty$.
\end{definition}
That is, the induced norm is simply the operator norm with respect to the trace norms on the respective spaces.
\begin{lemma}[Lemma 3.2.1 in \cite{wolf2023mathematical}]
Let $\mathcal{E} \in \m B(\m M_d, \m M_{d'})$ be completely bounded.
\begin{enumerate}
    \item[(i)] $\norm*{\mathcal{E}}_1 = \sup \left\{ \norm*{\mathcal{E}(|\varphi\rangle\langle\psi|)}_1 \mid \norm*{\varphi} = \norm*{\psi} = 1 \right\}$
    \item[(ii)] If $\m E$ preserves the adjoint, then $\norm*{\mathcal{E}}_\diamond = \sup_{n \leq d} \sup_{\norm*{\varphi}=1} \norm*{(\mathcal{E} \otimes \text{id}_n)(|\varphi\rangle\langle\varphi|)}_1$
    \item[(iii)] If $\mathcal{E}$ is completely positive, then $\norm*{\mathcal{E}}_\diamond = \norm*{\mathcal{E}}_1 = \sup_{\norm*{\varphi}=1} \text{tr}[\mathcal{E}(|\varphi\rangle\langle\varphi|)]$
    \item[(iv)] $\norm*{\mathcal{E}}_\diamond \le d \norm*{\mathcal{E}}_1$
\end{enumerate}
\end{lemma}
By statement (iii) from the lemma above, the induced and diamond norm of a quantum channel are equal to one. \textcolor{red}{mention relation to lipschitz constant}

\textcolor{red}{do i even need the stuff before wrt 1 norm and not states?}
\end{comment}

\begin{definition}
    Let $\m E \in C_d$ be a quantum channel, then its contraction coefficient with respect to density operators is defined as 
    \begin{align*}
        c_2(\m E) = \sup_{\rho,\sigma \in \m D(\C^d)}\frac{\norm{\m E(\rho - \sigma)}_2}{\norm{\rho - \sigma}_2}
    \end{align*}
\end{definition}
Note that one can write
\begin{align}
    \rho = \frac{\1}{d}+\underbrace{(\rho-\frac{\1}{d})}_{\eqqcolon X_\rho \in \m T},\label{eq:trace-less}
\end{align}
where $X_\rho$ is the trace-less part of $\rho$ and $\m T\coloneqq \{A=A^*:\trr{A} = 0\}$ the space of Hermitian trace-less matrices. With this, $c_2$ is equivalently the contraction on this trace-less subspace:
\begin{prop}\label{prop:contraction-traceless}
    For all $\m E \in C_d$ it is 
    \begin{align*}
        c_2(\m E) = \sup\left\{ \frac{\norm{\m E(A-B)}_2}{\norm{A-B}_2}: A,B \in \m T\right\}.
    \end{align*}
\end{prop}
\begin{proof}
    Call the right-hand side of the equation in the proposition (1). By using equation \eqref{eq:trace-less} one has that 
    \begin{align*}
        c_2(\m E) = \sup\left\{ \frac{\norm{\m E(A-B)}_2}{\norm{A-B}_2}: A,B \in \m T\cap \m D(\C^d)\right\}\eqqcolon (2)
    \end{align*}
     Let $A \in \mathcal{{T}}$, so $A = \sum_i \lambda_i \pr{\lambda_i} $ with $\sum_i \lambda_i = 0$, and $\1/d+ A = \sum_i \left( \1/d + \lambda_i \right) \pr{\lambda_i}$. Hence, there is an $\epsilon > 0$, so that $1/d + \epsilon \lambda_i\geq 0$ for all $i$, making $\1/d + \epsilon A \geq 0$. Since $\m E$ is linear, one may add such an $\epsilon$ in both the denominator and numerator and then add and subtract $\1/d$.
\end{proof}


\subsubsection{Spectral Properties}
Note that the set of quantum channels $C (\C^d, \C^{d'})$ is convex, closed and bounded. By Caratheodory's theorem, any quantum channel $\m E \in C_d$ can be written as a (pointwise) convex combination of $d^^4 - d^2 +1$ quantum channels. For mixed unitary channels this number furthermore decreases ot $d^4-2d^2+2$, see for exmample proposition 4.4 in \cite{watrousTheoryQuantumInformation2018}.

To investigate the spectral properties, consider a linear map $\m E:\m M_d \rightarrow \m M_d$. If $T$ is Hermiticity-preserving, then all eigenvalues come in complex conjugated pairs: $\m E(X) = \lambda X$ and $\m E(X^*)=\overline{\m E(X)}=\overline{\lambda}X^*$. There is the following relation between the spectrum of a map and its dual:
\begin{prop}Let $\m E:\m M_d\rightarrow \m M_d$ be linear, then $\sigma(T)=\overline{\sigma(T^*)}$
\end{prop}
\begin{proof}Follows from 
$\det(\lambda \1 - \m E)=\overline{\det(\overline \lambda \ 1 -\m E^*)}{=}0$
\end{proof}
\begin{prop}[Spectral radius of positive maps, see e.g. in proposition 6.1 in \cite{wolf2012quantum}]
    Let $\m E:\m M_d \rightarrow \m M_d$ be linear and positive, then $r(\m E) = \sup_{\lambda \in \sigma(\m E)}|\lambda| \leq \norm{\m E(\1)}_\text{op}$. If $\m E$ is trace-preserving or unital, then there exists an eigenvalue with absolute value one.
\end{prop}
In particular, the eigenvalues of quantum channels lie on the unit circle in the complex plane.

\subsubsection{Fixed points and possible properties of channels}
\begin{definition}[Stationary state]
    A state $\rho \in \mathcal{D}(\mathbb{C}^d)$ is called stationary with respect to a quantum channel $\m E$ if it is a fixed point, meaning $\m E(\rho) = \rho$.
\end{definition}
Because the set of density matrices forms a compact and convex subset within the real vector space of Hermitian matrices, Brouwer's fixed-point theorem guarantees that every quantum channel possesses at least one stationary state.
\begin{definition}[Ergodic, Faithful and Primitive Channels]\label{def:primitive-etc}
A quantum channel $\m E \in C_d$ is called ergodic if it has only one unique stationary state and faithful if that stationary state is strictly positive. It is called primitive if for all states $\rho \in \m D(\C^d)$ the following holds true: $\lim_{n\rightarrow \infty} \m E^n(\rho) = \sigma>0$. That is, the power series of the channel converges to a replacement channel onto a strictly positive state.
\end{definition}
Primitive channels have a trivial peripheral spectrum, meaning that there is only one eigenvalue of absolute value one. And the definition of primitivity is equivalent to the existence of any finite power $n \in \N$ for which $\m E^n(\rho)>0$ is strictly positve for all states $\rho$, see for example theorem 6.7 in \cite{wolf2012quantum}.

Interestingly, primitive channels are dense in the set of all channels: Consider a channel $\m E \in C_d$, then the sequence $(\m E_n)_{n\ in\n}$ defined via $\m E_n(X)\coloneqq (1-1/n)\m E(\rho) + \1/(nd)$ converges to $\m E$ in operator norm. The $\m E_n$ are primitive since already
\begin{align*}
    \mathcal{E}_n(\rho) &= \underbrace{\frac{1}{n} \frac{\1}{d}}_{> 0} + \underbrace{\left(1 - \frac{1}{n}\right) \mathcal{E}(\rho)}_{\ge 0} > 0.
\end{align*}

Primitive quantum channels are ergodic: Assume $\mathcal{E}$ not ergodic, i.e. $\mathcal{E}(\rho_i) = \rho_i$ for $i \in \{1,2\}$, then $
    \lim_{n \to \infty} \mathcal{E}^n(\rho_1) = \rho_1$ and $\lim_{n \to \infty} \mathcal{E}^n(\rho_2) = \rho_2 \quad \lightning$, which is a contradiction. 


\subsubsection{Exemplary Channels}

\begin{definition}
    Let $\m E \in C_d$ be a quantum channel. $\m E$ is
    \begin{enumerate}
        \item a replacement channel onto a density matrix $\sigma \in \m D(\C^d)$ if $\m E(\cdot) =\trr{\cdot}\sigma$. In this case, one writes $\m E \eqqcolon \m R_\sigma$,
        \item a unitary channel if $\m E = U \cdot U^*$ for a unitary matrix $U \in \m U(d)$ and a mixed unitary if it is a convex combination (defined pointwise) of unitary channels. A mixed Heisenberg-Weyl channel is a mixed unitary channel where the unitaries are Heisenberg-Weyl operators,
        \item a Schur\footnote{Also called Hadamard channel.} channel if there is a matrix $A \in M_d$ so that $ \m E(\cdot) = A *\cdot$,
        \item a pinching channel if there is a PVM $\{\Pi_i\}_i$ so that $\m E(\cdot ) = \sum_i \Pi_i \cdot \Pi_i$
        \item The depolarizing channel $
    \m C(\C^d)\ni \Phi_p =(1-p)\id + p \m R_{\1/d}$
    for some $p\in[0,1]$. The channel $\Phi_1$ is called the completely depolarizing channel.The depolarizing channel $\m C(\C^d)\ni \Phi_p =(1-p)\id + p \m R_{\1/d}$ for some $p\in[0,1]$. The channel $\Phi_1$ is called the completely depolarizing channel.
    \item The two quantum channels 
        \begin{align}
            \m E_\pm :\m M_d \rightarrow \m M_d, \ X \mapsto \frac{\tr[X]\1 \pm X^T}{d\pm 1}
        \end{align}
        are called the Werner-Holevo channels \cite{Werner2002}. They are self-dual, that is, $T_\pm ^* = T_\pm$.
        \item 
A channel $\mathcal{E}$ describes phase-damping if the $(\mathcal{E}^n)_{n \in \N}$ converges to
\begin{align*}
    \operatorname{diag}_{\{e_i\}_i} : X \mapsto \sum_i \bra{e_i}X\ket{e_i} \ket{e_i}\bra{e_i}
\end{align*}
for some basis $\{e_i\}_i$. The set $\{e_i\}_i$ is called the decoherence basis.

    \end{enumerate}
\end{definition}

The depolarizing channel is unital and its Choi matrix $C_{\Phi_p} = (1-p)\pr \Omega + p\1/d$ is isotropic. It is in self-dual and to find its Kraus operators one first needs the following lemma:

\begin{lemma}
Let $\Z_d \coloneqq \{0, \dots, d-1\}$ and $\{W_{kl}\}_{\Z_d^{\times 2}}$ the Heisenberg-Weyl matrices, then
    \begin{align*}
        \sum_{k,l \in \Z_d^{\times 2}}W_{kl}\rho W_{kl}^* = d\trr{\rho}\1
    \end{align*}{}
\end{lemma}
\begin{proof}
    Use the fact that the $d^2$ Heisenberg-Weyl matrices form an orthogonal basis for the space of $d \times d$ complex matrices under the Hilbert-Schmidt inner product: $
\text{tr}[W_{kl}^* W_{mn}] &= d \delta_{km} \delta_{ln}$.
\end{proof}
From this follows that the Kraus operators of the depolarizing channel are $K_0 = \sqrt{1-p+p/d^2}\1$ and $K_{kl}=\sqrt{p}K_{kl}/d$ for all $(k,l) \in \Z_d^{\times 2}\textbackslash \{(0,0)\}$.
\begin{prop}The depolarizing channel forms a a quantum dynamical semigroup in the following way: 
    $\{\Phi_{1-\e^{-t}}\}_{t\geq 0}$. From this also follows that $\Phi_p^n = \Phi_{1-(1-p)^n}$ for all $n\in \N$.{}
\end{prop}
\begin{proof}
Note that $\Phi_0 = \id$ and $(\Phi_s\circ\Phi_t)(\rho) = (1-t)(1-s)\rho + ((1-t)s+t)\1/d$. Hence $
        (\Phi_s\circ \Phi_t)(\rho) =  \e^{-(t+s)}\rho + (1-\e^{-(s+t)})\1/d$.
\end{proof}



\begin{lemma}\label{lem:heisenberg-choi}
    The Choi matrix of a mixed Heisenberg-Weyl channel $\m E = \sum_{kl}p_{kl}W_{kl}\cdot W_{kl}^*$ is a Bell-diagonal matrix $C_{\m E} = \sum_{kl} p_{kl}(\1 \otimes W_{kl})\pr\Omega (\1 \otimes W_{kl}^*) = \sum_{kl}p_{kl}\pr{\psi_{kl}}$
\end{lemma}



\chapter{Entanglement-Breaking Quantum Channels}
\label{chap:4_entanglement}




\subsubsection{Basics}

\begin{definition}[Entanglement breaking channel \cite{Horodecki2003}]\label{def:eb}
    Let $T \in \m C(\C^{d_A},\C^{d_B})$ be a quantum channel. It is entanglement-breaking if for all dimensions $d\in \N$ and states $\rho \in \m D(\C^d \otimes \C^{d_A})$ it holds true that $(\id_d\otimes T)(\rho)$ is separable. $T$ is said to have Holevo or measure-and-prepare form if there are density matrices $\sigma_k \in \m D(\C^{d_B})$ and a POVM $\{M_k\}_k$ on $\C^{d_A}$ so that $T(\cdot) = \sum_k \trr{M_k \cdot }\sigma_k$. The set of entanglement-breaking quantum channels $T\in \m C(\C^{d_A},\C^{d_B})$ will be denoted by $\text{EB}(\C^{d_A},\C^{d_B})$, $\text{EB}_d \coloneqq \text{EB}(\C^d)\coloneqq \text{EB}(\C^{d},\C^{d})$.
\end{definition}\noindent
Again, the order of the tensor product does not matter in this definition. By using the spectral decomposition, the density matrices and POVM effects can always be chosen to be rank-one (although the effect operators of the POVM may be multiples of each other.)
A closely related concept is that of an entanglement-annihilating channel, which destroys any entanglement within the system it acts on:
\begin{definition}[Entanglement-annihilating channel \cite{moravcikovaEntanglementannihilatingEntanglementbreakingChannels2010}]
    A quantum channel $\m E \in \m C_d$ is entanglement-annihilating if $\m E(\rho) $ is separable for all states $\rho \in \m D(\C^d)$.
\end{definition}
Clearly, a channel $\m E$ is entanglement-breaking if and only if $\id \otimes \m E$ is entanglement-annihilating.

An obvious example of a channel that is not entanglement-breaking is the identity channel $T=\id: \rho \mapsto \rho$ since then $(\id \otimes T)(\rho)$ is entangled whenever $\rho$ is.

\begin{lemma}
    A channel $T$ in Holevo form $T(\cdot) = \sum_k \trr{M_k \cdot }\rho_k$ has Kraus operators $K_{kmn}=\sqrt{\sigma_k}\ket m \bra n\sqrt{M_k}$ and Choi matrix $C_T = \sum_k M_k^\text T\otimes \sigma_k$. 
\end{lemma}
\begin{proof}
    \begin{align*}
\sum_{kmn} A_{kmn} \rho A_{kmn}^* &= \sum_{kmn} \sqrt{\sigma_k} |m\rangle\langle n| \sqrt{M_k} \rho \sqrt{M_k} |n\rangle\langle m| \sqrt{\sigma_k} \\
&= \sum_{km} \sqrt{\sigma_k} |m\rangle \text{Tr}[M_k \rho] \langle m| \sqrt{\sigma_k} = \sum_k \text{Tr}[M_k \rho] \sigma_k
\end{align*}
\end{proof}

The main result from \cite{Horodecki2003} is the following characterization of entanglement-breaking channels:
\begin{theorem}[\cite{Horodecki2003}, theorem 4]
    Let $\m E \in \m C(\C^{d_A},\C^{d_B})$ be a quantum channel, then the following statements are equivalent:
    \begin{enumerate}
        \item $\m E$ has Holevo form.
        \item $\m E$ is entanglement-breaking.
        \item The Choi matrix $C_{\m E}$ is separable.
        \item All Kraus operators of $\m E$ are rank-one.
        \item For all positive maps $\gamma$,  $\gamma \circ \mathcal{E}$ is completely positive.
        \item For all positive maps $\gamma$,  $\m E\circ\gamma$ is completely positive.
    \end{enumerate}
\end{theorem}
\begin{proof}
i) $\Rightarrow$ ii):
\begin{align*}
    (\text{id} \otimes \mathcal{E})(\rho) &= \sum_k (\text{id} \otimes \sigma_k) ( \text{id} \otimes \text{Tr}(M_k \cdot) ) (\rho) = \sum_k \text{tr}_2 \left[ (\text{id} \otimes M_k) \rho \right] \otimes \sigma_k \\
    &= \sum_k \gamma_k A_k \otimes \sigma_k,
\end{align*}
where $\gamma_k := \text{tr}[M_k \rho] \geq 0$, $\sum_k \gamma_k = 1$, and $A_k := \text{tr}_2 \left[ (\text{id} \otimes M_k) \rho \right] / \gamma_k  = \text{tr}_2[\sqrt{M_k} \rho \sqrt{M_k}]/\gamma_k\geq 0$,
implying that $(\text{id} \otimes \mathcal{E})\rho$ is separable.\\
ii) $\Rightarrow$ iii): trivial
iii) $\Rightarrow$ i):
\begin{align*}
    (\text{id} \otimes \mathcal{E})(\ket{\Omega}\bra{\Omega}) &= \sum_{ij} (\text{id} \otimes \mathcal{E})(\ket{i}\bra{j} \otimes \ket{i}\bra{j}) \\
    &= \sum_{ij} \ket{i}\bra{j} \otimes \mathcal{E}(\ket{i}\bra{j}) \\
    &\overset{3)}{=} \sum_n \tilde{p}_n \sigma_n^A \otimes \sigma_n^B = \sum_n p_n \ket{v_n}\bra{v_n} \otimes \ket{w_n}\bra{w_n}
\end{align*}
for suitably chosen normalized $\ket{v_n}, \ket{w_n}$ and $\sum p_n = 1$. Consequently, one has
\begin{align*}
    T(\rho) &= \text{tr}_1[C (\rho^T \otimes \1)] = \sum_n p_n \text{tr}_1 \left[ \ket{v_n}\bra{v_n} \rho^T \otimes \ket{w_n}\bra{w_n} \right] \\
    &= \sum_n p_n \bra{v_n} \rho^T \ket{v_n} \ket{w_n}\bra{w_n} = \sum_n p_n \bra{\bar{v}_n} \rho \ket{\bar{v}_n} \ket{w_n}\bra{w_n} \\
    &= \sum_n \text{tr} \left[ \rho \, p_n \ket{\bar{v}_n}\bra{\bar{v}_n} \right] \ket{w_n}\bra{w_n}.
\end{align*}
Now, i) - iii) are equivalent. iii) $\Rightarrow$ iv): Since
\begin{align*}
    \mathcal{E}(\rho) &= \sum_n \ket{w_n}\bra{w_n} \text{Tr}[p_n \rho \ket{v_n}\bra{v_n}] = \sum_n p_n \ket{w_n}\bra{w_n} \bra{v_n} \rho \ket{v_n} \\
    &= \sum_n p_n \ket{w_n}\bra{v_n}\rho \ket{v_n}\bra{w_n},
\end{align*}
the Kraus operators $K_n \coloneqq \sqrt{p_n}\ket{w_n}\bra{v_n}$ are rank-one.\\
iv) $\Rightarrow$ i):
\begin{align*}
    \mathcal{E}(\rho) &= \sum_k A_k \rho A_k^* = \sum_k \gamma_k \ket{v_k}\bra{w_k} \rho \ket{w_k}\bra{v_k} = \sum_k \text{Tr}[\gamma_k \rho \ket{w_k}\bra{w_k}] \ket{v_k}\bra{v_k}
\end{align*}
ii) $\Leftrightarrow$ v): Theorem \ref{thm:positive-maps} states that $\rho \text{ separable } \Leftrightarrow \forall \text{ positive } \gamma : (\text{id} \otimes \gamma)(\rho) \ge 0$ holds true for all states $\rho$. Now, one only needs to use that a channel is entanglement-breaking if and only if its Choi matrix is separable.\\
v) $\Leftrightarrow$ vi) 
Assume v) holds and let $\gamma$ be a positive map. By assumption, the composition $\gamma \circ \mathcal{E}$ is completely positive. Taking the adjoint preserves complete positivity, so that $
    (\gamma \circ \mathcal{E})^* = \mathcal{E}^* \circ \gamma^* $. Since $\gamma$ is positive, its adjoint $\gamma^*$ is also positive. Thus, for any positive map $\gamma^*$, the composition $\mathcal{E}^* \circ \gamma^*$ is completely positive. Because a channel is entanglement-breaking if and only if its adjoint is entanglement-breaking (as transposition preserves the separability of the Choi matrix), applying this reasoning to $\mathcal{E}$ directly yields vi).

\end{proof}
Point iii) of the theorem above implies that one only needs to check for the dimension $d_B$ in the identity $\id_d$ in the definition \ref{def:eb} of entanglement-breaking channels. An example of an entanglement-breaking channel is $\m E \in C_d, \rho \mapsto \trr{\rho}\1_{d'}/d'$ since its Choi matrix is $C_{\m E} = \1_d \otimes \1_{d'}$.
\begin{ex}\label{ex:depol-eb}
    The depolarizing channel $\Phi_p$ is entanglement-breaking if and only if $p\geq d/(d+1)$.
\end{ex}
\begin{proof}
The Choi matrix of the depolarizing channel is isotropic, for which separability is equivalent to the PPT condition. Parameterizing the channel as $\Phi_p(\rho) = (1-p)\rho + p \text{Tr}[\rho]\frac{\1}{d}$, its unnormalized Choi matrix evaluates to:
\begin{align*}
    C_{\Phi_p} &= (1-p)\ket{\Omega}\bra{\Omega} + p \left( \1 \otimes \frac{\1}{d} \right)
\end{align*}
Applying the partial transpose transforms $\ket{\Omega}\bra{\Omega}$ into the swap operator $\mathbb{F}$. The minimum eigenvalue of the swap operator is strictly $-1$ (occurring on the antisymmetric subspace). Enforcing the PPT condition dictates that the minimum eigenvalue of the fully transposed Choi matrix must remain non-negative:
\begin{align*}
    -(1-p) + \frac{p}{d} \ge 0 \Leftrightarrow
    p \left(1 + \frac{1}{d}\right) \ge 1 \Leftrightarrow
    p \ge \frac{d}{d+1}
\end{align*}
This directly yields the exact parameter regime where the channel becomes entanglement-breaking.
\end{proof}

The Choi matrix of a channel may have a positive partial transpose, but not necessarily be separable. The corresponding channels are called completely copositive or also PPT channels sometimes:
\begin{definition}[Complete copositivity/PPT channel]
    Let $\m E \in \m C(\C^d,\C^{d'}) $ be a quantum channel. If its Choi matrix is PPT, $(\id_d \otimes \vartheta_{d'})(C_T)\geq 0$, then the channel is called completely copositive. The set of all completely copositive channels will be denoted by PPT$(\C^d,\C^{d'})$ and PPT$_d\coloneqq \text{PPT}(\C^d)\coloneqq \text{PPT}(\C^d,\C^{d})$.
\end{definition}
This is equivalent to $(\id \otimes \m E)(\rho)$ being a PPT state again, for every input.
\begin{prop}
    The set of entanglement-breaking channels EB$_d$ is closed under composition with other entanglement-breaking quantum channels, as well as with other arbitrary channels, making it an ideal in the semigroup of quantum channels. The same is true for the set of PPT channels.
\end{prop}
\begin{proof}Let $\rho \in \m D(\C^d)$ be an aribrary state, then
    $(\id_d \otimes \Phi \circ \Psi)(\rho) =(\id_d \otimes \Phi)((\id_d \otimes \Psi)(\rho))$ is separable since $\Phi$ is entanglement-breaking. On the other hand, $(\id_d \otimes \Psi \circ \Phi)(\rho)$ is separable since product channels like $\id_d \otimes \Psi$ preserve separability. Similar reasoning proves the statement about PPT channels.
\end{proof}

The following conjecture was stated in \cite{christandl2012ppt}:
\begin{conjecture}[PPT$^2$ conjecture]. For any PPT quantum channel $\m E \in \text{PPT}_d$, its square is entanglement-breaking: $\m E^2 \in \text{EB}_d$.
\end{conjecture}
The PPT$^n$ conjecture is accordingly about $n$-time compositions. Of course, the trace-preservation of $\m E$ is irrelevant for the conjecture, so that it may be equivalently formulated for completely positive maps.



There is the following simple necessary channel for a channel to be entanglement-breaking:
\begin{lemma}\label{lem:rank}
    Let $\m E \in \m C_d$ be a quantum channel. If $\rk(C_{\m E}) < d$, then $C_T$ is entangled.
\end{lemma}
\begin{proof}
    Note that $C_{\m E}^\text{A} = \tr_\text{B}[C_{\m E}] = \1$, so that $\rk(C_{\m E}^\text{A}) = d$. Thus by proposition \ref{prop:rank-condition} it is $\rk(C_T) > d$.
\end{proof}
\begin{prop}[Necessary Kraus rank, theorem 6 in \cite{horodecki1999general}]
    Let $\m E \in \m C_d$. If $\m E$ can be written with fewer than $d$ Kraus operators, i.e., $\rk(C_{\m E})<d$, then $\m E$ is not entanglement-breaking.
\end{prop}
\begin{proof}
    The minimal number of required Kraus operators is the rank of the Choi matrix. Now use lemma \ref{lem:rank}.
\end{proof}
Unfortunately, there is no simple closed formula for the rank of the Choi matrix of a channel composition. Thus, this result does not provide an easy condition for the PPT$^n$ conjecture, for example.
\begin{prop}
    Let $\m E \in \m C_d$ be a quantum channel. If $\rk(C_{\m E}) = d$, then the channel $\m E$ is entanglement-breaking if and only if it is PPT.
\end{prop}
\begin{proof}
 Using proposition \ref{prop:low-rank} for the Choi matrix of the channel yields the statement.
\end{proof}

\subsubsection{Entanglement-Breaking Indices and Times}
Since entanglement-breaking channels are closed under composition with other quantum channels, in particular themselves, it may be interesting to investigate the smallest number of compositions needed for a channel to become entanglement-breaking (if that is the case). This concept was first investigated in \cite{lamiEntanglementBreakingIndices2015}. To take powers of a quantum channel, the channel must clearly be mapping between the same Hilbert space.
\begin{definition}[Entanglement-breaking index and time \cite{lamiEntanglementBreakingIndices2015, lamiEntanglementSavingChannels2016}]
    Let $\m E \in \m C_d$ be a quantum channel and $\{\Phi_t \in \m C_d\}_{t\geq 0}$ a quantum dynamical semigroup. The entanglement-breaking index of the channel and the entanglement-breaking time of the QDS are 
    \begin{align*}
        n_\text{EB}(T) &\coloneqq \min\{n_0 \in \N: \forall n \in \N, n\geq n_0 :\Phi^n \in \text{EB}_d\} \text{ and }\\
        t_\text{EB}(\Phi)&\coloneqq \inf\{t\geq 0: \forall s \geq t : \Phi_s \in \text{EB}_d \},
    \end{align*}
    respectively. If the above sets are empty, the respective entanglement-breaking index or time is set to $\infty$. Channels (QDSs) with a finite entanglement-breaking index (time) are called eventually entanglement-breaking. The set of eventually entanglement-breaking channels $\m E\in \m C_d$ will be denoted by EEB$_d$. If a channel is not eventually entanglement-breaking, it is entanglement-saving: $\text{ES}_d\coloneqq \m C_d $\textbackslash $\text{EEB}_d$
    \end{definition}\noindent
Since $t\mapsto \Phi_t$ is (strongly) continuous and EB$_d$ is closed, the entanglement-breaking time is achieved as a minimum. The entanglement-breaking indices of the elements of a QDS are closely linked to its entanglement-breaking time:
\begin{prop}
Let $(\Phi_t)_t$ be a QDS with entanglement-breaking time $t_\Phi$. Let $\delta \geq 0$, then the entanglement-breaking index of 
$\Phi_\delta$ is $n(\Phi_\delta) \in [\lfloor t_\Phi / \delta \rfloor, \lceil t_\Phi / \delta \rceil]$.
\end{prop}
\begin{proof}

Let $\Phi_\delta^n$ be entanglement-breaking, then $n\delta \geq t_\text{EB}(\Phi)$ and thus $ n \geq \lceil {t_\Phi}/{\delta} \rceil$ since $n \in \mathbb{N}$, which provides the desired upper bound for the entanglement-breaking index $n_\text{EB}(\Phi_\delta)$. On the other hand, clearly, 
$t_\text{EB}(\Phi) \geq \lfloor \frac{t_\text{EB}(\Phi)}{\delta} \rfloor \cdot \delta$ holds true.  Consequently, $\Phi_\delta^{\lfloor t_\Phi / \delta \rfloor}$ is entanglement-breaking if and only if ${t_\Phi}/{\delta} \in \mathbb{N} \Rightarrow n_0 \geq \lfloor \frac{t_\Phi}{\delta} \rfloor$.
\end{proof}\noindent
A QDS is, in particular, eventually entanglement-breaking if and only if any and thus all of its elements are.

\begin{ex} The entanglement-breaking index of the depolarizing map $\Phi_p$ is $n(\Phi_p) = \left\lceil - {\ln(d+1)}/{\ln(1-p)} \right\rceil$ for all $p \in\ ]0, 1[\ $, as well as $n(\Phi_0) = \infty$ and $n(\Phi_1) = 1$. 
\end{ex}
\begin{proof}Use example \ref{ex:depol-eb} and assume $p \in\ ]0, 1[\ $:
\begin{align*}
1 - (1-p)^n \ge \frac{d}{d+1} &\Leftrightarrow (1-p)^n \le 1 - \frac{d}{d+1} = \frac{d+1-d}{d+1} = \frac{1}{d+1} \\
&\Leftrightarrow n \ge \log_{1-p} \left(\frac{1}{d+1}\right) = \frac{\ln\left(\frac{1}{d+1}\right)}{\ln(1-p)} = -\frac{\ln(d+1)}{\ln(1-p)}.
\end{align*}
\end{proof}


If a channel $\m E \in \m C_d$ is not eventually entanglement-breaking, one might still consider the accumulation points of the sequence $(\m E^n)_{n\in \N}$.
\begin{definition}[AES channels \cite{lamiEntanglementSavingChannels2016} definition 3]
    Let $\m E \in \m C_d$ be a quantum channel. If all accumulation\footnote{Note that in \cite{lamiEntanglementSavingChannels2016} accumulation points are called limit points.} points of the sequence $(\m E^n)_{n \in \N}$ are entanglement-breaking, then $\m E$ is asymptotically entanglement-breaking. Otherwise, the channel is called asymptotically entanglement-saving. The set of asymptotically entanglement-breaking and entanglement-saving channels will be denoted by AEB$_d$ and AES$_d$, respectively.
\end{definition}
Interestingly, it is proven in proposition 29 of \cite{lamiEntanglementSavingChannels2016} that either all or none of the accumulation points are entanglement-breaking. The broad proof strategy is as follows: i) the set of all accumulation points forms a compact group, i.e., they are in particular invertible ii) by inverting the entanglement-breaking accumulation point, all other points can be written as a composition of this particular entanglement-breaking accumulation point and another auxiliary channel iii) use the ideal property of entanglement-breaking channels.

Let it be noted that in \cite{lamiEntanglementSavingChannels2016} they also introduce the notion of universal entanglement-preserving channels (UEP), which are those channels $\m E$ for which $(\id \otimes \m E)(\rho)$ is entangled whenever $\rho$ is. Such UEP channels are shown to be equivalent to the unitary channels. Unitary channels are addiionally entanglement-saving.

\subsubsection{Primitive channels}
Primitive quantum channels provide an example of channels that are eventually entanglement-breaking. This is a well-known result, see e.g. theorem 2.1 \cite{hansonEventuallyEntanglementBreaking2020} for a stronger result, where one direction is proven similarly to the present work.

If $\m E \in C_d$ is a primitive quantum channel, then by definition \ref{def:primitive-etc} it is 
\begin{align}
    \tilde{ \m E }\coloneqq \lim_{n \rightarrow \infty}\m E^n = \m R_\sigma \text{ where } \m D(\C^d)\ni\sigma > 0.
\end{align}
Hence, $C_{\tilde {\m E}}= \sum_{ij}\ket{i} \bra j \otimes \delta_{ij}\sigma = \1 \otimes \sigma$ is a product operator, so the limit $\tilde  {\m E}$ is an entanglement-breaking channel. In principle, the proof that primitive channels are eventually entanglement-breaking works by showing that the limit lies in the interior of EB$_d$ so that there must exist finite powers that come arbitrarily close to $\tilde{\m E}$ and thus also lie in $EB_d$.
\begin{lemma}
    A product operator $O = O_A \otimes O_B$ is in the interior of the set of separable operators if and only if $O_A$ and $O_B$ have full-rank:
    \begin{align*}
        O_A \otimes O_B \in \inter(\sep_{d_A d_B}) \Leftrightarrow (O_A > 0 \wedge O_B > 0)
    \end{align*}
\end{lemma}
\begin{proof}
One needs to show that there exists $\varepsilon > 0 \text{ such that } B_\varepsilon(O_A \otimes O_B) \subset \text{SEP}$. Obviously $\text{int}(\text{SEP}) \subset \text{int}(\mathcal{D}(\mathcal{H})) = \{ \text{full rank states} \}$. Hence, if $O_A \otimes O_B \in \text{int}(\text{SEP})$, then in particular $
    O_A \otimes O_B > 0$. Given the spectral decompositions $O_A = \sum_i \lambda_i |\lambda_i\rangle\langle\lambda_i|$ and $O_B = \sum_j \mu_j |\mu_j\rangle\langle\mu_j|$, one has 
$O_A \otimes O_B = \sum_{i,j} \lambda_i \mu_j (|\lambda_i\rangle \otimes |\mu_j\rangle)(\langle\lambda_i| \otimes \langle\mu_j|) > 0$, so that $O_A,O_B > 0$ follows from $O_A \otimes O_B > 0$. Hence, $(\Rightarrow)$ is done. Now assume that $O_A, O_B > 0$. Clearly, the identity on a bipartite system is a product matrix. Define the map
\begin{align*}
    \Lambda : \mathcal{B}(\mathcal{H}_{AB}) \to \mathcal{B}(\mathcal{H}_{AB}), \quad X \mapsto (O_A \otimes O_B)^{1/2} X (O_A \otimes O_B)^{1/2},
\end{align*}
such that $\Lambda(\1) = O_A \otimes O_B = O$. Consider the restriction $\Lambda|_{\text{SEP}}$; then $\text{Im}(\Lambda|_{\text{SEP}}) \subset \text{SEP}$ because the map is linear and preserves tensor products. By theorem 1 in \cite{gurvitsLargestSeparableBalls2002}, the identity matrix $\1$ lies in the interior of the separable cone, meaning there exists an open neighborhood $U \subset \text{SEP}$ containing $\1$. Furthermore, because $O_A, O_B > 0$, their square roots are uniquely defined and strictly invertible. This makes $\Lambda$ an invertible linear map on a finite-dimensional vector space, which guarantees that it is an open map. Consequently, $\Lambda$ maps the open set $U$ to an open set $\Lambda(U)$ containing $\Lambda(\1) = O$. Since $\Lambda(U) \subset \text{SEP}$, there exists an open neighborhood around $O$ that is entirely contained within the separable set, proving that $O \in \text{int}(\text{SEP})$.

\begin{prop}
    Primitive channels are eventually entanglement-breaking. \label{prop:primitive-eeb}
\end{prop}
\begin{proof}
    Since the Choi matrix $C_{\tilde{\m E}}$ of the limit of a primitive channel $\m E$ is a product matrix, it lies in $\inter(\sep_d)$, which is an open subset of $\sep_d$. Thus, there is an $\epsilon_{\m E} > 0$, such that $B_{\epsilon} (C_{\tilde {\m E}})\in \sep_d$. By the definition of convergence, there is $N(\epsilon) \in \N$ such that for all $n> N$ it is $\norm{\m E^n- \tilde{\m E}}< \epsilon$. The Choi mapping $\m C: \m E \mapsto C_{\m E}$ is linear map, thus bounded continuous. Consequently, one has  $\norm{C_{\m E}-C_{\tilde{\m E}}} = \norm{\m C({\m E}-{\tilde{\m E}})} \leq \norm {\m C} \norm{{\m E}-{\tilde{\m E}}}$. Now, for all $n> N(\epsilon_{\m E}/\norm {\m C})$ it holds that $\norm{C_{\m E^n}-C_{\tilde{\m E}}} < \epsilon_{\m E}$ and thus $C_{\m E^n} \in \sep_d$ and $\m E^n \in EB_d$.
\end{proof}

One potentially useful observation is that to prove the $\text{PPT}^2$ conjecture, one may restrict to a dense set:

\begin{prop}[Corollary 3.2 in \cite{hansonEventuallyEntanglementBreaking2020}] The $\text{PPT}^n$ conjecture holds if and only if every PPT primitive quantum channel $\m E$ has $\m E^n \in \text{EB}_d$.
\end{prop}
\begin{proof} Let $\m E \in \text{PPT}_d$ be a PPT quantum channel. By denseness of primitive quantum channels, there is a sequence $\m E_k \rightarrow \m E$ with each $\m E_k$ being PPT and primitive. By continuity of compositions, it is $\m E^n = \lim_{k \to \infty} \m E_k^n$. Since the set of entanglement-breaking channels is closed, $\m E^n \in \text{EB}_d$.
\end{proof}
By reading the proof, one readily sees that the statement holds for any dense set.

\subsubsection{Other results in the literature}
In the case of quantum dynamical semigroups, also the converse of proposition \ref{prop:primitive-eeb} holds, that is, every eventually entanglement-breaking QDS is primitive. This is shown in theorem 2.1 of \cite{lamiEntanglementSavingChannels2016}. This paper also goes further on to provide tighter bounds for the entanglement-breaking times of quantum dynamical semigroups.

There are many more recent developments on the PPT conjecture.  Notably, Rahaman et al. \cite{rahamanEventuallyEntanglementBreaking2018} rigorously demonstrated that every PPT channel is eventually entanglement-breaking (Theorem 4.4). Their work further characterizes the asymptotic behavior of unital maps: While the identity map serves as a trivial counterexample preventing a universal claim, they proved that almost all unital channels are indeed entanglement-breaking (Theorem 5.2). Additionally, they established that any unital irreducible channel is asymptotically entanglement-breaking (Corollary 6.1).  In \cite{christandlWhenComposedMaps2019} proved the conjecture for dimension $d=3$ using Schmidt number techniques and formulated several equivalent conjectures that provide alternative avenues for understanding this behavior.



  

\chapter{Results}\label{ch:own_results}
\section{SIC-POVMs}\label{sec:sic-povm}
Using a certain type of matrix basis, say $\{O_i\}_i$, any output of a quantum channel $\m E$ can be expanded as $\m E(\cdot) = \sum_i c_i(\cdot) O_i$ for suitable coefficients $c_i$ that depend on the input. Remember that a quantum channel is entanglement-breaking if and only if it can be written in a measure-and-prepare form $\m E (\cdot)= \sum_i \trr{M_i \cdot }\sigma_i$, where the $\sigma_i$ are density matrices and the matrices $M_i$ form a POVM. That is, if it is possible to choose an operator basis with positive semidefinite elements, it is
\begin{align}
    \m E (\cdot)=\sum_i \trr{O_i}c_i(\cdot) \frac{O_i}{\trr{O_i}},
\end{align}
where $O_i/\trr{O_i}$ is a density matrix. Hence, $\m E$ being entanglement-breaking depends entirely on the coefficients\footnote{Well, on what else would it depend...}. Unfortunately, the basis elements being positive semidefinite forces the basis to not be Hilbert-Schmidt orthogonal:
\begin{lemma}[SIC-POVMs, proposition 2.7 in \cite{wolf2012quantum}]\label{lem:sic}
Let $\{P_i\}_{i=1}^n$ be a set of positive semidefinite operators in $\mathcal{M}_d$, where $d \leq n$, which are normalized with respect to the Hilbert-Schmidt scalar product, i.e., $\operatorname{tr}[P_i^2] = 1$. Then
\begin{align*}
    \sum_{i \neq j} |\trr{P_iP_j}|^2 \geq \frac{(n-d)^2 n}{(n-1)d^2}, 
\end{align*}
with equality if and only if all $P_i$ are rank-one projections fulfilling $\trr{P_i P_j} = \frac{n-d}{(n-1)d}$ for all $i \neq j$, and $\sum_i P_i = \frac{n}{d}\1$. In case of equality the $P_i$'s are linearly independent.
\end{lemma}

Consequently, if a basis of positive semidefinite operators were pairwise Hilbert-Schmidt orthogonal, then the left-hand side in the equation above would be zero, but because of $n=d^2$ the right-hand side is equal to $d^2(d-1)/(d+1)$, leading to a contradiction. 

The basis case, where $n=d^2$, that achieves the equality in lemma \ref{lem:sic} leads to the following definition:
\begin{definition}[SIC-POVM]
A SIC-POVM is a POVM $\{M_i\}_{i=1}^{d^2}$, where each effect operator is rank-one (and thus a projection) and fulfills the symmetry condition 
\begin{align*}
    \trr{M_i M_j} = \frac{d \delta_{ij}+1}{d^2(d+1)}
\end{align*}
for all $i,j \in \{1,\dots, d^2\}$.
\end{definition}
That is, if $M_i$ are the effect operators of a SIC-POVM, then $P_i = dM_i$ are the matrices in lemma \ref{lem:sic}. 
The underlying mathematics of SIC-POVMs, that is, sets of complex equiangular lines, were studied back in the 1970s in the context of design theory \cite{Delsarte1975} and brought into the context of quantum mechanics in \cite{Zauner1999}\footnote{Although they were still not called SIC-POVMs.}, where it was conjectured that SIC-POVMs exist in all dimensions (Zauner's conjecture). 

\begin{lemma}\label{lem:sic}[Equation 2.33 in \cite{wolf2012quantum}] Let $\{P_i/d\}_{i=1}^{d^2}$ be a SIC-POVM. For all $X \in \m M_d$ it holds 
\begin{align*}
    X = \frac{1}{d}\sum_{i=1}^{d^2}[(d+1)\trr{X P_i}-\trr{X}]P_i.
\end{align*}
\end{lemma}
\begin{proof}
Since $\{P_i/d\}_i$ forms a basis, one has the following expansion:
\begin{align*}
    X = \sum_{i=1}^{d^2} c_i P_i,
\end{align*}
where $c_i$ are the coefficients. Now, take the Hilbert-Schmidt inner product of both sides with another $P_j$:
\begin{align*}
    \mathrm{tr}[X P_j] = \mathrm{tr}\left[ \left( \sum_{i=1}^{d^2} c_i P_i \right) P_j \right] = \sum_{i=1}^{d^2} c_i \mathrm{tr}[P_i P_j].
\end{align*}
Using that $
    \mathrm{tr}[P_i P_j] = (d\delta_{ij} + 1)/(d+1)$
yields
\begin{align*} \label{eq:trace_pj}
    \mathrm{tr}[X P_j] = \sum_{i=1}^{d^2} c_i  \frac{d\delta_{ij} + 1}{d+1}  = \frac{d}{d+1}c_j + \frac{1}{d+1}\sum_{i=1}^{d^2} c_i.
\end{align*}
Using $\trr{P_i}=1$ and taking the trace of the initial basis expansion of $X$ yields $
    \mathrm{tr}[X] = \mathrm{tr}\left[ \sum_{i=1}^{d^2} c_i P_i \right] = \sum_{i=1}^{d^2} c_i$.
Substitute this back into the equation from before, so that
\begin{align*}
    \mathrm{tr}[X P_j] &= \frac{d}{d+1}c_j + \frac{1}{d+1}\mathrm{tr}[X] \\\Leftrightarrow 
    c_j &= \frac{1}{d} \big[ (d+1)\mathrm{tr}[X P_j] - \mathrm{tr}[X] \big],
\end{align*}
which concludes the proof.
\end{proof}
\begin{prop}[Sufficient condition via SIC-POVMs]\label{prop:sic}
     If one exists, let $\left\{ P_i/d \right\}_{i=1}^{d^2}$ be a SIC-POVM, then the following statement holds true:
     \begin{align*}
        \text{For all }i \in \{1, \dots, d^2\} \text{ it is}(d+1)\mathcal{E}^*(P_i) \ge \1 \Rightarrow \m E\text{ is entanglement-breaking}
     \end{align*}
 If $\m E$ is measure-and-prepare with respect to the $\{d P_i\}_i$ as output states, then this condition is also necessary for entanglement-breaking.
\end{prop}
\begin{proof}
    Using lemma \ref{lem:sic} yields
    \begin{align*}
\mathcal{E}(\rho) &= \frac{1}{d} \sum_i \left[ (d+1) \operatorname{tr}[P_i \mathcal{E}(\rho)] - \operatorname{tr}[\rho] \right] P_i \\
&= \sum_i \operatorname{tr} \left[ \frac{d+1}{d} P_i \mathcal{E}(\rho) - \frac{1}{d} \rho \right] P_i \\
&= \sum_i \operatorname{tr} \Bigg[ \rho \underbrace{\left( \frac{d+1}{d} \mathcal{E}^*(P_i) - \frac{\1}{d} \right)}_{\eqqcolon M_i}\Bigg] P_i.
\end{align*}
If the $M_i$ constitute a POVM, then $\m E$ is in measure-and-prepare form. One may verify that $\sum_i M_i = \1$. Positivity $M_i \geq 0$ is equivalent to $(d+1)\m E^*(P_i) \geq \1$.
\end{proof}

The SIC-POVM condition is in general not necessary, which may be seen by considering the following example: $\m E (\cdot) = \pr 0 \trr{\cdot}$ with adjoint $\m E^*(\cdot) = \trr{\pr {0} \cdot}\1 = \bra{0}\cdot \ket{0}\1$ (which may be verified readily by plugging in the definition). $\m E$ is clearly entanglement-breaking since its Choi matrix is separable. For the Choi matrix of the adjoint one has $C_{\m E^*} = \sum_{ij}\ket{i}\bra{j} \otimes \delta_{i0}\delta_{j0}\1 = \pr 0 \otimes \1$ is separable, so that $C_{\m E}$ is entanglement-breaking. However, the SIC-POVM condition is not fulfilled since $(d+1)\m E^*(P_i) = (d+1)\bra 0 P_i \ket 0 \1 \ngeq \1$. To see the violation of the inequality, assume that it holds. Then $\bra 0 P_i \ket 0 < 1/(d+1) \Leftrightarrow  d = \sum_{i=1}^{d^2}\bra 0 P_i \ket 0 < d^2/(d+1)$, since $\{P_i/d\}_i$ is a POVM.

Note that, for higher dimensions, SIC-POVMs are generally only numerically known and calculated from so-called fiducial vectors, see \cite{Scott2010, Fuchs2017}. A numerical implementation of the SIC-POVM condition \ref{prop:sic} can be found in listing \ref{lst:code-sic}. 

As an example, consider the depolarizing channel again:
\begin{ex}[Depolarizing channel]
Since the depolarizing channel $\Phi_p$ is self-dual, the SIC-POVM condition evaluates to
\begin{align*}
    (d+1)\left( (1-p)P_i + \frac{p}{d}\1 \right) &\geq \1 \Leftrightarrow (d+1)(1-p)P_i \geq  \left(1-p\frac{d+1}{d}\right)\1.
\end{align*}
Since $P_i$ is rank-one, its eigenvalues are one and zero (with multiplicity $d-1$). Hence, the inequality is true if and only if $1-p(d+1)/d \leq 0 \Leftrightarrow p \geq d/(d+1)$, meaning that for the depolarizing channel the condition is also necessary, see \ref{ex:depol-eb}.
\end{ex}



\section{Contraction-Coefficient and Matrix-Norm Criteria}\label{sec:result:contraction}
In the proof that primitive channels are eventually entanglement-breaking, one uses the fact that the maximally mixed state ${\1}/{d}$ lies in the interior of the convex set of separable states. This guarantees the existence of a finite neighborhood, a separable ball of a specific radius with respect to the Hilbert-Schmidt norm, centered at the identity that consists entirely of separable states. 

The primary idea of this chapter is to leverage this geometric structure to derive explicit upper bounds on the entanglement-breaking index $n_{\text{EB}}(\cdot)$. Iterated applications of a quantum channel act as a contraction on the state space. By quantifying this rate of convergence using matrix norms or contraction coefficients, one can bound the maximum distance of any output state to the maximally mixed state after $n$ compositions.

To be more exact, there is the following result:
\begin{prop}[\cite{AliceBobMeet} theorem 9.15]\label{prop:ball}
    Let $\rho \in \m D(\C^d\otimes \C^{d'})$ and $n\coloneqq d_1 d_2$. If $\norm{\rho - \1/n}_2 \leq 1/\sqrt{n(n-1)}$, then $\rho$ is separable.
\end{prop}
\subsubsection{Entanglement-Annihilating}
\begin{prop}[Contraction Coefficient Criterion for Unital Channels]
Let $\mathcal{E} \in \mathcal{C}(\mathcal{H}_A \otimes \mathcal{H}_B)$ be a unital quantum channel, and let $ = d_A d_B$ denote the dimension of the joint Hilbert space. If the contraction coefficient $c_2(\mathcal{E})$ with respect to states satisfies
\begin{align*}
    c_2(\mathcal{E}) &\le \frac{1}{n - 1},
\end{align*}
then $\mathcal{E}$ is entanglement-annihilating.
\end{prop}

\begin{proof}
A channel is entanglement-annihilating if its image is entirely contained within the separable ball around the maximally mixed state. The radius of this separable ball in the Hilbert-Schmidt norm is $r_{\text{sep}} = {1}/{\sqrt{n(n-1)}}$. Thus, one requires $\|\mathcal{E}(\rho) - \frac{\1}{n}\|_2 \le r_{\text{sep}}$ for all states $\rho$. Using unitality and the definition of the contraction coefficient, one can bound the distance of the output state:
\begin{align*}
    \left\| \mathcal{E}(\rho) - \frac{\1}{n} \right\|_2 &= \left\| \mathcal{E}\left(\rho - \frac{\1}{n}\right) \right\|_2 \le c_2(\mathcal{E}) \left\| \rho - \frac{\1}{n} \right\|_2 \le c_2(\mathcal{E}) \sqrt{\frac{n-1}{n}}.
\end{align*}
Enforcing the separability condition by requiring $c_2(\mathcal{E}) \sqrt{\frac{n-1}{n}} \le \frac{1}{\sqrt{n(n-1)}}$ yields exactly $c_2(\mathcal{E}) \le \frac{1}{n-1}$.
\end{proof}

\begin{prop}[Contraction Coefficient Criterion for Arbitrary Channels]\label{prop:ea-contraction}
Let $\mathcal{E} \in \mathcal{C}(\mathcal{H}_A \otimes \mathcal{H}_B)$ be a general quantum channel with $n = d_A d_B$. Let the bias of the channel be defined as $n = \| \mathcal{E}(\frac{\1}{n}) - \frac{\1}{n} \|_2$. If the contraction coefficient satisfies
\begin{align*}
    c_2(\mathcal{E}) &\le \frac{1}{n - 1} - b \sqrt{\frac{n}{n - 1}},
\end{align*}
then $\mathcal{E}$ is entanglement-annihilating.
\end{prop}

\begin{proof}
Applying the triangle inequality by inserting the channel's action on the maximally mixed state:
\begin{align*}
    \left\| \mathcal{E}(\rho) - \frac{\1}{n} \right\|_2 &\le \left\| \mathcal{E}(\rho) - \mathcal{E}\left(\frac{\1}{n}\right) \right\|_2 + \left\| \mathcal{E}\left(\frac{\1}{n}\right) - \frac{\1}{n} \right\|_2 \\
    &\le c_2(\mathcal{E}) \left\| \rho - \frac{\1}{n} \right\|_2 + b \le c_2(\mathcal{E}) \sqrt{\frac{n-1}{n}} + b.
\end{align*}
Requiring this upper bound to be strictly less than or equal to the separable ball radius $r_{\text{sep}} = \frac{1}{\sqrt{D(D-1)}}$, one obtains:
\begin{align*}
    c_2(\mathcal{E}) \sqrt{\frac{n-1}{n}} + b &\le \frac{1}{\sqrt{n(n-1)}} \implies c_2(\mathcal{E}) \le \frac{1}{n-1} - b\sqrt{\frac{n}{n-1}}.
\end{align*}
\end{proof}

\subsubsection{Entanglement-Breaking: Natural Matrix Representation}

The Choi matrix of the power $\m E^n$ of a channel $\m E \in C_d$ is given by $(\id \otimes \m E^n)(\pr{\Omega}) \in \Im
(\id \otimes \m E^n)$. By using proposition \ref{prop:ball} with $d=d'$ one gets the following result: If 
\begin{align}
    \Im(\id \otimes \m E^n) \subset B_{\frac{1}{d\sqrt{d^2-1}},\norm \cdot_2}\left(\frac{\1}{d^2}\right),
\end{align}
then $\m E^n$ is entanglement-breaking.

A channel $\m E$ is entanglement-breaking if and only if $\id \otimes \m E$ is entanglement-annihilating. Thus, one may consider the contraction coefficient $c_2(\id \otimes \m E)$ and apply proposition \ref{prop:ea-contraction}. Unfortunately, this coefficient is always greater than one:
\begin{prop}
For all $ \mathcal{E} \in \mathcal{C}(\mathcal{H})$ one has $ c_2(\text{id} \otimes \mathcal{E}) \ge \max \{1, c_2(\mathcal{E})\}$
\end{prop}
\begin{proof}
It is
$$ c_2(\text{id} \otimes \mathcal{E}) = \sup_{A, B \in \mathcal{{E}}} \frac{\|(\text{id} \otimes \mathcal{E})(A - B)\|_{HS}}{\|A - B\|_{HS}} $$

Choose $A = \1 \otimes A_2$ and $B = \1 / 2 \otimes B_2$, so that $\text{tr}[A_2] = \text{tr}[B_2] = 0$

$$ c_2(\text{id} \otimes \mathcal{E}) \ge \frac{\|\1 / 2 \otimes \mathcal{E}(A_2 - B_2)\|_{HS}}{\|\1 / 2 \otimes (A_2 - B_2)\|_{HS}} = \frac{\|\mathcal{E}(A_2 - B_2)\|_{HS}}{\|A_2 - B_2\|_{HS}} $$

Taking the supremum over all $A$ and $B$ of this form yields $c_2(\text{id} \otimes \mathcal{E}) \ge c_2(\mathcal{E})$.
Quantum channels have a fixed point, so that $\m E(\sigma) = \sigma$ for at least one density matrix $\sigma$. Now let $A\in \m T$, then also $A \otimes \sigma \in \m T$ and $(\id \otimes \m E)(A \otimes \sigma) = A \otimes \sigma \Rightarrow c_2(\id \otimes \m E) \geq 1$.
\end{proof}
\begin{lemma}
    Let $\m E$ be a quantum channel with natural matrix representation 
    \begin{align*}
    \hat{\m E}= \begin{pmatrix}
         1 & 0\\
        v & \Delta
    \end{pmatrix}
    \end{align*}
    as explained in section \ref{sec:matrix-rep}. Then $c_2(\m E)$ is the largest singular value of of $\Delta$: $c_2(\m E) = s_1(\Delta)$
\end{lemma}
\begin{proof} Let $n_{p\rightarrow p}$ denote the contraction coefficient with respect to all matrices, and the Hilbert-Schmidt norm in the underlying spaces. It is $c_2(\m E) = n_{2 \rightarrow 2}(\m E|_{\m T})$, see proposition \ref{prop:contraction-traceless}. The corresponding natural vector representation of an $A \in \m T$ is given by $(0,\vec A)$ for some $\vec A \in \R^{d^2-1}$. Then the matrix representation of $\m E(A)$ is given by
    \begin{align*}
        \begin{pmatrix}
         1 & 0\\
        v & \Delta
    \end{pmatrix}\begin{pmatrix}
        0 \\ \vec A
    \end{pmatrix} = 
    \begin{pmatrix}
        0\\ \Delta \vec A
    \end{pmatrix}.
    \end{align*}
    Hence $\norm{\m E(A-B)}_2 = \norm{\Delta(\vec A-\vec B)}_2$. The proof is concluded by realizing that $n_{2 \rightarrow 2}$ is the largest singular value for said matrix.
\end{proof}
Since proposition \ref{prop:ball} deals with states, it can be convenient to introduce the normalized Choi matrix: $\tilde C_{\m E} \coloneqq C_{\m E}/d$. Note that now $\tilde C_{\m E}$ is a density operator since it has unit trace. The vector representation of the maximally mixed bipartite state $\1/{d^2}$ is given by 
    \begin{align}
        \begin{pmatrix}
            1 \\ 0
        \end{pmatrix},
    \end{align}
    , where the zero is $d^2-1$-dimensional, and since the completely depolarizing channel must map every input to this vector, its matrix representation is given by 
    \begin{align}
    \begin{pmatrix}
        1 & 0\\
        0 & 0
    \end{pmatrix}
    \end{align}
\begin{prop}\label{prop:contraction}
    Let $\m E$ be a quantum channel with natural matrix representation 
    \begin{align*}
    \hat{\m E}= \begin{pmatrix}
         1 & 0\\
        v & \Delta
    \end{pmatrix}.
    \end{align*} If $\norm{v}_2^2 + \norm{\Delta}_F^2 \leq 1/(d^2-1)$, then $\m E$ is entanglement-breaking. For a unital channels, this reduces to $\norm \Delta_\text{HS}^2 \leq 1/(d^2-1)$.
\end{prop}
\begin{proof}
The idea is to again use proposition \ref{prop:ball} on the normalized Choi matrix. Note that $\norm{\tilde C_{\m E} - \1/d^2}_2 = \norm{\hat{\m E}- \Phi_1}_\text{HS}/d$ since $\tilde C_{\Phi_1} = \1/d^2$. One has
\begin{align*}
    \norm{\hat{\m E}- \hat \Phi_1}_\text{HS} ^2 = \norm{\begin{pmatrix}
        0 & 0\\ 
        v & \Delta
    \end{pmatrix}}_\text{HS}^2  = \sum_{i=1}^{d^2-1}|v_i|^2 + \sum_{i,j = 1}^{d^2-1} \vert \Delta_{ij}\vert^2 = \norm{v}_2^2 + \norm \Delta_\text{HS}^2.
\end{align*}
As it is sufficient for entanglement-breaking, set 
\begin{align*}
    \norm{\hat{\m E}- \hat \Phi_1}_\text{HS} ^2 &\leq \frac{1}{d\sqrt{d^2-1}}
    \Leftrightarrow \frac{1}{d} \sqrt{\norm{v}_2^2 + \norm \Delta_\text{HS}^2} &\leq \frac{1}{d\sqrt{d^2-1}} \Leftrightarrow
    \norm{v}_2^2 + \norm \Delta_\text{HS}^2 &\leq \frac{1}{d^2-1},
\end{align*}
which concludes the proof.
\end{proof}

Proposition \ref{prop:contraction} may be used to find bounds on the entanglement-breaking index of the channel. Note that
\begin{align}
\begin{pmatrix}
1 & 0 \\
v_1 & \Delta_1
\end{pmatrix}
\begin{pmatrix}
1 & 0 \\
v_2 & \Delta_2
\end{pmatrix}
=
\begin{pmatrix}
1 & 0 \\
v_1 + \Delta_1 v_2 & \Delta_1 \Delta_2
\end{pmatrix}
\end{align}. 
By induction, this leads to the following result:
\begin{lemma}
    Let $\hat{\mathcal{E}} = \left(\begin{smallmatrix} 1 & 0 \\ v & \Delta \end{smallmatrix}\right)$ be a quantum channel, then 
\begin{align*}
\hat{\mathcal{E}}^n = \begin{pmatrix} 1 & 0^T \\ a & \Delta^n \end{pmatrix}
\end{align*}
with $a := \sum_{k=0}^{n-1} \Delta^k v$\label{lemma:power-matrix}
\end{lemma}
\begin{prop}
    Let $\hat{\mathcal{E}} = \left(\begin{smallmatrix} 1 & 0 \\ v & \Delta \end{smallmatrix}\right)$ be the matrix representation of a quantum channel. If 
    \begin{align*}
        \norm*{\sum_{k=1}^{n-1}\Delta^k v}_2^2 + \norm*{\Delta^n}_\text{HS}^2 \leq \frac{1}{d^2-1},
    \end{align*}
    then $\m E$ is entanglement-breaking.
\end{prop}
\begin{proof}
    Apply lemma \ref{lemma:power-matrix} onto proposition \ref{prop:contraction}.
\end{proof}

% \documentclass{article}
% \usepackage{amsmath}
% \usepackage{amssymb}
% \usepackage{amsthm}
% \usepackage{mathtools}

% % Define the norm command
% \DeclarePairedDelimiter{\norm}{\lVert}{\rVert}

% % Theorem environments
% \newtheorem{prop}{Proposition}
% \theoremstyle{remark}
% \newtheorem*{remark}{Remark}



\begin{prop}
Let $\mathcal{E}$ be a unital quantum channel acting on a $d$-dimensional Hilbert space, and let $\Delta$ denote the unital sub-block of its natural matrix representation, as usual.

i) If the norm of the iterated block satisfies
\begin{align*}
    \|\Delta^n\|_{\text{HS}}^2 &\le \frac{1}{d^2 - 1},
\end{align*}
then the concatenated channel $\mathcal{E}^n$ is entanglement-breaking.

ii) Consequently, if the norm strictly satisfies $\|\Delta\|_{\text{HS}} < 1$ and the index $n$ satisfies
\begin{align*}
    \|\Delta\|_{\text{HS}} \le (d^2 - 1)^{-\frac{1}{2n}} \iff n &\ge - \frac{1}{2} \frac{\ln(d^2 - 1)}{\ln(\|\Delta\|_{\text{HS}})},
\end{align*}
then $\mathcal{E}^n$ is entanglement-breaking.
\end{prop}

\begin{proof}
Statement i) is straightforward by using the procedure explained above. To prove ii), leverage the sub-multiplicativity of the Hilbert-Schmidt norm: $\|\Delta^n\|_{\text{HS}} \le \|\Delta\|_{\text{HS}}^n$. Assuming the initial condition $\|\Delta\|_{\text{HS}} \le (d^2 - 1)^{-\frac{1}{2n}}$, one can exponentiate both sides by $n$ to obtain:
\begin{align*}
    \|\Delta\|_{\text{HS}}^n &\le (d^2 - 1)^{-\frac{1}{2}}.
\end{align*}
Applying the sub-multiplicativity yields:
\begin{align*}
    \|\Delta^n\|_{\text{HS}} &\le \|\Delta\|_{\text{HS}}^n \le (d^2 - 1)^{-\frac{1}{2}}.
\end{align*}
Squaring this result yields:
\begin{align*}
    \|\Delta^n\|_{\text{HS}}^2 &\le \frac{1}{d^2 - 1}.
\end{align*}
By applying the result from part i), one concludes that $\mathcal{E}^n$ is entanglement-breaking.
\end{proof}

\textbf{Remark:} The inequality from part ii) yields a valid bound $n \ge 1$ if and only if the norm falls into the half-open interval $\|\Delta\|_{\text{HS}} \in \left[ \frac{1}{\sqrt{d^2 - 1}}, 1 \right)$. This provides a strict upper bound on the entanglement-breaking index $n_{\text{EB}}(\mathcal{E})$ (the minimum number of iterations required to break all entanglement):
\begin{align*}
    n_{\text{EB}}(\mathcal{E}) &< - \frac{1}{2} \frac{\ln(d^2 - 1)}{\ln(\|\Delta\|_{\text{HS}})}.
\end{align*}

Furthermore, because $\mathcal{E}$ is a unital channel, its affine translation vector is zero, establishing the norm equivalence:
\begin{align*}
    \|\mathcal{E}\|_{\text{HS}}^2 &= 1 + \|\Delta\|_{\text{HS}}^2.
\end{align*}
Substituting $\|\Delta\|_{\text{HS}} = \sqrt{\|\mathcal{E}\|_{\text{HS}}^2 - 1}$ allows expressing the index bound entirely in terms of the channel's total norm:
\begin{align*}
    n(\mathcal{E}) &< - \frac{1}{2} \frac{\ln(d^2 - 1)}{\ln \left( \sqrt{\|\mathcal{E}\|_{\text{HS}}^2 - 1} \right)}.
\end{align*}
This expression holds provided that the total norm strictly falls into the half-open interval:
\begin{align*}
    \|\mathcal{E}\|_{\text{HS}} &\in \left[ \frac{d}{\sqrt{d^2 - 1}}, \sqrt{2} \right).
\end{align*}

\begin{prop}
For every unital quantum channel $\mathcal{\m E}$, the action of the channel is a contraction with respect to the Hilbert-Schmidt norm. That is, $
\norm*{\mathcal{\m E}(X)}_{HS} \leq \norm*{X}_{HS}$ for all $X \in M_d$.
\end{prop}

\begin{proof}
The proof relies on the Kadison-Schwarz inequality (see for example [Wol12]), which holds for any completely positive unital map. For any operator $X$, the inequality states:
\begin{align*}
    \mathcal{T}(X^* X) &\ge \mathcal{T}(X)^* \mathcal{T}(X)
\end{align*}
Because the trace is a positive linear functional, applying the trace to both sides preserves the operator inequality. Separately, because quantum channels are trace-preserving maps, the left-hand side simplifies. By applying the definition of the Hilbert-Schmidt norm ($\|A\|_{\text{HS}}^2 = \text{Tr}[A^* A]$), one obtains:
\begin{align*}
    \text{Tr}[\mathcal{T}(X^* X)] = \text{Tr}[X^* X] = \|X\|_{\text{HS}}^2 &\ge \text{Tr}[\mathcal{T}(X)^* \mathcal{T}(X)] = \|\mathcal{T}(X)\|_{\text{HS}}^2
\end{align*}
Taking the square root of both sides concludes the proof.
\end{proof}

\begin{prop}
Let $\hat{\mathcal{E}} = \begin{pmatrix} 1 & 0 \\ 0 & \Delta \end{pmatrix}$ be the natural matrix representation of a unital quantum channel. Then the largest singular value of the matrix $\Delta$ is bounded by 1: $s_1(\Delta) \leq 1$.
\end{prop}
\begin{proof}
Let $X$ be a traceless Hermitian operator. In the matrix representation, its vector form is simply $(0, \vec{x})^T$, and its Hilbert-Schmidt norm corresponds exactly to the Euclidean norm of this vector, $\norm*{X}_{HS} = \norm*{\vec{x}}_2$. Applying the channel $\hat{\mathcal{E}}$ yields the output vector $(0, \Delta \vec{x})^T$. By Proposition 1, every unital channel is a contraction in the Hilbert-Schmidt norm, meaning the norm of the output cannot exceed the norm of the input: $
\norm*{\Delta \vec{x}}_2 \leq \norm*{\vec{x}}_2$.
Recall that the largest singular value $s_1(\Delta)$ is identical to the induced 2-norm of the matrix. Taking the supremum over all non-zero vectors $\vec{x}$ yields:
\begin{align*}
s_1(\Delta) = \sup_{\vec{x} \neq 0} \frac{\norm*{\Delta \vec{x}}_2}{\norm*{\vec{x}}_2} \leq 1.
\end{align*}
\end{proof}



\subsubsection{Entanglement-Breaking$: \chi$-Matrix}
Remember that proposition \ref{prop:norm-chi} provides a lower bound of $1$ and an upper bound of $d$ for the norm of the matrix representation of a quantum channel.
\begin{prop}
Let $\mathcal{E} \in C_d$ be a quantum channel. Then there is the following necessary condition for entanglement-breaking:
\begin{align*}
    \m E \text{ is entanglement-breaking} \Rightarrow \norm \chi_\text{HS} \leq \sqrt{d}
\end{align*}
Furthermore, there exists an entanglement-breaking channel that achieves the general lower bound $\|\chi\|_{\text{HS}} = 1$. Also, there is a completely co-positive channel that achieves the upper bound $\|\chi\|_{\text{HS}} = \sqrt{d}$ that is necessary for entanglement-breaking channels. 
\end{prop}

\begin{proof}
Assume $\m E$ to be entanglement-breaking. Then its Kraus operators $K_i \eqqcolon |w_i\rangle\langle v_i|$ are rank-one. Note that in this case the $\pr {w_i}$ are the 'prepare' states in the measure-and-prepare form, so that $\norm{w_i}=1$ for all $i$. Using
\begin{align*}
\text{tr}[K_i^* K_j] &= \text{tr}[|v_i\rangle\langle w_i|w_j\rangle\langle v_j|] = \langle w_i, w_j \rangle \langle v_j, v_i \rangle
\end{align*}
together with proposition \ref{prop:relation-norms} yields
\begin{align*}
\|\chi\|_{\text{HS}}^2 &= \sum_{i,j} |\text{tr}[K_i^* K_j]|^2 = \sum_{i,j} |\braket{ w_i}{w_j} |^2 |\braket{v_j}{ v_i }\rangle|^2 \\
&\le \sum_{i,j} \|w_i\|^2 \|w_j\|^2 |\braket{ v_j}{ v_i} |^2 = \sum_{ij}\vert \braket{v_j}{v_i}\vert ^2,
\end{align*} 
where the inequality follows from the Cauchy-Schwarz inequality. Trace-preservation of $\m E$ is in-forced by $\sum_i K_i^*K_i = \1$ which implies $\sum_i \ket{v_i}\braket{w_i}{w_i}\bra{v_i} = \sum_i \pr{v_i} = \1$ since $\{\pr {v_i}\}_i$ needs to be a POVM, so that 
\begin{align*}
    \norm \chi_\text{HS}^2 &\leq \sum_{i,j} \braket{v_j}{v_i} \braket{v_i}{v_j} = \sum_j \bra{v_j}, \left( \sum_i |v_i\rangle\langle v_i| \right) \ket{v_j } \\ &= \sum_j \braket{ v_j}{v_j}  = \text{tr}\left[ \sum_j |v_j\rangle\langle v_j| \right] = \text{tr}[\1] = d, 
\end{align*}
Taking the square root yields the upper bound $\|\chi\|_{\text{HS}} \le \sqrt{d}$. For saturation of the general lower bound with an entanglement-breaking channel consider the completely depolarizing map $ \mathcal{R}_{\1/d}$. Using the general formula for the norm of depolarizing channels yields
\begin{align*}
\|\chi\|_{\text{HS}} &= d \sqrt{1 - 1 \cdot \frac{d^2-1}{d^2}} = d \sqrt{\frac{d^2 - d^2 + 1}{d^2}} = d \sqrt{\frac{1}{d^2}} = 1
\end{align*}
which confirms the existence of an entanglement-breaking channel with $\|\chi\|_{\text{HS}} = 1$. Now onto the saturation of the upper bound with a PPT channel: 
\end{proof}

\begin{prop}
Let $\mathcal{E}$ be a quantum channel and $\vec{\lambda}$ a vector of its eigenvalues. If $
    \|\vec{\lambda}\|_{2n} > d^{\frac{1}{2n}}$, 
then $\mathcal{E}^n$ is not entanglement-breaking. Let $g$ be the number of peripheral eigenvalues, $\lambda_2$ be an eigenvalue with the largest modulus strictly smaller than 1, and $D$ be its algebraic multiplicity. If $
    g + D|\lambda_2|^{2n} > d$, 
then $\mathcal{E}^n$ is not entanglement-breaking. Consequently, if the iteration index satisfies
\begin{align*}
    n &< \frac{\ln(d - g) - \ln(D)}{2\ln(|\lambda_2|)},
\end{align*}
then $\mathcal{E}^n$ is not entanglement-breaking. Furthermore, if $d > g$, the entanglement-breaking index is bounded from below by
\begin{align*}
    n_{\text{EB}}(\mathcal{E}) &\ge \frac{\ln\left(\frac{d-g}{D}\right)}{2\ln(|\lambda_2|)} = \frac{\ln\left(\frac{D}{d-g}\right)}{2\ln(1/|\lambda_2|)}.
\end{align*}
This result requires $D > d - g$ to yield a nontrivial (positive) bound.
\end{prop}

\begin{proof}
By Proposition 6.2.10, if $\mathcal{E}^n$ is entanglement-breaking, the Hilbert-Schmidt norm of its Choi matrix must satisfy $\|\chi_{\mathcal{E}^n}\|_{\text{HS}} \le \sqrt{d}$. Let $\hat{\mathcal{E}}$ denote the natural representation (transfer matrix) of the channel. Because the elements of the Choi matrix are a permutation of the elements of the transfer matrix, their Hilbert-Schmidt norms are identical. Thus, the necessary condition becomes $\|\hat{\mathcal{E}}^n\|_{\text{HS}} \le \sqrt{d}$. By Schur's inequality, the squared Hilbert-Schmidt norm is bounded from below by the sum of the squared moduli of the eigenvalues:
\begin{align*}
    \|\hat{\mathcal{E}}^n\|_{\text{HS}}^2 &\ge \sum_i |\lambda_i(\mathcal{E})|^{2n}.
\end{align*}
Let $\vec{\lambda}$ be the vector of eigenvalues of $\mathcal{E}$. The condition for strict violation becomes $\|\vec{\lambda}\|_{2n}^{2n} > d$. Separating the $g$ peripheral eigenvalues (modulus 1) from the rest, the sum can be lower-bounded:
\begin{align*}
    \sum_i |\lambda_i|^{2n} &= g \cdot 1 + \sum_{i: |\lambda_i|<1} |\lambda_i|^{2n} \ge g + D|\lambda_2|^{2n}.
\end{align*}
Enforcing the strict violation condition $g + D|\lambda_2|^{2n} > d$ yields:
\begin{align*}
    D|\lambda_2|^{2n} &> d - g.
\end{align*}
Taking the natural logarithm and isolating $n$ (noting that $\ln(|\lambda_2|)$ is negative, which flips the inequality) results in:
\begin{align*}
    n &< \frac{\ln(d - g) - \ln(D)}{2\ln(|\lambda_2|)}.
\end{align*}
Therefore, if this inequality holds, the necessary norm condition is strictly violated, meaning $\mathcal{E}^n$ cannot be entanglement-breaking.
\end{proof}
\begin{prop}
Let $g$ be the number of peripheral eigenvalues. If $g > d$, then the quantum channel $\mathcal{E}$ is not eventually entanglement-breaking (EEB).
\end{prop}

\begin{proof}
As established previously, a necessary condition for any iterated channel $\mathcal{E}^n$ to be entanglement-breaking is that the Hilbert-Schmidt norm of its transfer matrix satisfies the upper bound $\|\hat{\mathcal{E}}^n\|_{\text{HS}}^2 \le d$. By Schur's inequality, the squared norm is bounded from below by the sum of the squared moduli of its eigenvalues. Isolating the $g$ peripheral eigenvalues, which have a modulus of exactly 1, yields the following lower bound for all $n \in \mathbb{N}$:
\begin{align*}
    \|\hat{\mathcal{E}}^n\|_{\text{HS}}^2 &\ge \sum_i |\lambda_i|^{2n} \ge g \cdot 1^{2n} = g.
\end{align*}
By assumption, $g > d$. Therefore, the necessary upper bound is strictly violated for all possible powers $n$: $\|\hat{\mathcal{E}}^n\|_{\text{HS}}^2 \ge g > d.$
\end{proof}







\section{See-Saw SDP optimization: Testing the PPT conjecture}\label{ch:result-numerics}
The aim of this section is to test the PPT$^2$ conjecture, i.e., checking if for two completely copositive quantum channels $\m E_1, \m E_2 \in C_d$, the Choi matrix $C_{\m E_1\circ \m E_2}$ of their composition is separable. As explained in chapter \ref{ch:entanglement}, one can do this by investigating whether $(\id \otimes \phi)(C_{\m E_1\circ \m E_2})$ has a negative eigenvalue, i.e., if there is a vector $\ket \psi$ such that $\bra \psi (id \otimes \phi)(C_{\m E_1\circ \m E_2}) \ket \psi < 0$. Here, $\phi:\m M_d \rightarrow \m M_d$ may be any positive, but not completely positive map. By the spectral decomposition, one may equivalently minimize the expression
\begin{align}
    \min_{\mathcal{E}_1, \mathcal{E}_2, \rho} \tr \Big[ \rho \, (\text{id} \otimes \phi)(C_{\mathcal{E}_1 \circ \mathcal{E}_2}) \Big],
\end{align}
where the minimization is over density matrices $\rho$ and PPT channels $\m E_1$ and $\m E_2$, and hope for a negative value:
\begin{enumerate}
    \item Fix a starting density matrix $\rho^{(1)}$ and PPT channel $\m E_1^{(1)}\in \text{PPT}_d$ and optimize over
    \begin{align*}
        \m E_2^{(2)}\coloneqq \argmin_{\m E_2\in\text{PPT}_d}\tr[\rho^{(1)} (\id \otimes \phi)(C_{\m E_1^{(1)}\circ \m E_2})]
    \end{align*}
    \item Fix $\m E_2^{(2)}$ and optimize over $\m E_1$:
    \begin{align*}
        \m E_1^{(2)}\coloneqq \argmin_{\m E_1\in\text{PPT}_d}\tr[\rho^{(1)} (\id \otimes \phi)(C_{\m E_1\circ \m E_2^{(2)}})]
    \end{align*}
    \item Optimize over the density matrix: 
    \begin{align*}
        \rho^{(2)}\coloneqq \argmin_{\rho \in \m D(\C^d)}\tr[\rho (\id \otimes \phi)(C_{\m E_1^{(2)}\circ \m E_2^{(2)}})]
    \end{align*}
\end{enumerate}
This process may be repeated arbitrarily many times, until the value of $\tr[\rho (\id \otimes \phi)(C_{\m E_1\circ \m E_2})]$ does not change significantly anymore (or sooner for weaker results). Note that, by the Rayleigh-Ritz theorem (\cite{Horn2012} theorem 4.4.2), the trace $\trr{M \rho}$ over the state space is minimized when $\rho$ is the pure state corresponding to the eigenvector of the smallest eigenvalue of $M$.

Fortunately, all three optimization steps may be performed by semidefinite programming (SDP). An SDP program consists of a variable, an objective for that variable and a constraint: The variable needs to be a Hermitian matrix $X$, the objective has to be formulated by minimizing $\trr{XC}$, where $C$ is a constant matrix. Lastly, the constraint on the variable needs to be either linear or positive semidefinite (hence the name). If these conditions are met, then the SDP is guaranteed to find this minimum. See for example \cite{Vandenberghe1996} for a more comprehensive overview.

In the case of the optimization problem for the PPT$^2$ conjecture, the objectives are clearly in the desired trace format. The optimization over density matrices obviously obeys SDP constraints. For the quantum channels, their Choi matrix (which is in linear dependence on the channel) needs to be positive semidefinite and have the appropriate partial trace over the second subsystem (depending on the normalization of the maximally entangled state used in the definition of the Choi matrix). Finally, the channel being PPT is a constraint on the partial trace of the Choi matrix, which is simply a reshuffling of a subsystem and thus linear.

An implementation of the SDP program with the Breuer-Hall map\footnote{As mentioned, there are many other positive maps to choose from. The Breuer-Hall map is a good choice to start since it is indecomposable and have a tunable parameter, as mentiond in chapter \ref{ch:entanglement}.} can be found in \ref{lst:breuer}, as well as its use in testing the PPT$^2$ conjecture in \ref{lst:testing-ppt-2}. Note that there are many different Breuer-Hall maps depending on the specific unitary and skew-symmetric matrix that is used. Unfortunately, another optimization over this unitary is not possible by using SDP programs, since the unitarity constraint, for example, is not linear (and neither SDP). Consequently, it is useful to (manually) change the Breuer-Hall maps for a more exhaustive search.

The author of this thesis executed this code for even dimensions $d=4$ up to $d=12$\footnote{On a standard computer in Python. Using lower-level programming languages could make testing $d=14$ possible, but no significantly higher dimensions due to the computational complexity staying the same.}, where no higher dimensions were possible due to the numerical complexity.
\begin{comment}
\section{\textcolor{red}{Complementary channels}}\label{ch:result-complement}
\begin{prop}
    Let $T\in \m C(\hi_A,\hi_B)$ be a quantum channel with Kraus rank $d_A$ then $T$ is entanglement breaking if and only if there is is a complementary channel $T^c$ that is a Schur map wrt some basis. In this case the Schur basis is the measurement basis in the Holevo form of the quantum channel \textcolor{red}{is that result already known}
\end{prop}

\begin{proof}
    \begin{align*}
    \mathcal{E} &= \sum_{n=1}^{d_A} \langle v_n, \cdot \, v_n \rangle |w_n\rangle\langle w_n| \text{, so } K_n = |w_n\rangle\langle v_n| \\
    (K_n)_{ij} &= (w_n)_i (v_n)_j = (\widetilde{K}_i)_{nj}
\end{align*}

$\mathcal{E} : B(\mathcal{H}_A) \to B(\mathcal{H}_B)$, so we choose $\{ |v_n\rangle / \|v_n\| \}_n$ as a basis for $\mathcal{H}_A$ (since $\sum_n |v_n\rangle\langle v_n| = \mathbb{1}_A$), so $(v_n)_j = \delta_{nj}$.

\begin{align*}
    \widetilde{K}_i &: \mathcal{H}_A \to \mathcal{H}_E \\
    (\widetilde{K}_i)_{nj} &= (w_n)_i \delta_{nj} \quad \text{diagonal } \square
\end{align*}

($\Leftarrow$) Let $\mathcal{E}^c$ be a complement s.t. $\mathcal{E}^c(X) = A * X$ for ($A \ge 0, A_{ii} = 1$), then every Kraus decomp $(\widetilde{K}_i)$ is diagonal.

\[
    (\widetilde{K}_i)_{nj} = \delta_{nj} \cdot c_{in} \quad \text{for some } c_{in} \in \mathbb{C} \text{ wrt a basis } \{|e_k\rangle\}_{k=1}^{d_A=d_E}
\]

Here $n,j \in \{1, \dots, d_A=d_E=\text{Kraus rank}\}$ and $i \in \{1, \dots, d_B\}$.

\[
    \text{Then } (K_n)_{ij} = (\widetilde{K}_i)_{nj} = \underbrace{\delta_{nj} c_{in}}_{= (w_n)_i (v_n)_j} = |w_n\rangle\langle v_n| \text{ wrt basis } \{|e_k\rangle\}
\]

\textit{(Note: Kraus rank is upper bounded by $d_B$, so we may extend number to $d_B$)}
\end{proof}
\begin{prop} Let $\mathcal{E} \in \mathcal{C}(\mathcal{H})$ be a q.c. with Kraus rank $= \dim \mathcal{H} = d$ \\
Then $\exists$ a complement $\mathcal{E}^c$ that is e.b. (and thus all are) \\
iff $\mathcal{E}(X) = \1 * X$ wrt a basis
\end{prop}
\begin{proof}
Firstly complements are unique up to composition with \\
unitary channels $\rightarrow$ can neither create entanglement nor \\
destroy it in $(id \otimes \mathcal{E}) (|\Omega\rangle\langle\Omega|)$

\begin{align*}
    (\Rightarrow) \quad & \mathcal{E}^c \text{ e.b., then } \exists \; \mathcal{E}^c \text{ Schur map, if } \mathcal{E}^c = \mathbb{1} * \cdot \\
    & \text{wrt a basis } \{v_i\}_i \Rightarrow \mathcal{E}^c(X) = (\delta_{ij} X_{ii})_{ij} = \sum_i \langle v_i, X v_i \rangle |v_i\rangle\langle v_i| \\
    & \Rightarrow \widetilde{K}_i = |v_i\rangle\langle v_i| \text{ are Kraus op., } (\widetilde{K}_i)_{jl} = \delta_{ij}\delta_{il} \\
    & \Rightarrow (K_j)_{il} = \delta_{ij}\delta_{il} \Rightarrow K_j = \sum_{i,l} (\delta_{ij} \delta_{il}) |v_i\rangle\langle v_l| \\
    & \hspace{4cm} = |v_j\rangle\langle v_j| = \widetilde{K}_j \\
    & \Rightarrow \mathcal{E}(X) = \mathbb{1} * X
\end{align*}

($\Leftarrow$) Let $\mathcal{E}(X) = \mathbb{1} * X$, then $\mathcal{E}^c = \mathcal{E}$ e.b.
\end{proof}
\end{comment}
\section{Entanglement-Breaking Indices of Sets}\label{ch:result-sets}
It may be interesting to consider sets of quantum channels and calculate an entanglement-breaking index of this set, i.e., the supremum of all indices of channels in that set. Firstly, let it be noted, that even if one considers a sequence of channels inside the set of eventually entanglement-breaking channels EEB$_d$, then a limit of such a sequence must not be entanglement-breaking again. As an example for this, consider the depolarizing channel
\begin{align}
    \Phi_{\frac{1}{n}}= \left(1-\frac{1}{n}\right)\id + \frac{1}{n}R_{\1/d}.
\end{align}
This channel converges to the identity, which is clearly not contained in EEB$_d$. In particular, there is no upper bound on the whole set EEB$_d$:
\begin{lemma}
    The set of entanglement-breaking channels $\text{EB}_d$ is closed
\end{lemma}
\begin{proof}
    $\text{EB}_d$ is the pre-image of the set $\text{SEP}_d$ of separable states under the Choi-Jamilkowski-mapping. This mapping is continuous and $\text{SEP}_d$ is closed.
\end{proof}
\begin{prop}
    There is no $n_\text{EB}(d)$, such that for all $\m E \in \text{EEB}_d$ and for all $n\geq n_\text{EB}(d)$ one has $\m E^n \in \text{EB}_d$.
\end{prop}
\begin{proof}
    Assume there is such a $n_\text{EB}(d)\eqqcolon C$, then $\m E^C \in \text{EB}_d$ for all $ \m E \in \text{EEB}_d$. Consider the set $\{\Phi_{1/n}\}_{n\in \N} \subset \text{EEB}_d$ of depolarizing channels. By assumption it is $\Phi_{1/n}^C \in \text{EB}_d$. Since $\text{EB}_d$ is closed, one has that
    \begin{align*}
        \lim_{n\rightarrow \infty} \Phi_{1/n}^C = \left( \lim_{n\rightarrow \infty} \Phi_{1/n}\right)^C  = \id_{C_d}^C\in \text{EEB}_d,
    \end{align*}
    which is a contradiction. Here, the first '=' is truholds because since the composition is a continuous operation.
\end{proof}
\begin{prop}
    Let $S \subset \text{EEB}_d$ be compact, such that for all $\m E \in S$ there exists $n \in \N$ such that $\m E^n \in \inter(\text{EB})$, then there is an index $n_S(d)$ such that for all $\m E \in S$ it holds $n_\text{EB}(\m E)\leq n_S(d)$.\label{prop:set-index}
\end{prop}
\begin{proof}
    Let $U_n\coloneqq \{\Phi \in S: \Phi^n \in \inter(\text{EB})\}$, then $S \subset \bigcup_{n\in \N}U_n$. Let $f_n:C_d \rightarrow C_d, \m E \mapsto \m E^n$ be the $n$-fold composition (which is continuous), then $U_n = f_n^{-1}(\inter(EB))$ is open since the interior of any set is open. Since $S$ is compact, it has a finite open subcover $S \subset \bigcup_{i=1}^N U_{n_i}$. Hence, $n_S(d) \leq \max_{i=1}^N \{n_i\}< \infty$.
\end{proof}
\begin{ex}
    Let $S \in \text{EEB}_d$ be a compact set, then there might not be an index $n_S$ with the properties from proposition \ref{prop:set-index}. As an example consider the channel $\m R_{\pr 0}$, which is entanglement-breaking, i.e., $n_{EB}(\m R_{\pr 0})=1$. In particular, this channel lies on the boundary of $EB_d$: $\m R_{\pr 0} \in \partial \text{EB}_d$.
\end{ex}
\begin{comment}
\section{Other}
\subsubsection{\textcolor{red}{convolution of heisenberg weyl channel stuff}}

\subsubsection{\textcolor{red}{missing}}

\begin{comment}

\begin{prop}
A quantum channel $\mathcal{E}$ is entanglement-breaking if and only if its natural representation $\hat{\mathcal{E}}$ satisfies:
\begin{align*}
\hat{\mathcal{E}} \in \text{span} \{ P \otimes \bar{P} : \text{rank}(P) = 1 \}
\end{align*}
\end{prop}

\begin{proof}
Let $\mathcal{E}$ be a quantum channel. The natural representation (or Liouville representation) of the channel in terms of its Kraus operators $\{K_i\}$ is universally given by:
\begin{align*}
\hat{\mathcal{E}} &= \sum_i K_i \otimes \bar{K}_i
\end{align*}

First, assume $\mathcal{E}$ is entanglement-breaking. By definition, an entanglement-breaking channel admits a Kraus decomposition where every Kraus operator is rank-one, meaning we can write $K_i = \ket{\psi_i}\bra{\phi_i}$ for some vectors $\ket{\psi_i}$ and $\ket{\phi_i}$. Substituting these into the natural representation yields:
\begin{align*}
\hat{\mathcal{E}} &= \sum_i (\ket{\psi_i}\bra{\phi_i}) \otimes \overline{(\ket{\psi_i}\bra{\phi_i})}
\end{align*}
Since each $K_i$ is a rank-one operator $P_i$, each term in the sum is exactly of the form $P_i \otimes \bar{P}_i$. Because $\hat{\mathcal{E}}$ can be written as a linear combination (specifically, a sum) of such elements, it directly follows that $\hat{\mathcal{E}} \in \text{span} \{ P \otimes \bar{P} : \text{rank}(P) = 1 \}$.

Conversely, assume that $\mathcal{E}$ is a quantum channel and $\hat{\mathcal{E}} \in \text{span} \{ P \otimes \bar{P} : \text{rank}(P) = 1 \}$. This implies that the action of the channel can be expanded using a set of rank-one operators $\{P_j\}$:
\begin{align*}
\hat{\mathcal{E}} &= \sum_j c_j (P_j \otimes \bar{P}_j)
\end{align*}
for some coefficients $c_j$. Because $\mathcal{E}$ is a completely positive map, its associated Choi matrix must be positive semi-definite. The canonical isomorphism between the natural representation and the Choi matrix ensures that the physically valid completely positive decomposition of $\hat{\mathcal{E}}$ corresponds to terms where the coefficients are strictly positive ($c_j > 0$). By absorbing these positive coefficients into the operators as $K_j = \sqrt{c_j} P_j$, we construct a valid Kraus decomposition $\hat{\mathcal{E}} = \sum_j K_j \otimes \bar{K}_j$ where every $K_j$ is still a rank-one operator. Having a Kraus decomposition consisting entirely of rank-one operators is the exact definition of an entanglement-breaking channel.
\end{proof}
\end{comment}
\end{comment}

\chapter{Conclusion and Outlook}
\label{chap:conclusion}
\subsubsection{Conclusion}
This thesis has investigated the concept of entanglement-breaking quantum channels, with a particular focus on their behavior under composition. To establish the necessary framework for analyzing these phenomena, chapter \ref{chap:2_general_preliminaries} provided a brief introduction to and overview of the core concepts of quantum information theory, as well as the basic mathematical foundations needed to study the topic. Chapter \ref{ch:3_preliminaries_channels} subsequently provided a more detailed overview of properties of quantum channels and special types of channels. A particular focus is on the different matrix representations of a quantum channel as a linear map, and recent developments are outlined. Chapter \ref{chap:4_entanglement} introduces the core concepts of (eventually) entanglement-breaking channels and entanglement-breaking indices. Some of the existing literature on is reviewed and recent developments are states. Primitive quantum channels, for example, are eventually entanglement-breaking since their limit of the power series does not lie on the boundary of the set of entanglement-breaking channels.

In chapter \ref{ch:own_results}, some necessary or separately sufficient analytical conditions for a channel to be entanglement-breaking are provided: Using the property, that SIC-POVMs are a basis and the measure-and-prepare form of an entanglement-breaking quantum channel, section 6.1 introduces a sufficient condition for a channel to be entanglement-breaking. It is shown that this criterion, is in general, only sufficient. As an example, the criterion is applied to the well-studied depolarizing channel. 

Section 6.2 derives sufficient conditions for a channel to be entanglement-breaking using its contraction coefficient and the Hilbert-Schmidt norm of its matrix representations. Specifically, it establishes an analytically computable upper bound on a unital channel's entanglement-breaking index, yielding:
\begin{align*}
    n_\text{EB}(\mathcal{E}) &< -\frac{\ln(d^2-1)}{2\ln(||\Delta||_{HS})}.
\end{align*}
Furthermore, a necessary criterion is introduced, proving that the Frobenius norm of an entanglement-breaking channel's x-matrix is strictly upper-bounded by $\sqrt{d}$. Lastly, it is demonstrated that these fundamental norm bounds can be saturated, as the completely depolarizing channel achieves the general lower bound of $1$, while certain completely copositive channels achieve the necessary upper bound of $\sqrt{d}$.

Section 6.3 details a numerical investigation of the PPT$^n$ conjecture using a See-Saw Semidefinite Programming optimization technique. Using the Breuer-Hall map as the positive map, the procedure was implemented to test the PPT$^2$ conjecture. No violations of the conjecture were found for even dimensions up to $d=12$ tested.  
\begin{comment}
\textcolor{red}{Section 6.4 establishes a structural equivalence between an entanglement-breaking channel of full Kraus rank and its complementary channel. Specifically, it proves that the original channel is entanglement-breaking if and only if its complementary channel acts identically to a Schur multiplier map with respect to the associated measurement basis.}
\end{comment}

Section 6.4 extends the concept of the entanglement-breaking index to arbitrary sets of quantum channels. It demonstrates that while the entire space of eventually entanglement-breaking channels lacks a universal, finite upper bound for this index, strict bounds can be rigorously established for specific compact subsets where the iterative compositions eventually enter the interior of the entanglement-breaking set.

\subsubsection{Outlook}
While the conditions on the matrix norm of the quantum channel were presented in chapter \ref{sec:result:contraction}, the thesis mostly fails to apply them to exemplary quantum channels apart from the depolarizing channel. This is something that could be done in the future.

In search of a violation of the PPT$^n$ criterion, the SDP program from section \ref{ch:result-numerics} could be performed with other positive maps, for example with Choi-type maps.

Ultimately, while this work provides some new perspectives on entanglement-breaking of iterated channels, the complete classification of eventually entanglement-breaking operations and the resolution of the PPT$^2$ conjecture remain open challenges.



\appendix

\chapter{Codes}
\label{cha:app1_codes}
\begin{lstlisting}[language=python, caption = {General utility functions}, label = {lst:code-utility}, basicstyle=\ttfamily\scriptsize, % Shrinks the font size
    lineskip=-1pt,                   % Pulls the lines of code closer together
    aboveskip=10pt,                  % Reduces white space above the box
    belowskip=10pt                   % Reduces white space below the box
]
import numpy as np
from scipy.linalg import block_diag
from scipy.linalg import eigvalsh
from toqito.channel_ops import apply_channel
from toqito.matrices import gen_pauli_z as clock, gen_pauli_x as shift

np.set_printoptions(precision=5, suppress=True)
from fid_vectors import fid_vector_list

import jax
import jax.numpy as jnp

import numpy as np

def depol_jax(rho, p):
    d = rho.shape[0]
    return (1-p)* rho + p* (jnp.eye(d) / d) * jnp.trace(rho)

def compute_ptm_jax(channel_func_single_arg, d):
    basis_matrices = get_orthonormal_basis_first_identity(d) # Deine Funktion von vorhin
    basis_stack = jnp.array(basis_matrices)
    
    basis_flat = basis_stack.reshape(basis_stack.shape[0], -1)
    outputs_flat = outputs_stack.reshape(outputs_stack.shape[0], -1)
    
    ptm = jnp.conj(basis_flat) @ outputs_flat.T
    return ptm


def get_orthonormal_basis_first_identity(d: int):
    """
    Generates the d^2 Orthonormal Generalized Gell-Mann matrices.
    IMPORTANT: The normalized Identity is placed at index 0.
    
    Args:
        d (int): Dimension of the Hilbert space.
        
    Returns:
        list[np.ndarray]: List of basis matrices E_0, ..., E_{d^2-1}.
    """
    basis = []

    # E_0 = I / sqrt(d)
    basis.append(np.eye(d, dtype=complex) / np.sqrt(d))
    
    # 2. Symmetric Off-Diagonal
    for k in range(d):
        for j in range(k):
            mat = np.zeros((d, d), dtype=complex)
            mat[j, k] = 1
            mat[k, j] = 1
            basis.append(mat / np.sqrt(2))
            
    # 3. Anti-Symmetric Off-Diagonal
    for k in range(d):
        for j in range(k):
            mat = np.zeros((d, d), dtype=complex)
            mat[j, k] = -1j
            mat[k, j] = 1j
            basis.append(mat / np.sqrt(2))
            
    # 4. Diagonal (Traceless)
    for l in range(1, d):
        mat = np.zeros((d, d), dtype=complex)
        coeff = np.sqrt(1 / (l * (l + 1)))
        for m in range(l):
            mat[m, m] = coeff
        mat[l, l] = -l * coeff
        basis.append(mat)
        
    return basis

def compute_ptm(channel_func, d: int) -> np.ndarray:
    """
    Computes the Pauli Transfer Matrix (PTM) for a black-box channel.
    
    Args:
        channel_func (callable): A function T(rho) -> rho_prime.
        d (int): Dimension of the system.
        
    Returns:
        np.ndarray: The d^2 x d^2 matrix M where M_ij = Tr(E_i @ T(E_j)).
    """
    # 1. Get the basis E_0, ..., E_{d^2-1}
    basis = get_orthonormal_basis_first_identity(d)
    n = len(basis) # Should be d^2
    
    # 2. Initialize the empty PTM
    ptm = np.zeros((n, n), dtype=complex)
    
    # 3. Compute column by column (to minimize calls to channel_func)
    for j in range(n):
        E_j = basis[j]
        
        output_rho = channel_func(E_j)
        
        for i in range(n):
            E_i = basis[i]
            
            # M_ij = Tr(E_i @ T(E_j))
            # Note: Since basis is Hermitian, E_i^dagger = E_i.
            val = np.trace(E_i @ output_rho)
            ptm[i, j] = val
            
    return ptm

def compute_ptm_vectorized(channel_func, d: int) -> np.ndarray:
    """
    Computes the Pauli Transfer Matrix (PTM) using vectorization.
    """
    basis_flat = np.array([b.flatten() for b in basis_matrices])

    output_states = [channel_func(E_j) for E_j in basis_matrices]

    outputs_flat = np.array([rho.flatten() for rho in output_states])

    ptm = basis_flat.conj() @ outputs_flat.T
    
    return ptm

    


def heisenberg_weyl_toqito(x: int, z: int, d: int) -> np.ndarray:
    """
    Implements the Heisenberg-Weyl operator using toqito's generalized Pauli function.
    
    Toqito's `gen_pauli(k, 1, d)` corresponds to X.
    Toqito's `gen_pauli(k, 2, d)` corresponds to Z.
    
    However, toqito does not have a single function taking both x and z directly 
    as a combined W operator in one call, so we construct it similarly.
    """
    try:
        from toqito.matrices import gen_pauli
    except ImportError:
        return "Toqito not installed. Run `pip install toqito`."
    
    X_op = gen_pauli(x, 1, d)
    
    Z_op = gen_pauli(z, 2, d)
    
    return X_op @ Z_op


def depol_channel(rho):
    return apply_channel(rho,choi_depolarizing)

def top_two_abs_values(arr: np.ndarray, decimals: int = 5):
    """
    Analyzes a 1D numpy array to find the counts of the highest 
    and second-highest absolute values, robust to floating point errors.

    Args:
        arr (np.ndarray): Input array of shape (n,).
        decimals (int): Number of decimal places to round to. 
                        Default is 10, which handles most numerical noise.

    Returns:
        dict: Statistics about the top two absolute values.
    """
    # 1. Compute absolute values
    abs_vals = np.abs(arr)
    
    # 2. CRITICAL STEP: Round to handle floating point noise
    # This groups 0.300000001 and 0.299999999 together as 0.3
    abs_vals = np.round(abs_vals, decimals=decimals)
    
    # 3. Find unique values and their counts
    unique_vals, counts = np.unique(abs_vals, return_counts=True)
    
    # 4. Handle Edge Case: Empty Array
    if len(unique_vals) == 0:
        return {
            "first_max_val": None, "first_max_count": 0,
            "second_max_val": None, "second_max_count": 0
        }

    # 5. Extract First Highest (Last element in sorted array)
    first_max_val = unique_vals[-1]
    first_max_count = counts[-1]

    # 6. Extract Second Highest (Second to last element)
    if len(unique_vals) >= 2:
        second_max_val = unique_vals[-2]
        second_max_count = counts[-2]
    else:
        # Edge Case: All numbers have the same absolute value
        second_max_val = None
        second_max_count = 0

    return {
        "first_max_val": first_max_val,
        "first_max_count": first_max_count,
        "second_max_val": second_max_val,
        "second_max_count": second_max_count
    }
\end{lstlisting}

\begin{lstlisting}[language = python, caption = Implementation of the SIC-POVM condition, label = {lst:code-sic}, basicstyle=\ttfamily\scriptsize, % Shrinks the font size
    lineskip=-1pt,                   % Pulls the lines of code closer together
    aboveskip=10pt,                  % Reduces white space above the box
    belowskip=10pt                   % Reduces white space below the box
]

def get_sic_povm(d, fiducial_vector):
    """
    Generates a SIC-POVM from a fiducial vector using Weyl-Heisenberg operators.
    """
    X = shift(d)
    Z = clock(d)
    
    povm_elements = []
    # Displacement operators D_{a,b} = X^a Z^b (ignoring phase factors for POVM)
    for a in range(d):
        for b in range(d):
            # D = (X^a) @ (Z^b)
            op = np.linalg.matrix_power(X, a) @ np.linalg.matrix_power(Z, b)
            vec = op @ fiducial_vector
            # POVM element E_k = 1/d * |psi><psi|
            povm_elements.append((1/d) * np.outer(vec, vec.conj()))
            
    return povm_elements

def sic_condition_jax(channel_func, d):
    """
    channel_func: Deine 'channel_ready' Funktion (lambda), die nur ein Argument (rho) nimmt.
    """

    fid_vector = jnp.array(fid_vector_list[d-2]) 

    fid_vector = fid_vector / jnp.linalg.norm(fid_vector)
    
    sic_povm_np = get_sic_povm(d, fid_vector) 
    sic_povm = [jnp.array(M) for M in sic_povm_np]
    
    temp_sic = True

    for i in sic_povm:
        channel_out = channel_func(i)
        
        temp = (d + 1) * channel_out - jnp.eye(d) / d

        if not is_positive_semidefinite_jax(temp):
            temp_sic = False
            # Optional: break, um Zeit zu sparen, wenn einer fehlschlägt
            break 

    if temp_sic:
        print('SICs say channel is entanglement breaking')
    else:
        print('SICs can not make a statement')
        
    return temp_sic
\end{lstlisting}

\begin{lstlisting}[language = python, caption = {Optimization of one PPT channel with an SDP program using the Breuer-Hall map}, label = {lst:breuer}, basicstyle=\ttfamily\scriptsize, % Shrinks the font size
    lineskip=-1pt,                   % Pulls the lines of code closer together
    aboveskip=10pt,                  % Reduces white space above the box
    belowskip=10pt                   % Reduces white space below the box
]
def optimize_channel(target, J_fixed, rho_fixed, d, U):
    """
    Optimizes either T1 or T2 to minimize Tr(rho * M(T1, T2)).
    
    Parameters:
    - target: String, either 'T1' or 'T2' (the variable to optimize).
    - J_fixed: numpy array (d^2 x d^2), the Choi matrix of the fixed channel.
    - rho_fixed: numpy array (d^2 x d^2), the fixed state (eigenvector).
    - d: Integer, dimension of the subsystem (must be even, >= 4).
    
    Returns:
    - min_val: The minimum expectation value found.
    - J_opt: The Choi matrix of the optimized channel.
    """
    d2 = d**2
    P_reshuffle, P_PT, P_TrB = build_linear_maps(d)
    P_reshuffle_inv = P_reshuffle.T # Inverse of permutation is its transpose
    
    # Define the CVXPY Variable for the target Choi matrix
    J_var = cp.Variable((d2, d2), hermitian=True)
    
    # Convert fixed Choi to Superoperator
    K_fixed = (P_reshuffle @ J_fixed.flatten()).reshape((d2, d2), order = 'C')
    
    # Convert variable Choi to Superoperator
    K_var = cp.reshape(P_reshuffle @ cp.vec(J_var, order='C'), (d2, d2), order = 'C')
    
    # Compose channels (Matrix multiplication of superoperators)
    if target == 'T1':
        # T1 is var, T2 is fixed. K_composed = K_var @ K_fixed
        K_composed = K_var @ K_fixed
    elif target == 'T2':
        # T2 is var, T1 is fixed. K_composed = K_fixed @ K_var
        K_composed = K_fixed @ K_var
    else:
        raise ValueError("Target must be 'T1' or 'T2'")
        
    # Convert composed Superoperator back to Choi matrix
    J_composed = cp.reshape(P_reshuffle_inv @ cp.vec(K_composed, order='C'), (d2, d2))
    
    # --- Apply Breuer-Hall Map to J_composed ---
    #U = get_breuer_hall_U(d)
    I_U = np.kron(np.eye(d), U)
    
    # J_composed is TP, so Tr_B(J) x I = I_d^2
    term1 = np.eye(d2)
    term2 = J_composed
    
    # Apply (I x U) * J_composed^{T_B} * (I x U^dagger)
    J_comp_PT = cp.reshape(P_PT @ cp.vec(J_composed, order='C'), (d2, d2), order = 'C')
    term3 = I_U @ J_comp_PT @ I_U.conj().T
    
    M = term1 - term2 - term3
    
    # --- Objective ---
    objective = cp.Minimize(cp.real(cp.trace(rho_fixed @ M)))
    
    # --- Constraints for J_var to be a valid PPT Channel ---
    constraints = []
    constraints.append(J_var >> 0) # Completely Positive
    
    # Partial Transpose Constraint (PPT)
    J_var_PT = cp.reshape(P_PT @ cp.vec(J_var, order='C'), (d2, d2), order = 'C')
    constraints.append(J_var_PT >> 0)
    
    # Trace Preserving Constraint: Tr_B(J) = I_d
    TrB_J_var = cp.reshape(P_TrB @ cp.vec(J_var, order='C'), (d, d), order = 'C')
    constraints.append(TrB_J_var == np.eye(d))
    
    # --- Solve ---
    prob = cp.Problem(objective, constraints)
    prob.solve(solver=cp.SCS) # SCS or MOSEK recommended for large complex SDPs
    
    return prob.value, J_var.value

\end{lstlisting}

\begin{lstlisting}[language = python, caption = Testing the PPT$^2$ conjecture with SDPs, label={lst:testing-ppt-2}, basicstyle=\ttfamily\scriptsize, % Shrinks the font size
    lineskip=-1pt,                   % Pulls the lines of code closer together
    aboveskip=10pt,                  % Reduces white space above the box
    belowskip=10pt                   % Reduces white space below the box
]
    import numpy as np
import cvxpy as cp
from scipy.optimize import minimize
import autograd.numpy as anp
from autograd import value_and_grad
from scipy.linalg import qr
import scipy.sparse as sp

def get_random_breuer_hall_U(d):
    """
    Generates a random antisymmetric unitary matrix U for the Breuer-Hall map.
    """
    # 1. Construct the trivial antisymmetric unitary (your original Omega)
    Omega = np.zeros((d, d), dtype=complex)
    for i in range(d // 2):
        Omega[2*i, 2*i+1] = 1.0
        Omega[2*i+1, 2*i] = -1.0
        
    # 2. Generate a random complex Gaussian matrix
    G = np.random.randn(d, d) + 1j * np.random.randn(d, d)
    
    # 3. Extract a random unitary matrix V using QR decomposition
    V, _ = qr(G)
    
    # 4. Construct the random antisymmetric unitary
    # CRITICAL: Use V.T (transpose), NOT V.conj().T (Hermitian conjugate)
    U_random = V @ Omega @ V.T
    
    return U_random

def generate_random_density_matrix(dim):
    """
    Generiert eine zufällige Dichtematrix der Dimension (dim x dim).
    """
    G = np.random.randn(dim, dim) + 1j * np.random.randn(dim, dim)
    rho_unnormalized = G @ G.conj().T
    rho = rho_unnormalized / np.trace(rho_unnormalized)
    return rho

def generate_random_boundary_ppt_channel(d, P_PT, P_TrB):
    d2 = d**2

    G = np.random.randn(d2, d2) + 1j * np.random.randn(d2, d2)
    H = G + G.conj().T
    

    C_PT = cp.reshape(P_PT @ cp.vec(C, order='C'), (d2, d2), order='C')
    constraints.append(C_PT >> 0)
    
    TrB_C = cp.reshape(P_TrB @ cp.vec(C, order='C'), (d, d), order='C')
    constraints.append(TrB_C == np.eye(d))
    
    objective = cp.Minimize(cp.real(cp.trace(H @ C)))
    prob = cp.Problem(objective, constraints)
    prob.solve(solver=cp.SCS) 
    
    return C.value


def optimize_U_scalable(rho, J_PT, d, Omega):
    """
    Optimizes the antisymmetric unitary U for any dimension d using the Cayley Transform.
    """
    # Force all constants into autograd's tracked array format to prevent primitive leaks
    I_d = anp.eye(d, dtype=complex)
    Omega_anp = anp.array(Omega, dtype=complex)
    rho_anp = anp.array(rho, dtype=complex)
    J_PT_anp = anp.array(J_PT, dtype=complex)
    
    def objective(params):
        # 1. Map flat array to complex matrix Z
        Z_real = anp.reshape(params[:d**2], (d, d))
        Z_imag = anp.reshape(params[d**2:], (d, d))
        Z = Z_real + 1j * Z_imag
        
        # 2. Map Z to Hermitian H
        H = Z + anp.conj(anp.transpose(Z))
        
        # 3. Map Hermitian H to Unitary V via Cayley Transform (Bypassing expm)
        iH = 1j * H
        V = anp.matmul(I_d - iH, anp.linalg.inv(I_d + iH))
        
        # 4. Map V to Antisymmetric Unitary U
        U = anp.matmul(V, anp.matmul(Omega_anp, anp.transpose(V)))
        
        # 5. Construct tensor operator (autograd safe)
        I_U = anp.reshape(anp.einsum('ij,kl->ikjl', I_d, U), (d**2, d**2))
        I_U_dag = anp.conj(anp.transpose(I_U))
        
        # 6. Calculate objective trace
        val = anp.real(anp.trace(anp.matmul(rho_anp, anp.matmul(I_U, anp.matmul(J_PT_anp, I_U_dag)))))
        return -val

    x0 = np.random.randn(2 * d**2)
    
    res = minimize(value_and_grad(objective), x0, jac=True, method='BFGS')
    
    # Reconstruct optimal U purely in NumPy
    Z_real_opt = res.x[:d**2].reshape((d, d))
    Z_imag_opt = res.x[d**2:].reshape((d, d))
    Z_opt = Z_real_opt + 1j * Z_imag_opt
    H_opt = Z_opt + Z_opt.conj().T
    
    # Cayley Transform for the final output
    I_np = np.eye(d, dtype=complex)
    iH_opt = 1j * H_opt
    V_opt = (I_np - iH_opt) @ np.linalg.inv(I_np + iH_opt)
    U_opt = V_opt @ Omega @ V_opt.T
    
    return U_opt

def update_rho(M_matrix):
    """
    Updates the variational state rho by finding the eigenvector 
    corresponding to the smallest eigenvalue of M.
    
    Parameters:
    - M_matrix: A (d^2 x d^2) numpy array representing the Hermitian matrix M.
    
    Returns:
    - rho_new: The updated (d^2 x d^2) density matrix.
    - min_eval: The smallest eigenvalue (for tracking progress).
    """
    # 1. Calculate eigenvalues (w) and eigenvectors (v)
    # eigh guarantees w is sorted ascending, and v[:, i] is the i-th eigenvector
    w, v = np.linalg.eigh(M_matrix)
    
    # 2. Extract the smallest eigenvalue and its eigenvector
    min_eval = w[0]
    psi_min = v[:, 0]  # The column vector corresponding to w[0]
    
    # 3. Construct the new density matrix: rho = |psi_min><psi_min|
    # Using outer product of the vector and its complex conjugate
    rho_new = np.outer(psi_min, psi_min.conj())
    
    return rho_new, min_eval


def build_linear_maps(d):
    """
    Constructs the permutation matrices using ultra-fast, memory-safe 
    vectorized operations in COO format. No nested loops.
    """
    d2 = d**2
    d4 = d**4
    
    # -------------------------------------------------------------
    # 1. P_reshuffle (K_{kl, ij} = J_{ik, jl})
    # -------------------------------------------------------------
    # Generate all combinations of (i, j, k, l)
    i, j, k, l = np.mgrid[0:d, 0:d, 0:d, 0:d]
    
    # Flatten them to 1D arrays
    i = i.flatten()
    j = j.flatten()
    k = k.flatten()
    l = l.flatten()
    
    # Calculate row and column indices
    row_K = (k*d + l)*d2 + (i*d + j)
    col_J = (i*d + k)*d2 + (j*d + l)
    vals_K = np.ones(d4)
    
    # Build sparse matrix directly
    P_reshuffle = sp.coo_matrix((vals_K, (row_K, col_J)), shape=(d4, d4)).tocsr()
    
    # -------------------------------------------------------------
    # 2. P_PT (Partial Transpose: J^{T_B}_{ik, jl} = J_{il, jk})
    # -------------------------------------------------------------
    # The indices i, j, k, l are already defined above
    row_PT = (i*d + k)*d2 + (j*d + l)
    col_PT = (i*d + l)*d2 + (j*d + k)
    vals_PT = np.ones(d4)
    
    P_PT = sp.coo_matrix((vals_PT, (row_PT, col_PT)), shape=(d4, d4)).tocsr()
    
    # -------------------------------------------------------------
    # 3. P_TrB (Partial Trace on B: Tr_B(J)_{i,j} = sum_k J_{ik, jk})
    # -------------------------------------------------------------
    # For Partial Trace, we only iterate over i, j, k
    i_tr, j_tr, k_tr = np.mgrid[0:d, 0:d, 0:d]
    
    i_tr = i_tr.flatten()
    j_tr = j_tr.flatten()
    k_tr = k_tr.flatten()
    
    row_TrB = i_tr*d + j_tr
    col_TrB = (i_tr*d + k_tr)*d2 + (j_tr*d + k_tr)
    vals_TrB = np.ones(d**3) # Only d^3 elements
    
    P_TrB = sp.coo_matrix((vals_TrB, (row_TrB, col_TrB)), shape=(d2, d4)).tocsr()
    
    return P_reshuffle, P_PT, P_TrB

def get_breuer_hall_U(d):
    """Constructs the antisymmetric unitary U for the Breuer-Hall map."""
    U = np.zeros((d, d), dtype=complex)
    for i in range(d // 2):
        U[2*i, 2*i+1] = 1.0
        U[2*i+1, 2*i] = -1.0
    return U

def optimize_channel(target, J_fixed, rho_fixed, d, U):
    """
    Optimizes either T1 or T2 to minimize Tr(rho * M(T1, T2)).
    
    Parameters:
    - target: String, either 'T1' or 'T2' (the variable to optimize).
    - J_fixed: numpy array (d^2 x d^2), the Choi matrix of the fixed channel.
    - rho_fixed: numpy array (d^2 x d^2), the fixed state (eigenvector).
    - d: Integer, dimension of the subsystem (must be even, >= 4).
    
    Returns:
    - min_val: The minimum expectation value found.
    - J_opt: The Choi matrix of the optimized channel.
    """
    d2 = d**2
    P_reshuffle, P_PT, P_TrB = build_linear_maps(d)
    P_reshuffle_inv = P_reshuffle.T # Inverse of permutation is its transpose
    
    # Define the CVXPY Variable for the target Choi matrix
    J_var = cp.Variable((d2, d2), hermitian=True)
    
    # Convert fixed Choi to Superoperator
    K_fixed = (P_reshuffle @ J_fixed.flatten()).reshape((d2, d2), order = 'C')
    
    # Convert variable Choi to Superoperator
    K_var = cp.reshape(P_reshuffle @ cp.vec(J_var, order='C'), (d2, d2), order = 'C')
    
    # Compose channels (Matrix multiplication of superoperators)
    if target == 'T1':
        # T1 is var, T2 is fixed. K_composed = K_var @ K_fixed
        K_composed = K_var @ K_fixed
    elif target == 'T2':
        # T2 is var, T1 is fixed. K_composed = K_fixed @ K_var
        K_composed = K_fixed @ K_var
    else:
        raise ValueError("Target must be 'T1' or 'T2'")
        
    # Convert composed Superoperator back to Choi matrix
    J_composed = cp.reshape(P_reshuffle_inv @ cp.vec(K_composed, order='C'), (d2, d2), order = 'C')
    
    # --- Apply Breuer-Hall Map to J_composed ---
    #U = get_breuer_hall_U(d)
    I_U = np.kron(np.eye(d), U)
    
    # J_composed is TP, so Tr_B(J) x I = I_d^2
    term1 = np.eye(d2)
    term2 = J_composed
    
    # Apply (I x U) * J_composed^{T_B} * (I x U^dagger)
    J_comp_PT = cp.reshape(P_PT @ cp.vec(J_composed, order='C'), (d2, d2), order = 'C')
    term3 = I_U @ J_comp_PT @ I_U.conj().T
    
    M = term1 - term2 - term3
    
    # --- Objective ---
    objective = cp.Minimize(cp.real(cp.trace(rho_fixed @ M)))
    
    # --- Constraints for J_var to be a valid PPT Channel ---
    constraints = []
    constraints.append(J_var >> 0) # Completely Positive
    
    # Partial Transpose Constraint (PPT)
    J_var_PT = cp.reshape(P_PT @ cp.vec(J_var, order='C'), (d2, d2), order = 'C')
    constraints.append(J_var_PT >> 0)
    
    # Trace Preserving Constraint: Tr_B(J) = I_d
    TrB_J_var = cp.reshape(P_TrB @ cp.vec(J_var, order='C'), (d, d), order = 'C')
    constraints.append(TrB_J_var == np.eye(d))
    
    # --- Solve ---
    prob = cp.Problem(objective, constraints)
    prob.solve(solver=cp.SCS) # SCS or MOSEK recommended for large complex SDPs
    
    return prob.value, J_var.value

def calculate_M_matrix(J1, J2, U, d=6):
    d2 = d**2
    P_reshuffle, P_PT, P_TrB = build_linear_maps(d)
    
    # 1. Compose exactly like the SDP
    K1 = (P_reshuffle @ J1.flatten()).reshape((d2, d2), order='C')
    K2 = (P_reshuffle @ J2.flatten()).reshape((d2, d2), order='C')
    
    # K2 is applied first, then K1. Let's match your T1/T2 logic.
    # In your loop: min_val, J1 = optimize_channel('T1', J2...) -> K_composed = K_var @ K_fixed -> K1 @ K2
    # Wait, if T1 is applied after T2, it's K1 @ K2.
    K_composed = K1 @ K2 
    
    J_composed = (P_reshuffle.T @ K_composed.flatten()).reshape((d2, d2), order='C')
    
    # 2. PT exactly like the SDP
    J_comp_PT = (P_PT @ J_composed.flatten()).reshape((d2, d2), order='C')
    
    # 3. Breuer-Hall
    term1 = np.eye(d2)
    term2 = J_composed
    I_U = np.kron(np.eye(d), U)
    term3 = I_U @ J_comp_PT @ I_U.conj().T
    
    M = term1 - term2 - term3
    
    # Ensure Hermitian (SCS numerical noise can make it slightly non-Hermitian)
    M = (M + M.conj().T) / 2.0
    
    return M


d = 6
stop = False
# Initialize a random maximally mixed Choi matrix as a starting point
# J_fixed = np.eye(d**2) / d 
# rho_fixed = np.eye(d**2) / (d**2)
# J2 = J_fixed
# rho = rho_fixed
for j in range(10000):
    if stop: 
        break
    print(j)
    
    rho_fixed = generate_random_density_matrix(d**2)
    P_reshuffle, P_PT, P_TrB = build_linear_maps(d)
    J_fixed = generate_random_boundary_ppt_channel(d, P_PT, P_TrB)
    rho = rho_fixed
    J2 = J_fixed
    
    U = get_random_breuer_hall_U(d)
    n = 3
    #print('d=', d)
    for i in range(n):
        # Optimize T1 while keeping T2 fixed
        _, J1 = optimize_channel('T1', J2, rho, d, U) # Discard prob.value
        min_val = np.real(np.trace(rho @ calculate_M_matrix(J1, J2, U, d)))
        _, J2 = optimize_channel('T2', J1, rho, d, U) # Discard prob.value
        min_val = np.real(np.trace(rho @ calculate_M_matrix(J1, J2, U, d)))
        
        d2 = d**2
        K1 = (P_reshuffle @ J1.flatten()).reshape((d2, d2), order='C')
        K2 = (P_reshuffle @ J2.flatten()).reshape((d2, d2), order='C')
        K_composed = K1 @ K2
        J_composed = (P_reshuffle.T @ K_composed.flatten()).reshape((d2, d2), order='C')
        J_PT = (P_PT @ J_composed.flatten()).reshape((d2, d2), order='C')
        
        if i % 4 == 0 and i != 0:
            #print(i)
            U_new = optimize_U_scalable(rho, J_PT, d, get_breuer_hall_U(d))
            val_new = np.real(np.trace(rho @ calculate_M_matrix(J1, J2, U_new, d)))
            if val_new < min_val:
                U = U_new
        
        #print('old min val', min_val)
        rho, min_val = update_rho(calculate_M_matrix(J1, J2, U, d))
    
        if min_val < -1e-3:
            print('weeee ddidd ittt!!!!!!1')
            print(i, min_val)
            print(J1, J2, U, rho)
            stop = True
            break
    print(min_val)
\end{lstlisting}



\addcontentsline{toc}{chapter}{\protect Bibliography}
\printbibliography

\end{document}